\documentclass[11pt, titlepage, bibliography=totoc]{scrreprt}
\setcounter{tocdepth}{3}

\usepackage[a4paper, hmargin=2.5cm, vmargin=3cm]{geometry}
\usepackage[T1]{fontenc}
\usepackage{setspace}
% remove indentation
\setlength{\parindent}{0pt}   
% add spacing between paragraphs
\setlength{\parskip}{1em}     
\usepackage{lmodern}
\usepackage[utf8]{inputenc}
\usepackage[english]{babel}
\usepackage{microtype}
\usepackage{csquotes}

\usepackage{amsmath}
\usepackage{blkarray}
\usepackage{tikz}
\usepackage{booktabs}
\usepackage{tabularx}
\usepackage{pgfplots}
\usepackage{graphicx}
\usepackage{caption, ragged2e}
\usepackage{subcaption}
\usepackage{multirow}
\usepackage{float}
\usepackage{empheq}
\usepackage{amssymb, amsthm}
\usepackage[automark]{scrlayer-scrpage}
\usepackage{listings}
\usepackage{scrhack}
\usepackage[linesnumbered,ruled,vlined]{algorithm2e}
\usepackage{comment}
% needed for \color in listings
\usepackage{xcolor} 
\usepackage{hyperref}

\usepackage{array}

\hypersetup{
    colorlinks=true,
    linkcolor=blue,
    citecolor=blue,
    urlcolor=blue,
    pdftitle={Categorizing Incident Reports Using Traditional Machine Learning and Large Language Models},
    pdfauthor={Md Shariar Imroze Khan},
    pdfsubject={Master Thesis - University of Koblenz},
    pdfkeywords={Process Safety, Incident Classification, Machine Learning, NLP, BERT, LLM}
}

\thispagestyle{empty}

\usetikzlibrary{automata, positioning, arrows, intersections}
\pgfplotsset{compat=1.18}

\newcommand{\matindex}[1]{\mbox{\scriptsize#1}}
\newcolumntype{L}{>{\RaggedRight\arraybackslash}X}

\pagestyle{scrheadings}
\setkomafont{disposition}{\normalfont\bfseries}
\clearpairofpagestyles
\ohead{\headmark}
\ofoot*{\pagemark}

\lstset{
    language=Python,
    basicstyle=\ttfamily\small,
    keywordstyle=\color{blue},
    stringstyle=\color{red},
    commentstyle=\color{green},
    morecomment=[l][\color{magenta}]{\#},
    columns=fullflexible,
    breaklines=true,
    breakatwhitespace=true,
}

\usepackage[style=numeric,sorting=none,maxbibnames=4]{biblatex}
\addbibresource{Literature.bib}

\begin{document}

\begin{titlepage}
\newcommand{\trtitle}{Categorizing Incident Reports Using Traditional Machine Learning and Large Language Models: A Comparative Study}
\newcommand{\trtype}{Master Thesis}
\newcommand{\trdate}{Winter Term 2025/26, \today}
% \newcommand{\trdate}{Winter Term 2025/26, 20 April 2026} use a fixed date
\newcommand{\trauthor}{Md Shariar Imroze Khan (Matriculation No: 220202354)}

\newcommand{\trInstitute}{submitted in the degree programme \\ M.Sc.\ Mathematical Modelling, Simulation and Optimization \vspace{1em}}
\newcommand{\trUniversity}{Faculty 3: Mathematics / Natural Sciences\\ University of Koblenz}

\begin{flushright}
\includegraphics[height=50pt]{Images/uniLogo.pdf}
\end{flushright}
\vspace{2.5cm}
\begin{center}
  \textbf{\Huge \trtitle}
\end{center}
\vspace{2cm}

\begin{center}
        \textbf{\Large \trtype}\\[1em]
        \trInstitute \\
        \trUniversity \\[1em]
    Submitted by \\[1em]
        \textbf{\large \trauthor}
   
  \vspace{2em}
  \trdate \\ [4em]
  \begin{tabular}{rl}
    First Supervisor: & Prof. Dr. Andreas Mauthe\\
                & {\footnotesize Institut für Wirtschafts- und Verwaltungsinformatik}\\
                & {\footnotesize FG IT Security and Data Security}\\[0.8em]
    Second Supervisor: & Alexander Rosenbaum\\
                & {\footnotesize Institut für Wirtschafts- und Verwaltungsinformatik}\\
                & {\footnotesize FG IT Security and Data Security}\\[0.8em]
    Company Supervisor: & Dr. Hugh Burnham-Slipper\\
                & {\footnotesize Uniper Kraftwerke GmbH}\\
\end{tabular}

\end{center}
\vspace{1cm}
\vfill

\end{titlepage}

\chapter*{Abstract}
\addcontentsline{toc}{chapter}{Abstract}

Incident reporting is a critical process for maintaining safety and operational efficiency across industries such as energy, mining, and healthcare. However, manual classification of incident reports remains prone to human error, inconsistency, and inefficiency. This thesis investigates the application of machine learning and large language models (LLMs) for automated classification of process safety incidents in multilingual industrial datasets.

A comprehensive experimental framework was developed across Iterations~0--8, progressively advancing from a controlled BERT-based baseline with Support Vector Machines to more sophisticated multilingual transformer and LLM-based approaches. Iteration~8 is retained as a placeholder for a final integrated refinement stage that combines the strongest findings from the earlier experiments; no empirical results are reported for this stage in the present thesis.

Key contributions include a detailed preprocessing pipeline for multilingual incident data, comparative benchmarking of more than 10 transformer architectures, advanced class balancing strategies including SMOTE and language-aware undersampling, ensemble methods combining multiple classifiers, and LLM-based few-shot classification using Flan-T5 models.


Experimental results show that Iteration~0 compared three alternative English subset construction strategies under a common BERT-base-uncased embedding pipeline with feature standardisation and an RBF-kernel SVM classifier. The manually curated English subset was retained as the official baseline and achieved an accuracy of 0.8112 and a macro F1-score of 0.7825. The automatically constructed Lingua-based and langdetect-based subsets achieved macro F1-scores of 0.7714 and 0.7648, respectively. Although the automatically constructed subsets produced competitive results, the manually curated English subset was retained as the official baseline because it provides the most controlled and auditable reference point for subsequent comparisons. Building on this baseline, later iterations improved performance through multilingual embeddings, stronger transformer architectures, and advanced machine learning techniques. The study addresses four research questions concerning classification accuracy improvements, traditional ML versus LLM comparisons, hazard type extraction, and a comprehensive severity prediction framework for the business case.


\cleardoublepage

\chapter*{Zusammenfassung}
\addcontentsline{toc}{chapter}{Zusammenfassung}

Die Meldung von Vorfällen ist ein zentraler Prozess zur Gewährleistung von Sicherheit und betrieblicher Effizienz in Branchen wie der Energieversorgung, dem Bergbau und dem Gesundheitswesen. Die manuelle Klassifikation von Vorfallsberichten ist jedoch weiterhin anfällig für menschliche Fehler, Inkonsistenzen und Ineffizienz. Diese Arbeit untersucht den Einsatz von Machine Learning und Large Language Models (LLMs) zur automatisierten Klassifikation von Process-Safety-Vorfällen in mehrsprachigen industriellen Datensätzen.

Im Rahmen der Untersuchung wurde ein umfassendes experimentelles Framework über die Iterationen~0--8 entwickelt, das sich schrittweise von einer kontrollierten BERT-basierten Baseline mit Support Vector Machines hin zu anspruchsvolleren mehrsprachigen Transformer- und LLM-basierten Ansätzen weiterentwickelte. Iteration~8 wird dabei als Platzhalter für eine abschließende integrierte Verfeinerungsphase beibehalten, in der die stärksten Befunde der vorangegangenen Experimente zusammengeführt werden sollen; für diese Phase werden in der vorliegenden Arbeit keine empirischen Ergebnisse berichtet.

Zu den zentralen Beiträgen gehören eine detaillierte Vorverarbeitungspipeline für mehrsprachige Vorfallsdaten, ein vergleichendes Benchmarking von mehr als 10 Transformer-Architekturen, fortgeschrittene Strategien zum Klassen-Balancing einschließlich SMOTE und sprachbewusstem Undersampling, Ensemble-Methoden zur Kombination mehrerer Klassifikatoren sowie eine LLM-basierte Few-Shot-Klassifikation unter Verwendung von Flan-T5-Modellen.

Die experimentellen Ergebnisse zeigen, dass in Iteration~0 drei alternative Strategien zur Bildung englischsprachiger Teilmengen unter einer gemeinsamen Pipeline aus BERT-base-uncased-Embeddings, Merkmalsstandardisierung und einem SVM-Klassifikator mit RBF-Kernel verglichen wurden. Der manuell kuratierte englische Teil-Datensatz wurde als offizielle Baseline beibehalten und erreichte eine Accuracy von 0{,}8112 sowie einen Macro-F1-Score von 0{,}7825. Die automatisch konstruierten Lingua- bzw. langdetect-basierten Teilmengen erreichten Macro-F1-Scores von 0{,}7714 bzw. 0{,}7648. Obwohl auch die automatisch konstruierten Teilmengen konkurrenzfähige Ergebnisse erzielten, wurde der manuell kuratierte englische Teil-Datensatz als offizielle Baseline beibehalten, da er den kontrolliertesten und am besten nachvollziehbaren Referenzpunkt für nachfolgende Vergleiche darstellt. Aufbauend auf dieser Baseline verbesserten nachfolgende Iterationen die Leistung durch mehrsprachige Embeddings, leistungsfähigere Transformer-Architekturen und fortgeschrittene Machine-Learning-Techniken. Die Arbeit behandelt vier Forschungsfragen zu Verbesserungen der Klassifikationsgenauigkeit, zum Vergleich traditioneller ML-Verfahren mit LLMs, zur Extraktion von Gefährdungstypen sowie zu einem umfassenden Framework zur Vorhersage von Schweregraden für den Business Case.

\cleardoublepage

\spacing{1.15}
\tableofcontents
\thispagestyle{empty}
\singlespacing
\cleardoublepage
\pagenumbering{arabic}

% ============================================================================
% ACKNOWLEDGMENTS
% ============================================================================
\chapter*{Acknowledgments}
\addcontentsline{toc}{chapter}{Acknowledgments}

I would like to express my sincere gratitude to all those who have supported me throughout this research journey.

First and foremost, I extend my deepest appreciation to my supervisors, \textbf{Prof.\ Dr.\ Andreas Mauthe} and \textbf{Alexander Rosenbaum} from the University of Koblenz, as well as \textbf{Dr.\ Hugh Burnham-Slipper} from Uniper Kraftwerke GmbH. Their expert guidance, constructive feedback, and continuous encouragement have been invaluable in shaping this research.

I am particularly grateful to Uniper Kraftwerke GmbH for providing access to the industrial incident datasets that formed the foundation of this study. I would also like to thank my manager, \textbf{Craig Walker}, for his support and motivation throughout the thesis process. The opportunity to work with real-world safety data has been instrumental in ensuring the practical relevance of this work.

I would also like to thank my colleagues and fellow researchers, namely \textbf{Zahid Hasan Polash} and \textbf{Christopher Thiele}, for their valuable discussions and insights throughout this project.

Finally, I extend my heartfelt thanks to my family and friends for their unwavering support, patience, and encouragement during my studies.

\cleardoublepage

% ============================================================================
% ENGLISH VERSION
% ============================================================================

\chapter*{Declaration of authorship}
\addcontentsline{toc}{chapter}{Declaration of authorship}

I declare that I have written this thesis independently
and have not used any sources or aids other than those indicated
and that this thesis has not been submitted in the same or a similar form
to any other examination authority and accepted by it as part of an examination.

\vspace{10mm}

\noindent
\begin{tabular}{@{}p{0.82\linewidth}@{\hspace*{2ex}}r@{\hspace*{2ex}}r}
I agree to the inclusion of this thesis in the library.
& yes $\square$ & no $\square$ \\[1em]
I agree to the publication of this thesis on the Internet.
& yes $\square$ & no $\square$ \\
\end{tabular}

\vspace{20mm}

\noindent
\dotfill

\noindent
(Place, Date) \hspace{8.5cm} (Signature)

\cleardoublepage

% ============================================================================
% EIDESSTATTLICHE ERKLÄRUNG / DECLARATION
% ============================================================================

\chapter*{Eidesstattliche Erklärung}
\addcontentsline{toc}{chapter}{Eidesstattliche Erklärung}

Ich versichere, dass ich die vorliegende Arbeit selbständig verfasst
und keine anderen als die angegebenen Quellen und Hilfsmittel benutzt habe
und dass die Arbeit in gleicher oder ähnlicher Form noch keiner anderen Prüfungsbehörde
vorgelegen hat und von dieser als Teil einer Prüfungsleistung angenommen wurde.

\vspace{10mm}

\noindent
\begin{tabular}{@{}p{0.82\linewidth}@{\hspace*{2ex}}r@{\hspace*{2ex}}r}
Mit der Einstellung der Arbeit in die Bibliothek bin ich einverstanden.
& ja $\square$ & nein $\square$ \\[1em]
Der Veröffentlichung dieser Arbeit im Internet stimme ich zu.
& ja $\square$ & nein $\square$ \\
\end{tabular}

\vspace{20mm}

\noindent
\dotfill

\noindent
(Ort, Datum) \hspace{8.5cm} (Unterschrift)

\cleardoublepage

% ============================================================================
% LIST OF FIGURES
% ============================================================================
\listoffigures
\addcontentsline{toc}{chapter}{List of Figures}
\cleardoublepage

% ============================================================================
% LIST OF TABLES
% ============================================================================
\listoftables
\addcontentsline{toc}{chapter}{List of Tables}
\cleardoublepage

% ============================================================================
% LIST OF ABBREVIATIONS
% ============================================================================
\chapter*{List of Abbreviations}
\addcontentsline{toc}{chapter}{List of Abbreviations}

\noindent
\begin{tabular}{@{}p{2.5cm}p{12cm}@{}}
\textbf{ADASYN} & Adaptive Synthetic Sampling \\[0.3em]
\textbf{AI} & Artificial Intelligence \\[0.3em]
\textbf{API} & Application Programming Interface \\[0.3em]
\textbf{AUC} & Area Under the Curve \\[0.3em]
\textbf{BERT} & Bidirectional Encoder Representations from Transformers \\[0.3em]
\textbf{BoW} & Bag of Words \\[0.3em]
\textbf{CBOW} & Continuous Bag of Words \\[0.3em]
\textbf{CCPS} & Center for Chemical Process Safety \\[0.3em]
\textbf{CLS} & Classification Token \\[0.3em]
\textbf{CNN} & Convolutional Neural Network \\[0.3em]
\textbf{CSV} & Comma-Separated Values \\[0.3em]
\textbf{DeBERTa} & Decoding-enhanced BERT with Disentangled Attention \\[0.3em]
\textbf{ENN} & Edited Nearest Neighbors \\[0.3em]
\textbf{ESD} & Emergency Shutdown \\[0.3em]
\textbf{F1} & F1-Score (Harmonic Mean of Precision and Recall) \\[0.3em]
\textbf{FN} & False Negative \\[0.3em]
\textbf{FNR} & False Negative Rate \\[0.3em]
\textbf{FP} & False Positive \\[0.3em]
\textbf{FPR} & False Positive Rate \\[0.3em]
\textbf{GPT} & Generative Pre-trained Transformer \\[0.3em]
\textbf{GPU} & Graphics Processing Unit \\[0.3em]
\textbf{GT} & Gas Turbine \\[0.3em]
\textbf{HRSG} & Heat Recovery Steam Generator \\[0.3em]
\textbf{HSSE} & Health, Safety, Security, and Environment \\[0.3em]
\textbf{LightGBM} & Light Gradient Boosting Machine \\[0.3em]
\textbf{LIME} & Local Interpretable Model-agnostic Explanations \\[0.3em]
\textbf{LLM} & Large Language Model \\[0.3em]
\textbf{LOPC} & Loss of Primary Containment \\[0.3em]
\textbf{LSTM} & Long Short-Term Memory \\[0.3em]
\end{tabular}

\begin{tabular}{@{}p{2.5cm}p{12cm}@{}}
\textbf{MCC} & Matthews Correlation Coefficient \\[0.3em]
\textbf{ML} & Machine Learning \\[0.3em]
\textbf{MLM} & Masked Language Modeling \\[0.3em]
\textbf{mBERT} & Multilingual BERT \\[0.3em]
\textbf{NLP} & Natural Language Processing \\[0.3em]
\textbf{NPS} & Non-Process Safety \\[0.3em]
\textbf{NSP} & Next Sentence Prediction \\[0.3em]
\textbf{OSHA} & Occupational Safety and Health Administration \\[0.3em]
\textbf{PII} & Personally Identifiable Information \\[0.3em]
\textbf{PS} & Process Safety \\[0.3em]
\textbf{RAG} & Retrieval-Augmented Generation \\[0.3em]
\textbf{RBF} & Radial Basis Function \\[0.3em]
\textbf{RF} & Random Forest \\[0.3em]
\textbf{RNN} & Recurrent Neural Network \\[0.3em]
\textbf{RoBERTa} & Robustly Optimized BERT Pretraining Approach \\[0.3em]
\textbf{ROC} & Receiver Operating Characteristic \\[0.3em]
\textbf{RP} & Recommended Practice \\[0.3em]
\textbf{RQ} & Research Question \\[0.3em]
\textbf{SHAP} & SHapley Additive exPlanations \\[0.3em]
\textbf{SIL} & Safety Integrity Level \\[0.3em]
\textbf{SMOTE} & Synthetic Minority Over-sampling Technique \\[0.3em]
\textbf{SVM} & Support Vector Machine \\[0.3em]
\textbf{TF-IDF} & Term Frequency-Inverse Document Frequency \\[0.3em]
\textbf{TN} & True Negative \\[0.3em]
\textbf{TP} & True Positive \\[0.3em]
\textbf{XGBoost} & Extreme Gradient Boosting \\[0.3em]
\textbf{XLM} & Cross-lingual Language Model \\[0.3em]
\textbf{XLM-R} & Cross-lingual Language Model - RoBERTa \\[0.3em]
\textbf{XLNet} & Generalized Autoregressive Pretraining for Language Understanding \\[0.3em]
\end{tabular}

\cleardoublepage

% ============================================================================
% GLOSSARY
% ============================================================================
\chapter*{Glossary}
\addcontentsline{toc}{chapter}{Glossary}

\textbf{Attention Mechanism:} A neural network component that allows models to focus on relevant parts of the input sequence when producing output, enabling capture of long-range dependencies.

\textbf{Class Imbalance:} A situation in classification tasks where one class (minority) has significantly fewer samples than another class (majority), which can bias model predictions toward the majority class.

\textbf{[CLS] Token:} A special token prepended to input sequences in BERT-style models, whose final hidden state is used as the aggregate sequence representation for classification tasks.

\textbf{Embedding:} A dense vector representation of text (words, sentences, or documents) in a continuous vector space where semantic similarity is captured by vector proximity.

\textbf{False Negative:} An instance where the model incorrectly predicts the negative class when the true label is positive (e.g., missing a Process Safety incident).

\textbf{False Positive:} An instance where the model incorrectly predicts the positive class when the true label is negative (e.g., incorrectly flagging a non-PS incident as PS).

\textbf{Few-Shot Learning:} A machine learning paradigm where models learn to classify new examples with only a small number of labeled training examples per class.

\textbf{Fine-Tuning:} The process of adapting a pre-trained model to a specific downstream task by continuing training on task-specific data.

\textbf{Macro F1-Score:} The unweighted average of F1-scores across all classes, treating each class equally regardless of sample size.

\textbf{Mean Pooling:} An aggregation strategy that computes the element-wise mean of all token embeddings (weighted by attention mask) to produce a single document representation.

\textbf{Process Safety:} A disciplined framework for managing the integrity of operating systems and processes handling hazardous substances, focusing on prevention of uncontrolled releases that could result in fires, explosions, or environmental damage.

\textbf{SMOTE (Synthetic Minority Over-sampling Technique):} An oversampling method that generates synthetic samples for the minority class by interpolating between existing minority samples and their nearest neighbors.

\textbf{Tokenization:} The process of breaking text into smaller units (tokens) that serve as input to NLP models, using subword algorithms like WordPiece or BPE.

\textbf{Transfer Learning:} A machine learning technique where a model trained on one task is repurposed for a related task, leveraging learned representations to improve performance with limited task-specific data.

\textbf{Transformer:} A neural network architecture based entirely on self-attention mechanisms, eliminating recurrence and enabling parallel processing of sequences.

\textbf{Zero-Shot Classification:} Classification of inputs into categories that the model was not explicitly trained on, relying on the model's general knowledge and understanding of class descriptions.

\cleardoublepage

%%%%%%%%%%%%%%%%%%%%%%%%%%%%%%%%%%%%%%%%%%%%%%%%%%%%%%%%%%%%%%%%%%%%%%%%%%%%%%
\chapter{Introduction}

Incident reporting is a critical process in industries such as healthcare, aviation, energy, mining, construction, and information technology. Accurate classification and analysis of incident reports are essential for maintaining safety, ensuring operational efficiency, and meeting regulatory requirements enforced by agencies such as OSHA and the U.S. Chemical Safety Board \cite{1_Baybutt}.

Process safety incidents, though infrequent, can have catastrophic consequences. The Bhopal disaster of 1984, one of the deadliest industrial accidents in history, highlights the devastating impact of inadequate safety measures and incident response \cite{2_Broughton}. Insights from subsequent investigations emphasize the need for robust incident management systems to prevent, detect, and mitigate risks \cite{Stefana2024}.

Despite advancements in digital reporting protocols, manual classification of incidents remains prone to error. Human judgment introduces subjectivity, inconsistencies, and fatigue-related mistakes \cite{4_Kurian}. These challenges have been observed across construction \cite{Zhong2020}, oil and gas \cite{Sattari2021OilGas}, aviation \cite{Madeira2021}, mining \cite{8_Ganguli}, agriculture \cite{9_Kakhki}, and healthcare \cite{10_Mathew}. Such errors compromise the reliability of incident data, leading to flawed decision-making, suboptimal resource allocation, and delayed responses \cite{11_Mustapha}.

This research proposes the application of machine learning and large language models (LLMs) to automate and improve the classification of text-based incident reports. The aim is to design an ML model that reduces errors arising from manual processes \cite{12_Silva}, improves efficiency and consistency in classification \cite{13_Bugalia}, and adapts dynamically to evolving business and regulatory requirements \cite{14_Chen}.

Through the integration of classical ML techniques, modern deep learning approaches, and transformer-based language models, this research seeks to contribute to the development of scalable, explainable, and multilingual incident classification systems.

\newpage

%%%%%%%%%%%%%%%%%%%%%%%%%%%%%%%%%%%%%%%%%%%%%%%%%%%%%%%%%%%%%%%%%%%%%%%%
\section{Motivation}

Effective incident reporting is central to risk management, yet current systems face persistent challenges. Manual classification is subject to variability in human judgment, fatigue, and inconsistent criteria, which reduces data reliability and weakens safety interventions \cite{Sattari2021OilGas}. As industries grow more complex, incident datasets expand in volume and linguistic diversity, amplifying classification challenges \cite{13_Bugalia}.

Language barriers further complicate reporting, with multinational organizations dealing with reports in multiple languages and varying levels of detail \cite{macleod2024loss}. Short, ambiguous, or error-prone reports reduce the effectiveness of traditional classification methods, requiring advanced NLP approaches for meaningful analysis.

Machine learning has shown promise in domains such as construction \cite{Zhong2020}, aviation \cite{Madeira2021}, mining \cite{8_Ganguli}, and healthcare \cite{islam2024}, where it consistently outperforms manual classification inaccuracies and inefficiencies. ML also uncovers hidden patterns in near-miss and incident data, supporting proactive risk reduction \cite{Fang}.

The rise of explainable AI frameworks enables organisations to balance predictive performance with interpretability, which is essential in safety-critical contexts \cite{sarkar2020review}. Furthermore, ML enables multilingual processing, severity prediction, and adaptation to evolving categories (e.g., WCM, leaks, fires), making it a transformative tool for safety management \cite{balasubramanian2021safeone}.

By adopting ML-driven classification systems, organisations can enhance decision-making, strengthen safety outcomes, and implement proactive risk management \cite{Stefana2024}.

\newpage

%%%%%%%%%%%%%%%%%%%%%%%%%%%%%%%%%%%%%%%%%%%%%%%%%%%%%%%%%%%%%%%%%%%%%%%%%%%%%%%
\section{Research Objectives, Questions, and Expected Outcomes}
\label{chapterGeneral}

This research aims to address the shortcomings of manual incident reporting by proposing and evaluating an automated, machine learning-driven framework for incident classification. Manual classification processes are often affected by human error, inconsistency, and inefficiency, particularly when incident reports are produced across multilingual and domain-specific contexts. To address these limitations, this study investigates whether machine learning and large language model (LLM)-based approaches can improve the reliability, scalability, and practical utility of incident classification.

The research is guided by four core objectives. First, it seeks to analyse the limitations of current manual reporting processes and identify major sources of inconsistency, inefficiency, and error. Second, it aims to develop and evaluate models capable of accurately classifying incident types across multilingual datasets. Third, it investigates whether hazard types can be extracted automatically from incident titles and descriptions using LLM-based approaches. Fourth, it explores the suitability of a proposed classification framework for supporting incident severity prediction in a business-oriented use case.


Building upon Iteration~0, which compared three English subset construction strategies under a fixed BERT-base-uncased embedding pipeline with feature standardisation and an RBF-kernel SVM classifier, the manually curated English subset was retained as the official baseline. This official baseline achieved an accuracy of 0.8112 and a macro-average F1-score of 0.7825. The automatic Lingua-based and langdetect-based subsets achieved macro-average F1-scores of 0.7714 and 0.7648, respectively. Since the manually curated subset delivered the strongest macro-F1 performance, it is adopted as the official baseline for subsequent comparisons and target setting. The study therefore addresses the following research questions with measurable performance targets in order to establish a clear link between the experimental evaluation and the broader objectives of improving classification reliability, interpretability, and scalability.

\subsection{Research Questions}


    \textbf{RQ1:} Can the classification of \textit{Process Safety} incidents across multilingual and domain-specific datasets be improved by at least \textbf{6.75 percentage points} in macro-average F1-score (from approximately 0.7825 to at least 0.85) compared with the baseline BERT \& SVM model?\\

    \textbf{RQ2:} How do traditional machine learning models (e.g., SVM, XGBoost, LightGBM, Random Forest, and ensemble methods) compare with transformer-based and LLM-based models (e.g., BERT, RoBERTa, DeBERTa, GPT-based approaches, or LLaMA variants) in achieving \textbf{macro-F1 $\geq 0.7825$}?\\

    \textbf{RQ3:} Can large language models (LLMs) be used to \textbf{extract and identify underlying hazard types from process safety incidents} (e.g., \textit{Fire/Explosion, Loss of Containment, Mechanical/Material Failure, Exposure to Toxic Substances}) based on incident \textit{titles and descriptions}, with \textbf{macro-F1 $\geq 0.7825$} across multilingual datasets?\\

    %\textbf{RQ4:} How suitable is a proposed classification framework for predicting incident severity levels (e.g., low, medium, severe) in the business case?\\

    \textbf{RQ4:} What conceptual framework can be proposed for predicting incident severity levels (e.g., low, medium, severe) from multilingual incident reports in the business case, and how suitable is this framework in terms of methodological feasibility, interpretability, and practical applicability?


\subsection{Expected Outcomes}

The expected outcomes of this study are closely aligned with the research questions and serve as concrete evaluation targets for the proposed framework. First, the study aims to achieve \textbf{macro-F1 $\geq 0.85$} for Process Safety classification, thereby improving upon the official Iteration~0 baseline result of 0.7825 obtained on the manually curated English subset. In addition, class-level \textbf{precision} and \textbf{recall} will be reported for the Process Safety class, together with robustness analyses across country- and domain-specific dataset shifts. Second, the study seeks to provide a systematic benchmark comparing traditional machine learning and transformer-based approaches, with the expectation that at least one transformer-based or multilingual transformer-based method achieves \textbf{macro-F1 $\geq 0.85$}. This comparison will also consider computational efficiency, including \textbf{training time}, \textbf{peak memory usage}, and resource--performance trade-offs, supported by ablation studies on feature representations and document field selection. Third, the study aims to demonstrate the feasibility of automatically extracting underlying hazard types from incident titles and descriptions with \textbf{macro-F1 $\geq 0.7825$} across multilingual corpora. The analysis will include per-language performance reporting and an error analysis of frequent class confusions, such as \textit{Loss of Containment} versus \textit{Mechanical Failure}. Finally, the study proposes and assesses a multi-class severity classification framework for the business case, focusing on the categorisation of incidents into \textit{low}, \textit{medium}, and \textit{severe} severity levels. The emphasis is placed on evaluating the conceptual suitability and potential practical value of such a framework rather than claiming a production-ready implementation.

\newpage

\section{Thesis Structure}

This thesis is organized into six main chapters and an appendix as follows:

\textbf{Chapter 1: Introduction} provides the context and motivation for automated incident classification, presents the research objectives and questions, and outlines the expected outcomes.

\textbf{Chapter 2: Literature Review} surveys the evolution of AI and ML techniques for text processing, examines applications of NLP in incident classification across various industries, and identifies research gaps that this work addresses.

\textbf{Chapter 3: Methodology} describes the research design, including data collection and preprocessing, the proposed dataset structure, and the iterative experimental framework spanning Iterations~0--8.

\textbf{Chapter 4: Results} presents the experimental outcomes reported for Iterations~0--7 and includes a structured placeholder for Iteration~8 as the reserved final refinement stage.

\textbf{Chapter 5: Discussion} synthesizes findings in relation to the research questions, analyzes technical insights, discusses limitations, and outlines practical implications for deployment.

\textbf{Chapter 6: Conclusion} summarizes key contributions, evaluates research question outcomes, and proposes directions for future work.

\textbf{Appendix} provides supplementary materials including code snippets, additional experimental results, and data samples.

\newpage


%%%%%%%%%%%%%%%%%%%%%%%%%%%%%%%%%%%%%%%%%%%%%%%%%%%%%%%%%%%%%%%%%%%%%%%%%%%%%%%%

\chapter{Literature Review}

\section{Evolution of AI and ML Techniques for Text Processing}

\subsection{Historical Foundation and Development}
The field of artificial intelligence and natural language processing has undergone significant evolution since its inception in the 1950s. Turing's seminal work on ``Computing Machinery and Intelligence'' in 1950 introduced the concept of machine intelligence, while the term ``artificial intelligence'' was coined by John McCarthy during the 1956 Dartmouth Workshop \cite{mccarthy1955}. Early text analysis efforts were primarily rule-based systems that relied on manually crafted linguistic rules and dictionaries, exemplified by programs like \textsc{ELIZA} (1966), the first chatbot that simulated conversation using simple pattern matching and template responses \cite{weizenbaum1966}.

The late 1980s marked a revolutionary shift toward statistical approaches in natural language processing, transitioning from knowledge-driven to data-driven methodologies \cite{church1988}. This period saw the emergence of machine learning algorithms for language tasks, including the adoption of Bayesian methods for probabilistic inference and statistical approaches leveraging large corpora to understand language patterns. The 1990s witnessed significant progress with recurrent neural networks (RNNs) becoming popular for sequential data processing, though they faced limitations including the vanishing gradient problem and sequential processing requirements \cite{bengio1994}. To address these issues, Long Short-Term Memory (LSTM) networks were developed in 1997, incorporating memory cells and gating mechanisms to better handle long-term dependencies in text \cite{hochreiter1997}.

\subsection{The Deep Learning Revolution}
The 2000s brought the rise of more sophisticated machine learning techniques, including the development of foundational text representation methods. Traditional approaches such as Bag-of-Words (BoW) converted text into numerical format by representing documents as collections of word frequencies, while Term Frequency–Inverse Document Frequency (TF-IDF) improved upon BoW by weighting words based on their importance and rarity across the corpus \cite{salton1988}.

A breakthrough occurred in 2013 with the introduction of Word2Vec by Mikolov et al., which presented two key architectures: Continuous Bag-of-Words (CBOW) and Skip-gram for learning word embeddings \cite{mikolov2013a}. This was followed by GloVe (Global Vectors for Word Representation) in 2014, which used global word co-occurrence statistics to learn vector representations \cite{pennington2014}.

The 2017 publication of ``Attention Is All You Need'' by Vaswani et al. introduced the Transformer architecture, revolutionizing natural language processing by eliminating the need for recurrent processing through the self-attention mechanism \cite{vaswani2017}. This innovation enabled parallel processing of sequences, more effective capture of long-range dependencies, and scalability to larger datasets and model sizes. The landmark year 2018 saw the release of BERT (Bidirectional Encoder Representations from Transformers) by Google and GPT (Generative Pre-trained Transformer) by OpenAI, establishing the foundation for modern large language models \cite{radford2018}.

\section{Applications of ML and NLP in Incident Classification}

\subsection{Traditional Machine Learning Approaches}
Early applications of machine learning to incident classification predominantly employed traditional algorithms with handcrafted features. Support Vector Machines (SVM) have shown particular effectiveness in safety incident classification. Sattari et al. \cite{Sattari2021OilGas} demonstrated that Linear SVM achieved the highest accuracy (85.68\%) for assigning Process Safety Management elements to incident reports from oil and gas operations, outperforming other classifiers including Random Forest (76.52\%) and Multilayer Perceptron (83.16\%). Similarly, Zhang et al. \cite{zhang2019} applied text mining and ensemble techniques to construction site accident analysis using OSHA data, where their optimized ensemble model with Sequential Quadratic Programming (SQP) weight optimization achieved superior performance compared to individual classifiers.

Ensemble methods have gained prominence in incident classification due to their ability to combine multiple learning algorithms for enhanced performance. Ubeynarayana et al. \cite{ubeynarayana2017} evaluated ensemble approaches using popular machine learning algorithms including SVM, decision tree, linear regression, k-nearest neighbor, naive Bayes, and neural networks for construction accident classification. Their ensemble model achieved precision ranging from 0.6 to 1.0, recall from 0.18 to 0.96, and F1 scores from 0.14 to 0.96 across 11 accident type labels, with the highest average F1 score of 0.69.

\subsection{Deep Learning and Neural Network Approaches}
Convolutional Neural Networks (CNNs) and Recurrent Neural Networks have been successfully applied to incident text classification. Ren et al. \cite{ren2013} proposed using unsupervised feature learning algorithms to generate higher-level features for automatic incident detection in transportation systems. Their approach with enhanced features showed significant improvements in detection rate, false alarm rate, and mean time to detect across three representative test cases.

\subsection{Text Representation and Embeddings}
A central challenge in natural language processing is the conversion of textual input into numerical representations that can be processed by machine learning algorithms. Let a text sequence be denoted by
\begin{equation}
x = (w_1, w_2, \dots, w_n),
\end{equation}
where each $w_i$ represents a token drawn from a vocabulary $\mathcal{V}$.

Traditional text representation methods such as bag-of-words or TF--IDF encode documents as sparse vectors in a high-dimensional space. Although these approaches are simple and often effective, they do not explicitly capture semantic similarity or contextual relationships between words. In contrast, an embedding maps each token or document into a dense vector space of relatively low dimensionality, such that semantically or syntactically related units are positioned closer together.

Formally, an embedding can be defined as a function
\begin{equation}
f : \mathcal{V} \rightarrow \mathbb{R}^{d},
\end{equation}
where $d$ denotes the embedding dimension and $f(w_i)$ is the dense vector representation of token $w_i$. For a sequence of tokens, the corresponding embedding matrix is
\begin{equation}
\mathbf{E}(x) =
\begin{bmatrix}
f(w_1) \\
f(w_2) \\
\vdots \\
f(w_n)
\end{bmatrix}
\in \mathbb{R}^{n \times d}.
\end{equation}

A distinction must be made between \emph{static} and \emph{contextual} embeddings. Static embeddings assign one fixed vector to each word regardless of context. For example, a word with multiple meanings would be represented identically in all sentences. Contextual embeddings, by contrast, generate token representations conditioned on the surrounding sequence. This is particularly important in incident-report narratives, where the interpretation of a term may depend strongly on operational context, equipment references, or hazard descriptions.

For document-level classification tasks, token embeddings must be aggregated into a single fixed-length vector. If $\mathbf{h}_1, \mathbf{h}_2, \dots, \mathbf{h}_n$ denote contextual token representations, a document embedding $\mathbf{z}$ may be obtained by a pooling operation such as
\begin{equation}
\mathbf{z} = \mathrm{Pool}(\mathbf{h}_1, \mathbf{h}_2, \dots, \mathbf{h}_n),
\end{equation}
where $\mathrm{Pool}(\cdot)$ may correspond to mean pooling, max pooling, or a dedicated classification token representation. The resulting vector $\mathbf{z} \in \mathbb{R}^{d}$ can then be used as input to a downstream classifier such as an SVM, Random Forest, XGBoost model, or neural network.

Within the context of this thesis, embeddings are of particular importance because incident reports are short, domain-specific, and frequently multilingual. Dense semantic representations therefore provide a principled way to preserve contextual meaning while transforming free-text narratives into a format suitable for supervised classification.

\subsection{Transformer-Based Models in Safety Applications}
The emergence of transformer-based models has significantly advanced incident classification capabilities. Grigorev et al. \cite{grigorev2024} demonstrated the innovative integration of Large Language Models into machine learning workflows for traffic incident management, focusing on severity classification using accident reports. Their hybrid approach, leveraging features generated by modern language models alongside conventional data, showed improvements in classification accuracy across several machine learning algorithms when applied to datasets from the United States, United Kingdom, and Queensland, Australia.

In healthcare applications, recent research has shown the superior performance of BERT-based models. Wang and Magrabi \cite{Ying} evaluated BERT's transferability to patient safety incident classification, applying it to classify 10 incident types and 4 severity levels. Fine-tuned BERT achieved high F-scores above 89\% across all test datasets, significantly outperforming CNNs (81\%) and effectively generalizing to real-world settings, including rare incident types.

Zisad et al. \cite{zisad2025} conducted a comparative analysis of transformer models in disaster tweet classification for public safety applications. Their results showed that BERT achieved the highest accuracy (91\%), significantly outperforming traditional machine learning models like Logistic Regression and Naive Bayes (both at 82\%).

\section{Multilingual and Cross-Cultural Challenges}

\subsection{Language Barriers in Incident Reporting}
Multinational organizations face significant challenges in processing incident reports written in multiple languages with varying levels of detail and linguistic complexity. Hughes et al. \cite{hughes2019} developed an approach to extract meaning from multi-lingual free-text safety incident reports, analyzing 5065 safety incident reports in multiple languages. Their graph database approach achieved an overall performance of 98.5\% accuracy in extracting information.

The multilingual event detection research by Veyseh et al. \cite{veyseh2022} introduced MINION, a large-scale multilingual dataset for event detection covering 8 languages. This work highlighted the challenges and transferability of event detection across languages.

\subsection{Translation and Cross-Lingual Processing}
Steinberger et al.\cite{steinberger2006} provided insights on minimizing effort in developing multilingual NLP applications, demonstrating that minimalistic approaches using local grammars for information extraction can be effective for multilingual incident processing.

\section{Text Quality and Preprocessing Challenges}

\subsection{Data Quality Issues in Incident Reports}
Incident reports frequently exhibit quality issues that complicate automated processing. Kurian et al. \cite{4_Kurian} addressed these challenges by applying machine learning and keyword analysis to create a digitalized system for efficiently reporting incidents in oil and gas operations. Their research emphasized the importance of data preprocessing and quality improvement before applying machine learning algorithms.

The systematic review by Khairuddin et al. \cite{khairuddin2022} identified that text mining and NLP applications for occupational injury analysis require comprehensive preprocessing steps including stop-word removal, tokenization, stemming and lemmatization, POS tagging, N-grams, parsing, and semantic reasoning.

\subsection{Named Entity Recognition and Feature Extraction}
Advanced NLP techniques have been employed to extract structured information from unstructured incident reports. Tixier et al. \cite{tixier2016} focused on building NLP systems to automatically extract attributes and safety outcomes from unstructured reports using hand-coded rules and keyword dictionaries, achieving high precision of 0.95.

Xu et al. \cite{xu2021} developed an improved text mining approach to extract safety risk factors from construction accident reports, proposing a tailored domain lexicon to identify safety risk factors.

\section{Explainable AI and Interpretability}

\subsection{Need for Explainable Models}
The critical nature of safety applications necessitates interpretable machine learning models. Chen et al. \cite{chen2024} developed a machine learning approach with human-AI collaboration for patient safety event report classification, integrating LIME (Local Interpretable Model-Agnostic Explanations).

However, research by Slack et al. \cite{slack2020} revealed potential vulnerabilities in popular explanation techniques, demonstrating that adversarial classifiers can fool methods like LIME and SHAP.

\subsection{Practical Applications}
Despite limitations, explainable AI techniques continue to provide value in safety applications. Wen et al. \cite{wen2022} compared LIME, LSA, and PDP for crash risk analysis, showing their usefulness when properly applied. Yuan et al. \cite{yuan2022} applied explainable ML for real-time safety analysis using SHAP values to reveal influential factors in accident fatalities.

\section{Emerging Trends and Future Directions}

\subsection{Large Language Models and Few-Shot Learning}
Recent developments in large language models have opened new possibilities for incident classification with limited training data. Mylius and Bernardi \cite{mylius2025} developed a scalable AI incident classification tool using LLMs to process raw incident reports, achieving approximately 80\% consistency with expert human ratings.

\subsection{Multimodal and Real-Time Processing}
Current research indicates movement toward multimodal AI systems combining text with other data sources. Anik et al. \cite{anik2024} proposed a time-embedded attention-based Transformer model for crash prediction, demonstrating potential for integrating textual incident reports with temporal and spatial data.

\section{Research Gaps and Opportunities}
While significant progress has been made, comprehensive comparative studies evaluating traditional ML methods against modern transformer-based approaches remain limited. There is also a need for domain-specific adaptations, standardized evaluation frameworks, and practical deployment considerations.

This review highlights that although ML and NLP have advanced incident classification, opportunities remain in multilingual processing, explainability, and systematic comparison of classical ML and LLM-based approaches.

\chapter{Methodology}

\section{Data Collection}
Historical incident reports were gathered from multiple internal Uniper HSSE reporting streams in order to ensure diversity in reporting scope, operational context, and incident type. The collected data exhibited several quality issues, including overly brief descriptions, spelling mistakes, grammatical irregularities, inconsistent reporting styles, and multilingual content. These characteristics motivated the need for a structured preprocessing pipeline before machine learning experiments were conducted. In addition, the presence of reports in multiple languages complicated direct integration into a single monolingual modelling pipeline and motivated the later development of language-specific and multilingual experiments.

\section{Proposed Dataset}

The dataset used in this thesis was not taken from a single precompiled table. Instead, it was constructed by integrating multiple raw reporting streams from the internal Uniper HSSE environment. This design decision was necessary because Process Safety incidents, general HSSE incidents, and near-hit/observation reports were originally stored in separate source tables with partially overlapping but not fully identical schemas. To build a consistent experimental corpus, these tables were harmonised and merged into a single master dataset before any language filtering, class formulation, or modelling was carried out.

\subsection{Source Tables and Initial Data Integration}

Four raw JSONL source tables were used to construct the master dataset. Each table represented a different operational reporting stream within the broader Uniper HSSE ecosystem.

\begin{table}[htbp]
\centering
\caption{Raw source tables used to construct the Uniper HSSE master dataset.}
\label{tab:source_tables}
\small
\renewcommand{\arraystretch}{1.2}
\begin{tabular}{p{4.8cm}p{3.4cm}p{1.8cm}p{5.0cm}}
\toprule
\textbf{Source table / file} & \textbf{Reporting scope} & \textbf{Rows} & \textbf{Purpose in this study} \\
\midrule
\texttt{14K\_Reports RW\_ ACTUALS\_PS.jsonl} 
& Process Safety reports 
& 14,432 
& Preserved dedicated Process Safety cases for domain-relevant incident coverage. \\

\texttt{28k\_Reports RW\_ ACTUALS\_ALL.jsonl} 
& General HSSE incident reports 
& 28,323 
& Added broader operational incident coverage beyond PS-only records. \\

\texttt{30K\_Reports RW\_ ACTUALS\_ALL\_v2.0.jsonl} 
& Updated general incident reports 
& 30,609 
& Added supplementary incident data after schema harmonisation. \\

\texttt{700k\_reports RW\_ ACTUALS\_NH\_OBS.jsonl} 
& Near hits and observations 
& 717,041 
& Added large-scale near-hit and observation reporting context. \\

\midrule
\textbf{Total before deduplication} 
& 
& \textbf{790,405} 
& Raw integrated input to the master dataset. \\
\bottomrule
\end{tabular}
\end{table}

The first source table contained dedicated Process Safety records, while the second and third tables contained broader incident reports from the general HSSE reporting workflow. The fourth table, which was substantially larger than the others, contained near-hit and observation records. These four tables were concatenated to form the initial master dataset. Before deduplication, the combined table contained 790,405 rows.

\subsection{Schema Harmonisation and Column Alignment}

Because the four source tables were not fully identical in structure, a schema harmonisation step was required before merging. This was particularly important for ensuring that later preprocessing, embedding generation, and classification experiments operated on a consistent feature space.

To harmonise the source tables prior to integration, several schema adjustments were applied. In the 700k near-hit and observation table, the field \texttt{DATETIME} was renamed to \texttt{CASE\_OCCURENCE\_DATE}. In the updated 30k incident table, \texttt{EMPLOYMENT\_CATEGORY} was renamed to \texttt{COMPANY\_INVOLVED\_TYPE}, and \texttt{HAZARD\_PHYSICAL\_SECURITY\_EVENT} was renamed to \texttt{HAZARD}. In addition, the column \texttt{IMM\_ACTION\_TAKEN\_RECOM} was added to the 30k table because it was missing there but present in the other reporting streams. Several columns that did not fit the common target schema were removed from the 30k table, namely \texttt{PERSONAL\_INJURIES}, \texttt{ACTUAL\_WORKDAYS\_LTA}, \texttt{COMPANY\_NAME}, \texttt{COMPANY\_TYPE}, \texttt{CASE\_APPROVED\_DATE}, \texttt{OUTSIDE\_LEGAL\_INFLUENCE}, and \texttt{RETURNED\_TO\_WORK\_DATE}.

After harmonisation, all four source tables were aligned to a common schema containing 36 columns. This standardised structure ensured that the resulting master dataset remained internally consistent despite being assembled from four distinct operational reporting streams.

\begin{table}[htbp]
\centering
\caption{Common harmonised schema of the integrated Uniper HSSE master dataset.}
\label{tab:harmonised_schema}
\small
\renewcommand{\arraystretch}{1.15}
\begin{tabular}{p{0.47\textwidth}p{0.47\textwidth}}
\toprule
\textbf{Schema fields (1--18)} & \textbf{Schema fields (19--36)} \\
\midrule
\texttt{CASENO} &
\texttt{CASE\_TYPE} \\

\texttt{COMPANY} &
\texttt{CASE\_SEVERITY} \\

\texttt{FUNCTIONAL\_GROUP} &
\texttt{CASE\_DESCRIPTION} \\

\texttt{FUNCTION} &
\texttt{STATUS} \\

\texttt{FUNCTIONAL\_AREA} &
\texttt{HAZARD} \\

\texttt{FUNCTIONAL\_LOCATION} &
\texttt{IMM\_ACTION\_TAKEN\_RECOM} \\

\texttt{FUNCTIONAL\_SUB\_LOCATION} &
\texttt{FULL\_INVESTIGATION\_DONE} \\

\texttt{LOCATION\_SID} &
\texttt{CREATED\_DATE} \\

\texttt{LOCATION\_SHORT} &
\texttt{MODIFIED\_DATE} \\

\texttt{COUNTRY\_SHORT} &
\texttt{APPROVED\_WITHIN\_DEADLINE} \\

\texttt{SL\_COUNTRY} &
\texttt{CASE\_CLOSED\_DATE} \\

\texttt{SL\_LOCATION\_LVL\_1} &
\texttt{CASES\_NO\_OF\_REGISTRATIONS} \\

\texttt{SL\_LOCATION\_LVL\_2} &
\texttt{VALID\_FROM} \\

\texttt{SL\_LOCATION\_LVL\_3} &
\texttt{VALID\_TO} \\

\texttt{SL\_LOCATION\_LVL\_4} &
\texttt{POTENTIAL\_SEV\_LEVEL} \\

\texttt{CASE\_OCCURENCE\_DATE} &
\texttt{RISK\_AREA} \\

\texttt{TITLE} &
\texttt{LEARNINGS\_ACTUAL\_SEVERITY} \\

\texttt{COMPANY\_INVOLVED\_TYPE} &
\texttt{LEARNINGS\_POTENTIAL\_RISK} \\
\bottomrule
\end{tabular}
\end{table}

\subsection{Deduplication and Cleaning of the Master Dataset}

After concatenation, duplicate handling was necessary because the same case could appear multiple times across reporting updates. Duplicates were identified using the unique case identifier \texttt{CASENO}. For each duplicated case, the most recent record was retained by sorting the dataset in descending order of \texttt{MODIFIED\_DATE} and keeping only the first occurrence. This step reduced the dataset from 790,405 rows to 34,926 unique incident records.

A second cleaning step removed records with missing or empty values in fields that were considered essential for the downstream NLP and classification pipeline. Specifically, rows missing \texttt{TITLE}, \texttt{CASE\_DESCRIPTION}, \texttt{CASE\_TYPE}, or \texttt{HAZARD} were removed. After this filtering step, the final cleaned master dataset contained 34,576 records.

\begin{table}[h!]
\centering
\caption{Dataset size reduction during construction of the master dataset.}
\label{tab:dataset_reduction}
\begin{tabular}{p{9cm}r}
\toprule
\textbf{Processing step} & \textbf{Records remaining} \\
\midrule
Raw merged dataset (4 source tables combined) & 790,405 \\
After deduplication on \texttt{CASENO} (keeping latest \texttt{MODIFIED\_DATE}) & 34,926 \\
After removing missing \texttt{TITLE} & 34,925 \\
After removing missing \texttt{CASE\_DESCRIPTION} & 34,925 \\
After removing missing \texttt{CASE\_TYPE} & 34,576 \\
After removing missing \texttt{HAZARD} & 34,576 \\
\midrule
\textbf{Final cleaned master dataset} & \textbf{34,576} \\
\bottomrule
\end{tabular}
\end{table}

The strong reduction from 790,405 raw rows to 34,576 cleaned records is not a data loss in the usual sense, but primarily the result of deduplication across repeated case updates and the removal of records that lacked the minimum textual and categorical information required for reliable supervised learning.

\subsection{Distribution of Cases by Country and Case Type}

The cleaned master dataset is highly unevenly distributed across countries. Germany and Sweden account for the largest shares of the final dataset, followed by the UK and the Netherlands, while several locations such as UAE, International, Russia, and Poland contribute only a very small number of records. This geographical imbalance is important because it directly affects the linguistic composition of the data and the robustness of country-specific experiments.

\begin{figure}[htbp]
    \centering
    \includegraphics[width=0.85\textwidth]{Results/_iteration_0/Number of Cases by Country.png}
    \caption{Distribution of cleaned incident records by country.}
    \label{fig:cases_by_country}
\end{figure}

\begin{figure}[htbp]
    \centering
    \includegraphics[width=0.85\textwidth]{Results/_iteration_0/Number of Cases by Case Type.png}
    \caption{Distribution of cleaned incident records by case type.}
    \label{fig:cases_by_case_type}
\end{figure}

Figure~\ref{fig:cases_by_country} shows that Germany contributes 14,951 records and Sweden 12,224 records, making these two countries the dominant sources of incident reports in the cleaned master dataset. The UK contributes 4,673 records and the Netherlands 1,990 records, while all remaining country groups contribute comparatively few cases. This imbalance has methodological consequences for multilingual modelling, since the effective representation of languages and reporting styles is not uniform across the corpus.

Figure~\ref{fig:cases_by_case_type} shows that the case-type distribution is also strongly imbalanced. The largest class is \textit{Safety} with 25,164 records, followed by \textit{Process Safety} with 4,561 and \textit{Environment} with 2,768. The remaining categories are substantially smaller: \textit{Asset and Reputation damage/loss} (1,333), \textit{Operational loss} (456), \textit{Information Security} (186), \textit{Physical Security} (100), and \textit{Not classified} (8). This imbalance provides an important methodological justification for later binary reformulation, class balancing, and macro-F1-based evaluation.

\subsection{Relationship to the Iteration~0 English Subsets}

The cleaned master dataset served as the source pool for the English-language subsets used in Iteration~0. From this master dataset, three alternative English subsets were created using different subset construction strategies: a manually curated English subset, a \texttt{langdetect}-based subset, and a \texttt{lingua}-based subset. According to the current notebook output, these yielded 4,713, 6,010, and 6,222 records, respectively. The manually curated English subset was later selected as the official Iteration~0 baseline because it achieved the strongest macro-F1 score in the saved BERT \& SVM result artefacts.

\subsection{External Datasets Explored but Not Included}

In addition to the internal Uniper HSSE reporting data, several external datasets were explored during the early scoping phase of the thesis. These included OSHA IMIS narratives, MSHA injury narratives, Dutch health complaint datasets, and UK Environment Agency CICS summaries. Although these datasets were initially considered as potential supplements for enlarging the corpus or enriching category diversity, they were ultimately not included in the final experimental dataset.

The main reason for exclusion was insufficient alignment with the structure and business context of the Uniper HSSE reporting environment. The external datasets differ substantially from the Uniper corpus in at least four ways. First, their label systems and incident taxonomies are defined for different regulatory or organisational settings, which makes direct category mapping unreliable. Second, the narrative style and reporting granularity differ from the internal Uniper HSSE reports, reducing comparability at the text level. Third, the domain emphasis of these external datasets is not fully consistent with the industrial Process Safety framing of this study; some are more strongly oriented toward occupational injury, public complaints, or environmental compliance narratives. Fourth, integrating them would have introduced substantial domain shift, which would have weakened the internal validity of the comparative experiments.

For these reasons, the final thesis experiments were restricted to the harmonised Uniper HSSE master dataset and the language-specific subsets derived from it. This ensured that the classification task remained closely aligned with the actual business case, the internal reporting taxonomy, and the multilingual incident-reporting environment relevant to the thesis objectives.

%------------------Table-------------------------
\begin{table}[h!]
\centering
\caption{Examples of raw incident reports illustrating data quality issues.}
\label{tab:rawdata}
\begin{footnotesize}
\begin{tabular}{p{1.3cm} p{8.2cm} p{6.5cm}}
\toprule
\textbf{Location} & \textbf{Title and Description} & \textbf{Potential Concerns} \\
\midrule
NL & \textbf{Title:} Overflow buffer B\newline \textbf{Description:} Buffer B is over gelopen door een bedieningsfout op buffer A. Tijdens het overzetten van de buffers omdat buffer B vol was is de retourwaterafsluiter van buffer A open gezet terwijl de aanvoerwaterafsluiter dicht stond. Dit heeft geresulteerd in het overlopen van buffer B met heet water (ca.~96$^\circ$C). Doordat dit een stoompluim op het dak van de buffer gaf, was de brandweer verwittigd. & Title is in English, but the description is in Dutch. \\
\addlinespace
NL & \textbf{Title:} Object falling from height -- GT3 enclosure\newline \textbf{Description:} While working with the chain hoist to remove parts of the turbine, the hoist's chain came into contact with an electrical distribution box. The distribution box was sealed with an iron cover tied down with strings. Due to the contact, the cover became loose and fell down where employees were working. The cover narrowly missed hitting a worker at ground level. & Report from NL but written in English. \\
\addlinespace
UK & \textbf{Title:} Environmental agency not correctly informed of GT12 fire at Killingholme power station\newline \textbf{Description:} Following an uncontrolled fire at the Killingholme power station, the Environment Agency was not informed. Site procedures were not clear about the requirement. & Grammatical errors in the description. \\
\addlinespace
NL & \textbf{Title:} Geen BMC meer op Roca\newline \textbf{Description:} Door storing geen BMC op Roca. & Description too short. Recommended description length is at least 50--100 words. \\
\addlinespace
NO & \textbf{Title:} Se över allmän belysning i DC.\newline \textbf{Description:} Se över allmän belysning i DC. & Description too short. Similar issues observed in reports from Sweden (due to data protection regulations). \\
\addlinespace
UK & \textbf{Title:} Steam Turbine Safety System Transmitter changed without Modification Process Followed\newline \textbf{Description:} One of the condenser level transmitters was replaced with a different make and model without following the proper modification process. This transmitter is part of the Steam Turbine SIL (Safety Integrity Level) rated protection functions. The modification process is intended to ensure any new transmitter meets the required reliability calculations from the Safety Manual and maintains the integrity of the protective function. There is a risk that the safety system’s integrity has been compromised without due diligence. & Context-specific abbreviations (e.g., ``SIL'' for Safety Integrity Level). \\
\addlinespace
NIL & \textbf{Title:} Safety Walk Betrieb\newline \textbf{Description:} 47319 An dem Kran in der benachbarten Halle der Wasseraufbereitung ist zu prüfen, ob die Bedieneinheit korrekt positioniert ist. & Location missing in dataset. \\
\bottomrule
\end{tabular}
\end{footnotesize}
\end{table}

\section{Data Preprocessing}

Before modelling, several preprocessing steps were applied to improve data quality and prepare the incident narratives for downstream machine learning experiments:

\begin{itemize}
    \item \textit{Data Anonymization:} Sensitive or business-critical information was removed or masked where necessary in order to protect confidentiality and reduce the risk of exposing internal operational details.
    
    \item \textit{Cleaning and Class Balancing:} Standard cleaning operations were applied, including the removal of invalid or incomplete entries, the treatment of missing values, and the normalisation of inconsistent textual content where appropriate. Because the classification task exhibits substantial class imbalance, class balancing strategies were later considered as part of the modelling pipeline.
    
    \item \textit{Tokenization and Lemmatization:} Text was segmented into tokens and standardised for NLP processing. Where relevant, words were reduced to their base forms to minimise variation caused by inflectional differences.
    
    \item \textit{Dataset Splitting:} The cleaned dataset was divided into training and test subsets in order to enable reproducible model training and evaluation.
\end{itemize}

\newpage

Table~\ref{tab:rawdata} presents representative examples of raw incident reports drawn from the industrial dataset. These examples highlight common data quality challenges encountered in real-world reporting, including multilingual descriptions, overly brief or incomplete narratives, grammatical errors, and the use of context-specific abbreviations. Together with the dataset-level imbalances shown in Figures~\ref{fig:cases_by_country} and~\ref{fig:cases_by_case_type}, these issues illustrate the complexity of preparing incident reports for automated analysis and reinforce the need for robust preprocessing before reliable machine learning models can be trained.

\section{Model Development}

The model development strategy of this thesis follows an iterative experimental design. Rather than directly moving to advanced multilingual transformers or large language models, the research begins with a controlled baseline implementation and then progressively introduces additional complexity in later iterations. This approach makes it possible to identify the effect of individual methodological choices, isolate the contribution of specific components, and build a reproducible performance trajectory from a simple baseline to more advanced architectures.

The development process starts with a baseline binary classification setting in which incident reports are classified as either \textit{Process Safety} or \textit{Non-Process Safety}. The binary formulation was chosen intentionally, because Process Safety incidents represent the most business-critical class and therefore provide a focused first benchmark before moving to more complex multi-class and multilingual settings.

Iteration~0 serves as the experimental foundation of the thesis. It establishes the first end-to-end implementation pipeline, starting from cleaned and merged incident reports, continuing through English subset construction, exploratory data analysis, text embedding generation, and classifier training, and ending with an evaluation of three alternative English-language subset construction strategies. The strongest class-balanced strategy is then adopted as the official baseline for subsequent iterations.

\section{Iteration 0: Baseline Implementation}

\subsection{Objective and Experimental Scope}

The purpose of Iteration~0 was to establish a reproducible baseline for Process Safety classification and to evaluate whether the method used to construct the English-language subset materially influences downstream model performance. Rather than relying on a single English subset, three alternative strategies were implemented and compared within an otherwise fixed modelling pipeline. The first strategy used a \textbf{manual English subset}, created by combining country-level files that were considered appropriate for an English-language baseline. The second strategy used a \textbf{\texttt{langdetect}-based English subset}, generated automatically through probabilistic language detection. The third strategy used a \textbf{\texttt{lingua}-based English subset}, generated automatically through a restricted high-speed language detector. After construction, all three subsets were processed identically. This experimental design made it possible to isolate the effect of English subset construction while keeping feature extraction and classification constant.

\subsection{Integration with the Master Dataset}

Iteration~0 operates on the cleaned master dataset introduced in the Proposed Dataset section. The full merged and cleaned master dataset contains 34,576 records after deduplication and removal of entries with missing \texttt{TITLE}, \texttt{CASE\_DESCRIPTION}, \texttt{CASE\_TYPE}, or \texttt{HAZARD}. From this cleaned dataset, English-language subsets were derived using three alternative strategies.

The manually constructed English subset was created by merging the country-level files corresponding to the United Kingdom, the United Arab Emirates, Russia, and Poland. In total, this subset comprised 4,673 records from the UK, 29 records from the UAE, 10 records from Russia, and 1 record from Poland. The resulting aggregation produced a manually curated English-language dataset of 4,713 records intended to support the baseline classification experiments.

The two automatic language detection approaches were then applied to the entire cleaned master dataset using the concatenated \texttt{TITLE} + \texttt{CASE\_DESCRIPTION} text.

For the \texttt{langdetect} approach, the language detector was restricted to four target languages, namely English, German, Swedish, and Dutch. Any cases for which the detector did not produce a confident assignment within this predefined set were mapped to the category \texttt{Unknown}. The resulting language distribution in the current notebook output showed that German was the most frequent language with 15,188 records, followed by Swedish with 9,794 records and English with 5,989 records. Dutch accounted for 1,990 records, while 1,615 records remained in the \texttt{Unknown} category. This distribution shows that the \texttt{langdetect} procedure captured substantial multilingual variation in the dataset while also preserving a separate category for uncertain or ambiguous predictions.

For the \texttt{lingua} approach, the detector was applied to the cleaned master dataset in order to construct a second automatically derived English subset. In the current notebook output, this procedure yields an English subset containing 6,222 records, making it the largest of the three English subsets used in Iteration~0. This larger automatically derived subset provides an additional point of comparison for assessing whether language-detection strategy influences downstream classification performance.

Thus, the three English subsets evaluated in Iteration~0 differed in size. The manually constructed subset contained 4,713 records, while the \texttt{langdetect}-based subset comprised 6,010 records and the \texttt{lingua}-based subset contained 6,222 records. These differences reflect the varying behaviour of manual and automated language selection strategies and are relevant for interpreting subsequent classification performance across the three experimental settings.

\subsection{Exploratory Text Analysis}

Before modelling, an exploratory analysis of text length and lexical sparsity was performed in order to understand the quality of the available report narratives and the potential limitations for transformer-based embedding generation.

\subsubsection{Character Count Analysis}

For the full cleaned master dataset, the character-length statistics indicate clear differences between the title, description, and combined text fields. Incident titles were comparatively short, with a minimum length of 3 characters, a maximum of 662 characters, a mean length of 54.7086 characters, and a median of 46 characters. By contrast, the case descriptions exhibited substantially greater variation, ranging from 1 character to 6,583 characters, with a mean of 140.0750 characters and a median of 82 characters. When the title and description fields were combined, the resulting text length ranged from 6 to 6,708 characters, with a mean of 194.7836 characters and a median of 140 characters. Overall, these statistics show that incident titles are generally concise, whereas descriptions vary considerably in length and display a pronounced right-skewed distribution, indicating the presence of a subset of particularly detailed incident reports.


\subsubsection{Word Count Analysis}

The word-count statistics further confirm the concise nature of many incident reports in the cleaned master dataset. For the title field, the number of words ranged from a minimum of 1 to a maximum of 92, with a mean of 7.9015 words and a median of 6 words. The description field was notably longer and more variable, ranging from 1 to 953 words, with a mean of 21.2266 words and a median of 12 words. When both title and description were combined, the resulting text length ranged from 2 to 969 words, with a mean of 29.1281 words and a median of 20 words. Taken together, these results indicate that a substantial proportion of incident reports are relatively short, particularly at the title level, which is an important consideration for embedding quality, contextual representation, and subsequent classification performance.

\subsubsection{Low Word Count Analysis}

A dedicated low-word-count analysis was performed to estimate the proportion of narratives with very limited semantic content.

For word count thresholds of 3, 5, and 10 words, the following proportions were observed:

\begin{table}[h!]
\centering
\caption{Low word count analysis for the cleaned master dataset.}
\label{tab:iter0_low_word_count}
\begin{tabular}{lrrr}
\toprule
\textbf{Threshold} & \textbf{Title} & \textbf{Description} & \textbf{Combined Text} \\
\midrule
$\leq 3$ words  & 6,374 (18.4348\%) & 3,263 (9.4372\%) & 267 (0.7722\%) \\
$\leq 5$ words  & 13,952 (40.3517\%) & 7,189 (20.7919\%) & 1,230 (3.5574\%) \\
$\leq 10$ words & 26,613 (76.9696\%) & 15,432 (44.6321\%) & 7,288 (21.0782\%) \\
\bottomrule
\end{tabular}
\end{table}

The very high proportion of short titles and relatively short descriptions highlights an important limitation of the raw data: some reports provide only weak semantic signal for contextual embedding models. This observation is relevant for interpreting baseline performance and motivates later experiments with alternative subset construction strategies and more robust multilingual representations.

\begin{figure}[htbp]
    \centering
    \includegraphics[width=0.95\textwidth]{Results/_iteration_0/Word Count Distribution - (Log Scale).png}
    \caption{Word count distribution for titles and descriptions on a logarithmic scale.}
    \label{fig:iter0_word_count_log}
\end{figure}

Figure~\ref{fig:iter0_word_count_log} further illustrates the strong right-skew in the text-length distribution. Most titles are very short, while description lengths vary widely and include a long tail of highly detailed reports.

\subsection{Binary Classification Formulation}

After constructing the three English subsets, each record was reformulated as a binary classification instance. Let $y_i$ denote the target label of report $i$. The label was defined as:

\begin{equation}
y_i =
\begin{cases}
1, & \text{if } \texttt{CASE\_TYPE}_i = \text{Process Safety},\\[0.5em]
0, & \text{otherwise}.
\end{cases}
\end{equation}

This binary formulation collapses all non-Process-Safety categories into a single negative class and yields the first operational baseline for the business problem.

The class distributions across the three English subsets reveal a consistent imbalance between Process Safety and Non-Process Safety records. In the manually constructed English subset, the dataset comprised 1,409 Process Safety cases and 3,304 Non-Process Safety cases. For the \texttt{langdetect}-based subset, the class distribution consisted of 1,528 Process Safety records and 4,461 Non-Process Safety records. Similarly, the \texttt{lingua}-based subset contained 1,528 Process Safety records and 4,694 Non-Process Safety records. Overall, all three subsets were characterised by a substantially larger number of Non-Process Safety cases, highlighting the class imbalance that must be taken into account during model training and evaluation.

The combined text feature for each record was created as:

\begin{equation}
\texttt{text\_features}_i = \texttt{TITLE}_i \oplus \text{``.''} \oplus \texttt{CASE\_DESCRIPTION}_i,
\end{equation}

where $\oplus$ denotes string concatenation. This ensures that both the concise title and the richer narrative description contribute to the document representation.

\subsection{Mathematical Foundations of Document Embeddings}
After concatenating the title and case description, each incident report is represented as a text sequence
\begin{equation}
x_i = (t_{i1}, t_{i2}, \dots, t_{im_i}),
\end{equation}
where $m_i$ denotes the number of tokens in report $i$. The objective of the embedding step is to transform this variable-length token sequence into a fixed-dimensional continuous vector representation suitable for downstream classification.

Let $\mathcal{V}$ denote the tokenizer vocabulary and let
\begin{equation}
\phi : \mathcal{V} \rightarrow \mathbb{R}^{d}
\end{equation}
be an embedding function that maps each token to a $d$-dimensional vector. For a tokenised report $x_i$, the embedding layer produces an ordered sequence of vector representations
\begin{equation}
\mathbf{H}_i^{(0)} =
\left[
\phi(t_{i1}), \phi(t_{i2}), \dots, \phi(t_{im_i})
\right]
\in \mathbb{R}^{m_i \times d}.
\end{equation}

In transformer-based models, these initial token embeddings are not used directly as final representations. Instead, they are iteratively contextualised through multiple self-attention and feed-forward layers. If $L$ denotes the number of transformer layers, the final hidden representation can be written as
\begin{equation}
\mathbf{H}_i^{(L)} = \mathrm{Transformer}\!\left(\mathbf{H}_i^{(0)}\right),
\end{equation}
where
\begin{equation}
\mathbf{H}_i^{(L)} =
\left[
\mathbf{h}_{i1}^{(L)},
\mathbf{h}_{i2}^{(L)},
\dots,
\mathbf{h}_{im_i}^{(L)}
\right].
\end{equation}

To apply a conventional classifier, the sequence of contextual token vectors must be reduced to a single document representation. In general form, this can be expressed as
\begin{equation}
\mathbf{z}_i = g\!\left(\mathbf{H}_i^{(L)}\right),
\end{equation}
where $g(\cdot)$ denotes a pooling or selection function. Depending on the encoder architecture, this may correspond to mean pooling across all tokens, max pooling, or the representation of a designated classification token. Thus, each incident report $x_i$ is transformed into a fixed-dimensional vector
\begin{equation}
\mathbf{z}_i \in \mathbb{R}^{d},
\end{equation}
which serves as the feature vector for the downstream classifier:
\begin{equation}
\hat{y}_i = f_{\mathrm{clf}}(\mathbf{z}_i).
\end{equation}

This formulation is particularly suitable for the present study because incident reports vary considerably in length, lexical richness, and language. The embedding process therefore acts as a representation-learning step that converts heterogeneous free-text safety narratives into a common semantic vector space in which similar incidents can be positioned closer together and subsequently classified more effectively.

\subsection{BERT Embedding Generation}

For all three English subsets, contextual document representations were generated using the pre-trained \texttt{bert-base-uncased} transformer model. The tokenizer and model were loaded from Hugging Face Transformers, and the model was used in evaluation mode.

The implementation employed the \texttt{bert-base-uncased} model, using the \texttt{[CLS]} token representation from the final hidden layer as the pooling strategy to obtain a fixed-length document embedding. This produced embeddings of dimensionality 768. During embedding generation, a batch size of 8 was used, and the maximum input sequence length was set to 512 tokens. The recorded notebook execution was performed on a CPU.

Formally, the document embedding was defined as:

\begin{equation}
\mathbf{h}_{\text{[CLS]}} = \text{BERT}\big([\texttt{[CLS]}] \oplus \text{tokens} \oplus [\texttt{[SEP]}]\big)[0] \in \mathbb{R}^{768}.
\end{equation}

The generated embedding matrices reflected the size of each English subset while preserving a fixed embedding dimensionality of 768. Specifically, the manually constructed subset produced an embedding matrix of size $(4713, 768)$, the \texttt{langdetect}-based subset yielded a matrix of size $(6010, 768)$, and the \texttt{lingua}-based subset produced a matrix of size $(6222, 768)$. These dimensions confirm that each incident report was represented by a 768-dimensional dense vector derived from the underlying BERT model.

Embeddings and labels were saved as \texttt{.npy} files for reproducibility and reuse in later experiments.

\subsection{Classifier Training and Evaluation Pipeline}

For each English subset, an identical train--test pipeline was applied in order to ensure a fair comparison across the three dataset construction strategies. First, the data were divided using an 80/20 stratified train--test split in order to preserve the class distribution in both partitions. Next, the embedding features were standardised using \texttt{StandardScaler}. The standardised embeddings were then used to train an SVM classifier. Finally, model performance was evaluated using accuracy, macro-F1, weighted-F1, detailed classification reports, and confusion matrices.

%%%%%%% The following needs to be corrected. Kept as a place holder for now. %%%

The train--test split counts were:

\begin{table}[h!]
\centering
\caption{Train--test split statistics for the three Iteration~0 English subsets.}
\label{tab:iter0_train_test_split}
\begin{tabular}{lrrrr}
\toprule
\textbf{Subset} & \textbf{Train} & \textbf{Test} & \textbf{Train PS / Non-PS} & \textbf{Test PS / Non-PS} \\
\midrule
manual     & 3,770 &   943 & 1,127 / 2,643 & 282 / 661 \\
langdetect & 4,808 & 1,202 & 1,222 / 3,586 & 306 / 896 \\
lingua     & 4,977 & 1,245 & 1,222 / 3,755 & 306 / 939 \\
\bottomrule
\end{tabular}
\end{table}

%%%%%%%%%% The above needs to be corrected %%%%%%%%%%%

\subsubsection{Actual Implemented Classifier}

The notebook implementation employed \texttt{sklearn.svm.SVC} as the classifier, using an \texttt{rbf} kernel and a regularisation parameter of $C = 1.0$. To account for class imbalance, the model was configured with \texttt{class\_weight='balanced'}. In addition, a fixed random state of 42 was used to ensure reproducibility of the experimental procedure.

Thus, the implemented baseline classifier is an \textbf{RBF-kernel SVM}, not a linear SVM. Since the notebook also scales the embeddings before training, the effective Iteration~0 classifier is more precisely described as:

\begin{quote}
\textit{BERT-base-uncased [CLS] embeddings + StandardScaler + RBF SVM}
\end{quote}

This is the formulation that should be reflected consistently throughout the thesis if the notebook implementation is treated as authoritative.


\subsection{Iteration~0 Results}

Table~\ref{tab:iter0_comparison_summary} summarises the final performance of the three English subset construction strategies under the common BERT \& SVM pipeline.

\begin{table}[h!]
\centering
\caption{Comparison of the three English subset construction strategies in Iteration~0.}
\label{tab:iter0_comparison_summary}
\begin{tabular}{lrrrr}
\toprule
\textbf{Subset} & \textbf{Total} & \textbf{Accuracy} & \textbf{Macro-average F1} & \textbf{Weighted-F1} \\
\midrule
manual     & 4,713 & 0.8112 & \textbf{0.7825} & 0.8143 \\
langdetect & 6,010 & 0.8128 & 0.7648 & 0.8170 \\
lingua     & 6,222 & 0.8185 & 0.7714 & \textbf{0.8241} \\
\bottomrule
\end{tabular}
\end{table}

The results show that the three subset construction strategies lead to materially different performance profiles. The manually constructed subset achieves the highest macro-average F1-score and therefore provides the strongest class-balanced baseline. By contrast, the \texttt{lingua}-based subset attains the highest accuracy and weighted-F1, while the \texttt{langdetect}-based subset performs below both alternatives on macro-average F1.

\subsubsection{Detailed Classification Reports}

\begin{table}[h!]
\centering
\caption{Per-class performance across the three English subset construction strategies.}
\label{tab:iter0_per_class_metrics}
\begin{tabular}{llccccc}
\toprule
\textbf{Subset} & \textbf{Class} & \textbf{Precision} & \textbf{Recall} & \textbf{F1-score} & \textbf{Macro-average F1} & \textbf{Support} \\
\midrule
\multirow{2}{*}{manual}
    & Non-Process Safety & 0.8864 & 0.8381 & 0.8616 & \multirow{2}{*}{0.7825} & 661 \\
    & Process Safety     & 0.6635 & 0.7482 & 0.7033 &  & 282 \\
\midrule
\multirow{2}{*}{langdetect}
    & Non-Process Safety & 0.8952 & 0.8482 & 0.8711 & \multirow{2}{*}{0.7648} & 896 \\
    & Process Safety     & 0.6147 & 0.7092 & 0.6586 &  & 306 \\
\midrule
\multirow{2}{*}{lingua}
    & Non-Process Safety & 0.9093 & 0.8435 & 0.8751 & \multirow{2}{*}{0.7714} & 939 \\
    & Process Safety     & 0.6070 & 0.7418 & 0.6676 &  & 306 \\
\bottomrule
\end{tabular}
\end{table}

These results indicate that the strongest subset is not determined solely by raw accuracy. The manual subset achieves the highest macro-average F1-score and the strongest Process Safety F1-score, which makes it the most appropriate official baseline when class-balanced performance is prioritised. Although the \texttt{lingua}-based subset achieves slightly higher overall accuracy and weighted-F1, its macro-average F1 remains below that of the manual subset.

\subsubsection{Confusion Matrix Analysis}

The confusion matrices for the three subset construction strategies are shown in Figures~\ref{fig:iter0_confusion_matrices_a} and~\ref{fig:iter0_confusion_matrices_b}. They provide a direct view of the trade-off between false positives and false negatives across the three English baselines.

\begin{figure}[p]
    \centering

    \begin{subfigure}[b]{0.82\textwidth}
        \centering
        \includegraphics[width=\textwidth]{Results/_iteration_0/confusion_matrix_manual.png}
        \caption{Manual English subset: TN=554, FP=107, FN=71, TP=211.}
        \label{fig:iter0_cm_manual}
    \end{subfigure}

    \vspace{1em}

    \begin{subfigure}[b]{0.82\textwidth}
        \centering
        \includegraphics[width=\textwidth]{Results/_iteration_0/confusion_matrix_langdetect.png}
        \caption{Langdetect English subset: TN=760, FP=136, FN=89, TP=217.}
        \label{fig:iter0_cm_langdetect}
    \end{subfigure}

    \caption{Confusion matrices for the manual and \texttt{langdetect}-based English subset construction strategies in Iteration~0.}
    \label{fig:iter0_confusion_matrices_a}
\end{figure}

\begin{figure}[p]
    \centering

    \begin{subfigure}[b]{0.82\textwidth}
        \centering
        \includegraphics[width=\textwidth]{Results/_iteration_0/confusion_matrix_lingua.png}
        \caption{Lingua English subset: TN=792, FP=147, FN=79, TP=227.}
        \label{fig:iter0_cm_lingua}
    \end{subfigure}

    \caption{Confusion matrix for the \texttt{lingua}-based English subset construction strategy in Iteration~0.}
    \label{fig:iter0_confusion_matrices_b}
\end{figure}

Figures~\ref{fig:iter0_cm_manual}--\ref{fig:iter0_cm_lingua} show the confusion matrices obtained for the three English subset construction strategies in Iteration~0. Presenting these matrices across two dedicated float pages improves readability and ensures that all axis labels, cell counts, and colour scales remain fully visible in the compiled thesis.

For the \texttt{manual} subset, the classifier correctly identifies 554 Non-Process Safety and 211 Process Safety reports, while producing 107 false positives and 71 false negatives.

For the \texttt{langdetect} subset, the classifier correctly identifies 760 Non-Process Safety and 217 Process Safety reports, while producing 136 false positives and 89 false negatives.

For the \texttt{lingua} subset, the classifier correctly identifies 792 Non-Process Safety and 227 Process Safety reports, while producing 147 false positives and 79 false negatives.

\subsection{Selection of the Official Iteration~0 Baseline}

According to the saved Iteration~0 classification reports, the manually constructed English subset represents the strongest baseline when macro-average F1 is used as the primary model selection criterion. This configuration achieved an accuracy of 0.8112, a macro-F1 score of 0.7825, and a weighted-F1 score of 0.8143. Although the \texttt{lingua}-based subset obtained slightly higher accuracy and weighted-F1, the manual subset provided the best class-balanced performance and is therefore retained as the official Iteration~0 baseline. The corresponding baseline configuration consists of a manually constructed English subset, encoded using BERT-base-uncased \texttt{[CLS]} embeddings, standardised with \texttt{StandardScaler}, and classified using an RBF-kernel SVM.

\subsection{Gap to Research Target}

Using the official Iteration~0 baseline, the gap to the RQ1 target of 0.85 is:

\begin{equation}
0.85 - 0.7825 = 0.0675.
\end{equation}

Thus, the required improvement from the official Iteration~0 baseline is 6.75 percentage points in macro-average F1.

\begin{table}[h!]
\centering
\caption{Gap between the Iteration~0 subset strategies and the RQ1 target.}
\label{tab:iter0_gap_analysis}
\begin{tabular}{lccc}
\toprule
\textbf{Subset} & \textbf{Macro-F1} & \textbf{RQ1 Target} & \textbf{Gap} \\
\midrule
manual     & 0.7825 & 0.85 & 0.0675 \\
langdetect & 0.7648 & 0.85 & 0.0852 \\
lingua     & 0.7714 & 0.85 & 0.0786 \\
\bottomrule
\end{tabular}
\end{table}

\subsection{Summary of Iteration~0}

Iteration~0 establishes the first fully specified experimental pipeline of the thesis and provides the initial empirical reference point for subsequent model development. The results show, first, that the method used to construct the English-language subset has a measurable impact on downstream classification performance. Second, they indicate that short and low-information incident narratives constitute an important data-quality limitation for embedding-based classification, as such reports provide only restricted semantic context for representation learning. Third, the experiments demonstrate that BERT-based document embeddings combined with an SVM classifier already yield a strong baseline for the binary Process Safety classification task. Finally, among the three English subset construction strategies evaluated in Iteration~0, the manually constructed English subset provides the strongest class-balanced baseline when assessed using macro-average F1. Taken together, this baseline implementation serves as the methodological reference point for all subsequent iterations, which progressively extend the experimental framework toward stronger transformer architectures, multilingual modelling, improved class-balancing strategies, ensemble methods, and LLM-based approaches.

\subsection{Iteration 1: BERT-base-cased Baseline Implementation}

\subsubsection{Objective}
Establish a country-stratified baseline for process safety incident classification using BERT-base-cased embeddings across the cleaned master dataset and its major language/country subsets.

\subsubsection{Implementation Steps}

\paragraph{1. Data Consolidation and Preparation}
Iteration~1 did not construct a new master dataset from raw reporting streams. Instead, it operated on the cleaned master dataset established in Iteration~0. This cleaned master dataset contained 34,576 incident records after deduplication and the removal of entries with missing \texttt{TITLE}, \texttt{CASE\_DESCRIPTION}, \texttt{CASE\_TYPE}, or \texttt{HAZARD}.

For country-stratified evaluation, six datasets were processed. These comprised the Germany subset with 14,951 records, the UK subset with 4,673 records, the manually curated English subset with 4,713 records, the Sweden subset with 12,224 records, the Netherlands subset with 1,990 records, and the full master dataset containing 34,576 records.

The English dataset corresponds to the manually curated English subset already introduced in Iteration~0.

\paragraph{2. Feature Engineering}
Binary labels were defined as Process Safety (1) and Non-Process Safety (0). Text features were constructed by concatenating the incident fields \texttt{TITLE} and \texttt{CASE\_DESCRIPTION}. Formally, the combined textual representation was defined as
\begin{equation}
\text{text\_features} = \text{TITLE} \oplus \text{``.''} \oplus \text{CASE\_DESCRIPTION},
\end{equation}
where \(\oplus\) denotes string concatenation. This representation combines the concise incident title with the more detailed narrative description in a single textual input for downstream embedding generation and classification.

\paragraph{3. Embedding Generation}
Contextual document representations were generated using the pre-trained \texttt{bert-base-cased} transformer model. The embedding process produced fixed-length vectors of dimensionality 768 by using the \texttt{[CLS]} token representation from the final hidden layer as the document-level embedding. To ensure efficient processing, embeddings were generated in batches with \texttt{batch\_size = 8} and progress tracking was applied throughout execution. The resulting embeddings were then serialized and stored as pickle files together with the corresponding metadata for reproducibility and later reuse.

\paragraph{4. Model Training}
For model training, a Linear Support Vector Machine (SVM) was employed as the classifier. To address class imbalance, the model was configured with \texttt{class\_weight='balanced'}. The data were divided using an 80/20 stratified train--test split with \texttt{random\_state=42} to ensure reproducibility while preserving the class distribution in both partitions. An independent classifier was then trained for each dataset in order to evaluate performance separately across the different country-specific and language-specific subsets.

\paragraph{5. Evaluation Metrics}
For each dataset, model performance was evaluated using a comprehensive set of classification metrics. These included the confusion matrix, reporting True Negatives (TN), False Positives (FP), False Negatives (FN), and True Positives (TP), as well as per-class precision, recall, and F1-score for both Process Safety and Non-Process Safety. In addition, macro-averaged precision, recall, and F1-score were computed, with macro-F1 serving as the primary evaluation metric because it gives equal importance to both classes despite class imbalance. To further assess classification behaviour, the False Positive Rate (FPR) and False Negative Rate (FNR) were also calculated, together with overall accuracy.

\subsubsection{Results and Outputs}
The outputs of Iteration~1 were organized systematically to support reproducibility, analysis, and later integration into the thesis. Generated embeddings were saved in \texttt{/Embeddings/\_iteration\_1/bert-base-cased/}, following the naming convention \texttt{\{dataset\}\_bert-base-cased\_embeddings.pkl}. In addition, confusion matrix heatmaps were exported as high-resolution PNG files at 300 DPI for visual inspection of classification performance. A CSV summary file containing all evaluation metrics across the processed datasets was also produced, together with formatted Markdown tables that could be incorporated into the LaTeX thesis document.

\subsubsection{Key Outcomes}
Six datasets were evaluated in Iteration~1, namely Germany, the UK, the manually curated English subset, Sweden, the Netherlands, and the full master dataset. The strongest result was obtained on the UK dataset, which achieved an accuracy of 0.7722 and a macro-F1 score of 0.7449. The manually curated English dataset performed similarly, with an accuracy of 0.7699 and a macro-F1 score of 0.7420. Performance on the non-English datasets was substantially weaker, particularly for Sweden and Germany, which achieved macro-F1 scores of 0.5825 and 0.5983, respectively. The Netherlands dataset reached a macro-F1 score of 0.6586, while the full master dataset achieved 0.6292. Across all datasets, False Negative Rates ranged from 0.2518 on the UK dataset to 0.3983 on the Netherlands dataset, indicating that recall for Process Safety incidents remained a substantial challenge in Iteration~1.

\subsection{Iteration 2: Multi-Model Comparative Evaluation}

\subsubsection{Objective}
Conduct a comprehensive comparative analysis of multiple pre-trained transformer models across diverse datasets to identify the optimal architecture for process safety incident classification.

\subsubsection{Implementation Steps}

\paragraph{1. Model Selection and Configuration}
Ten state-of-the-art transformer models were selected for evaluation in Iteration~2. The BERT family comprised \texttt{bert-base-cased}, with 110 million parameters and 768-dimensional embeddings, and \texttt{bert-large-cased}, with 340 million parameters and 1024-dimensional embeddings. The RoBERTa family included \texttt{roberta-base}, a 125 million-parameter model representing a robustly optimized BERT architecture, and \texttt{roberta-large}, a 355 million-parameter variant with enhanced pre-training. In addition, several specialized architectures were included: \texttt{distilbert-base-cased}, a distilled 66 million-parameter model designed for efficiency; \texttt{xlnet-base-cased}, a 110 million-parameter model based on permutation language modeling; \texttt{modernbert-base} and \texttt{modernbert-large}, with 139 million and 395 million parameters, respectively, representing modernized architectural variants; and \texttt{deberta-v3-base} and \texttt{deberta-v3-large}, with 184 million and 435 million parameters, respectively, incorporating disentangled attention mechanisms.

\paragraph{2. Dataset Selection}
Iteration~2 operated on the cleaned master dataset established in Iteration~0 and focused on four manually derived language- and country-specific subsets stored as JSON files, namely \texttt{english\_manual}, \texttt{german\_manual}, \texttt{swedish\_manual}, and \texttt{dutch\_manual}. On this basis, the verified experimental scope of Iteration~2 consisted of 10 transformer models evaluated across 4 datasets, yielding a total of \(10 \times 4 = 40\) planned model--dataset combinations.

\paragraph{3. Pipeline Architecture with Checkpointing}
A robust two-phase pipeline was implemented with comprehensive error handling and progress tracking. In Phase~1, embedding generation was performed in batches with \texttt{batch\_size=4} to ensure memory efficiency. GPU acceleration was used when available, with automatic fallback to CPU if necessary. Memory usage was controlled through periodic execution of \texttt{torch.cuda.empty\_cache()} after each batch, and all input sequences were truncated to a maximum length of 512 tokens. A checkpoint was saved after each successfully completed dataset--model combination. In Phase~2, classification and evaluation were carried out using \texttt{LinearSVC}, chosen for computational efficiency. Large training sets exceeding 60,000 samples were reduced through stratified subsampling, and a training timeout of 300 seconds per model was introduced to prevent stalled execution. To address class imbalance, the classifier was configured with \texttt{class\_weight='balanced'}, and the principal hyperparameters were set to \texttt{max\_iter=1000}, \texttt{dual=False}, and \texttt{random\_state=42}.

\paragraph{4. Checkpoint System Implementation}
A JSON-based checkpoint system enabled resumable execution throughout the experimental pipeline. The checkpoint file, \texttt{checkpoint.json}, tracked completed \((\text{model}, \text{dataset})\) tuples and made it possible to skip already processed combinations when the pipeline was restarted. Additional utility functions, namely \texttt{view\_progress()}, \texttt{reset\_checkpoint()}, and \texttt{complete\_rerun()}, supported monitoring and control of the execution process. This design ensured that progress was preserved across kernel restarts and unexpected failures.

\paragraph{5. Memory Management Strategy}
Aggressive memory management was applied in order to prevent out-of-memory errors during large-scale experimentation. This included batch-level GPU memory clearing, periodic garbage collection every 50 batches, unloading models between experiments, and the explicit deletion of large variables using \texttt{del} followed by \texttt{gc.collect()}. Where necessary, dataset subsampling was also employed during SVM training to reduce computational load while preserving statistical validity.

\paragraph{6. Comprehensive Evaluation Metrics}
For each model--dataset combination, a comprehensive set of evaluation metrics was computed. Overall performance was assessed using accuracy, macro-averaged F1-score, and weighted F1-score. Accuracy was defined as \(\frac{TP + TN}{TP + TN + FP + FN}\), macro-F1 as \(\frac{F1_{\text{PS}} + F1_{\text{NPS}}}{2}\), and weighted F1 as \(\frac{\sum_{c} F1_c \times \text{support}_c}{\sum_{c} \text{support}_c}\). In addition, per-class precision, recall, F1-score, and support were calculated separately for the Process Safety and Non-Process Safety classes. To provide further experimental context, additional metadata were recorded, including training time in seconds, a subsampling flag indicating whether the training set had been reduced, and the individual confusion-matrix elements TN, FP, FN, and TP.

\paragraph{7. Visualization and Reporting}
To support interpretation and documentation, confusion matrix heatmaps were generated for every model--dataset combination and saved as high-resolution PNG files at 300 DPI using the naming convention \texttt{cm\_\{dataset\}\_\{model\}.png}. In addition, all computed metrics were exported into a comprehensive CSV file, formatted classification reports were generated in a style analogous to \texttt{sklearn}, and aggregate statistics were produced both by model and by dataset.

\subsubsection{Results and Outputs}

\paragraph{Storage Structure}
The outputs of Iteration~2 were stored in a structured directory hierarchy to ensure reproducibility and ease of access. Generated embeddings were saved in \texttt{/Embeddings/\_iteration\_2/\{model\_name\}/}, while all result files were written to \texttt{/Results/\_iteration\_2/}. The checkpoint state was stored in \texttt{/Results/\_iteration\_2/checkpoint\_iteration\_2.json}, and the consolidated summary of all model results was exported to \texttt{/Results/\_iteration\_2/all\_models\_svm\_detailed\_results.csv}.

\paragraph{Aggregate Analysis}
Several aggregate analytical views were generated in order to interpret the experimental outcomes systematically. At the model level, mean values and standard deviations of accuracy and F1-scores across all datasets were computed to identify consistently high-performing architectures and assess model stability. At the dataset level, mean values and standard deviations across all models were used to estimate dataset difficulty and reveal language-specific performance patterns, particularly the contrast between English and non-English subsets. In addition, best-performing combinations were identified according to overall accuracy, Process Safety F1-score, and efficiency measured as the trade-off between predictive performance and training time.

\subsubsection{Key Innovations}

\paragraph{1. Scalability}
The pipeline successfully processed 24 completed model--dataset combinations, corresponding to 6 completed models evaluated across 4 datasets. Four models, namely \texttt{modernbert-base}, \texttt{modernbert-large}, \texttt{deberta-v3-base}, and \texttt{deberta-v3-large}, encountered compatibility issues and were therefore excluded from the final completed set. Automatic checkpoint-based recovery ensured that completed combinations did not need to be recomputed after interruptions.

\paragraph{2. Reproducibility}
Experimental reproducibility was ensured through fixed random seeds, particularly \texttt{random\_state=42}, comprehensive metadata tracking, and consistent version-controlled naming of models and outputs. Together, these measures allowed the experimental pipeline to be rerun under the same conditions and to produce consistent results.

\paragraph{3. Efficiency}
Efficiency was improved through several design decisions. The checkpoint system eliminated redundant computation by allowing interrupted runs to resume from the last completed step. Strategic subsampling reduced the computational burden of large datasets while maintaining statistical validity, and the combination of batch processing with explicit memory management enabled experimentation on resource-constrained hardware.

\paragraph{4. Robustness}
The pipeline was also designed for robustness. Timeout protection prevented infinite training loops, exception handling with detailed logging captured and documented model-specific failures, and graceful degradation ensured that failed combinations were skipped while the remaining experiments continued uninterrupted.

\subsubsection{Limitations and Observations}
Despite these strengths, several limitations were observed in Iteration~2. Large models such as \texttt{deberta-v3-large} and \texttt{modernbert-large} required substantial processing time, often exceeding three hours per dataset. The necessary reduction of training sets through subsampling may have slightly reduced model performance in some cases. Four models, specifically \texttt{modernbert-base}, \texttt{modernbert-large}, \texttt{deberta-v3-base}, and \texttt{deberta-v3-large}, encountered compatibility issues with certain datasets and could not be fully evaluated. In addition, the use of \texttt{LinearSVC} instead of a kernel-based SVM may have limited the complexity of the decision boundaries that could be learned from the embeddings. Finally, the application of English-pretrained models to German, Swedish, and Dutch datasets introduced a language mismatch that likely constrained performance on the non-English subsets. These limitations motivated the design of Iteration~3, which focuses on multilingual transformer architectures to address the observed weaknesses.

\subsection{Iteration 3: Multilingual Transformer Evaluation}

\subsubsection{Objective}
Evaluate specialized multilingual transformer models to address the language mismatch limitation identified in Iteration 2, where English-trained models were applied to non-English datasets (German, Swedish, Dutch).

\subsubsection{Rationale}
Previous iterations demonstrated that monolingual English models (BERT-base-cased, RoBERTa) achieved strong performance on English datasets but exhibited suboptimal results on non-English data. Multilingual models, pre-trained on diverse language corpora, provide native cross-lingual capabilities that may better capture semantic relationships in multilingual safety incident reports.

\subsubsection{Implementation Steps}

\paragraph{1. Multilingual Model Selection}
Five state-of-the-art multilingual transformer models were selected for Iteration~3 on the basis of language coverage and architectural design. Two multilingual BERT variants were included, namely \texttt{bert-base-multilingual-uncased} and \texttt{bert-base-multilingual-cased}, each comprising 110 million parameters and trained on 104 languages, with the former using case-insensitive tokenization and the latter preserving case information. In addition, three multilingual E5 embedding models were evaluated: \texttt{intfloat/multilingual-e5-base}, which is optimized for semantic similarity across more than 100 languages; \texttt{intfloat/multilingual-e5-large}, a scaled version with 560 million parameters and enhanced embedding quality; and \texttt{intfloat/multilingual-e5-large-instruct}, an instruction-tuned variant designed to improve task-specific performance. The E5 family represents a class of multilingual text embedding models optimized for semantic search and similarity tasks through explicit multilingual training objectives.

\paragraph{2. Dataset Configuration}
Iteration~3 followed the dataset design introduced in Iteration~2 and operated on four manually derived datasets, namely \texttt{english\_manual}, \texttt{german\_manual}, \texttt{swedish\_manual}, and \texttt{dutch\_manual}. These four datasets were used to evaluate the five multilingual transformer models, yielding a verified experimental scope of \(5 \times 4 = 20\) model--dataset combinations. Several additional datasets, including Hungary, International, Russia, and UAE, were excluded because they did not contain sufficient samples for reliable stratified splitting.

\paragraph{3. Pipeline Architecture}
The implementation followed a two-phase architecture broadly consistent with Iteration~2, but with optimizations tailored to multilingual tokenization and embedding generation. In Phase~1, tokenizer initialization used \texttt{trust\_remote\_code=True} where necessary in order to support model-specific implementations. For models lacking a default padding token, padding was configured automatically according to
\begin{equation}
\text{tokenizer.pad\_token} = \text{tokenizer.eos\_token}.
\end{equation}
Embeddings were generated with a batch size of 8, increased from Iteration~2 due to the smaller memory demands of the selected models, and all inputs were truncated to a maximum sequence length of 512 tokens. Device-aware processing enabled automatic GPU utilization when available, while periodic cache clearing every 50 batches reduced the risk of memory overflow. In Phase~2, classification and evaluation were carried out using \texttt{LinearSVC} with \texttt{dual=False}, \texttt{class\_weight='balanced'}, \texttt{max\_iter=1000}, and \texttt{random\_state=42}. Where necessary, training sets were subsampled to a maximum of 60,000 samples through stratified resampling in order to maintain computational feasibility.

\paragraph{4. Feature Engineering}
Text features were constructed in the same way as in previous iterations by concatenating the \texttt{TITLE} and \texttt{CASE\_DESCRIPTION} fields. Formally, the combined textual representation was defined as
\begin{equation}
\text{text\_features} = \text{TITLE} \oplus \text{``.''} \oplus \text{CASE\_DESCRIPTION},
\end{equation}
where \(\oplus\) denotes string concatenation. This strategy ensured that both the concise summary information in the title and the richer contextual information in the description were jointly provided to the multilingual encoders.

\paragraph{5. Checkpoint and Recovery System}
A robust checkpoint system enabled resumable execution throughout Iteration~3. Completed \((\text{model}, \text{dataset})\) combinations were serialized to JSON using \texttt{save\_checkpoint(processed\_items)}, while \texttt{load\_checkpoint()} restored previous progress from the checkpoint file. Additional helper functions, including \texttt{view\_progress()} and \texttt{reset\_checkpoint()}, supported monitoring and control of the execution process. A checkpoint was saved after each successful embedding-generation step, already completed combinations were skipped automatically on restart, and progress was tracked in the form \texttt{processed/total}. Persistent checkpoint storage was maintained in \texttt{/\_Iterations/\_iteration\_3/checkpoint.json}.

\paragraph{6. Embedding Generation Function}
A unified embedding function was implemented to support both standard multilingual BERT models and specialized multilingual E5 models. This function extracted the \texttt{[CLS]} token representation while incorporating batch-level memory management, as shown in Listing~\ref{lst:unified-embedding} in Appendix~\ref{sec:code-listings}. Fallback logic was included to account for differences in output structure across architectures, thereby ensuring compatibility across the multilingual BERT and E5 model families.

\paragraph{7. Comprehensive Evaluation Metrics}
The evaluation protocol in Iteration~3 mirrored that of Iteration~2. Overall performance was assessed using accuracy, weighted F1-score, and macro F1-score, with macro-F1 retained as the primary evaluation metric. Accuracy was defined as \(\frac{TP + TN}{TP + TN + FP + FN}\), weighted F1-score as \(\sum_{c \in \{\text{PS, NPS}\}} F1_c \times \frac{\text{support}_c}{\text{total support}}\), and macro F1-score as \(\frac{F1_{\text{PS}} + F1_{\text{NPS}}}{2}\). In addition, per-class precision, recall, F1-score, and support were computed for both classes. To support a more complete comparison of experimental behaviour, training time in seconds, a subsampling indicator, and the confusion-matrix components TN, FP, FN, and TP were also recorded.

\paragraph{8. Results Storage and Visualization}
The outputs of Iteration~3 were stored in a structured directory hierarchy. Embeddings were saved in \texttt{/Embeddings/\_iteration\_3/\{model\_key\}/}, while results and visual outputs were stored in \texttt{/Results/\_iteration\_3/}. Confusion matrices for each model--dataset combination were exported in that directory, with embeddings named according to the pattern \texttt{\{dataset\}\_\{model\_key\}\_embeddings.pkl} and confusion matrices following the convention \texttt{cm\_\{dataset\}\_\{model\_key\}.png}.

\subsubsection{Results Analysis Framework}

\paragraph{Multi-Dimensional Comparative Analysis}
The results were analyzed along three complementary dimensions. First, model-level comparison was performed by computing aggregate statistics, including means and standard deviations of accuracy, F1-scores, and per-class metrics across all datasets, thereby enabling the identification of consistently strong architectures and the assessment of cross-dataset stability. Second, dataset-level comparison examined average performance across all models in order to reveal language-specific patterns, compare English, German, Swedish, and Dutch subsets, and rank datasets by difficulty on the basis of average F1-scores. Third, best-performing combinations were identified according to overall accuracy, Process Safety F1-score, and computational efficiency measured as performance relative to training time.

\subsubsection{Key Innovations}

\paragraph{1. Language-Aware Architecture}
A central innovation of Iteration~3 was the shift to multilingual architectures that directly addressed the language-mismatch problem identified in Iteration~2. These models provided native support for German, Swedish, and Dutch incident reports without requiring translation or language-specific fine-tuning.

\paragraph{2. E5 Embedding Models}
The inclusion of multilingual E5 models introduced contrastive learning-based embeddings optimized for semantic similarity, offering the potential to capture incident similarity patterns more effectively than conventional masked language models.

\paragraph{3. Instruction Tuning}
The \texttt{multilingual-e5-large-instruct} model added an instruction-tuned component that potentially enables improved task-specific adaptation through natural-language instructions, even though such instruction-following capabilities were not explicitly exploited in this iteration.

\paragraph{4. Unified Embedding Interface}
The generalized embedding function abstracted model-specific implementation details and enabled a uniform processing interface across multilingual BERT and E5 architectures, thereby simplifying experimentation and facilitating future model extensions.

\paragraph{5. Enhanced Checkpoint Granularity}
The checkpoint system tracked individual \((\text{model}, \text{dataset})\) combinations rather than waiting for full model completion, which reduced wasted computation in the presence of partial failures and improved overall execution resilience.

\subsubsection{Expected Outcomes}
Based on the capabilities of multilingual transformer models, several outcomes were anticipated. First, multilingual BERT variants were expected to outperform monolingual English models on the German, Swedish, and Dutch datasets by approximately 5--10 percentage points in macro-F1. Second, performance on English datasets was expected to remain comparable to, or only slightly below, the Iteration~2 monolingual baselines because multilingual pretraining distributes representational capacity across many languages. Third, the larger E5 models were expected to demonstrate stronger semantic understanding, particularly in capturing incident similarity relationships. Finally, it was anticipated that multilingual architectures would be especially beneficial for mixed-language or cross-lingual settings by enabling more effective transfer across reporting languages.

\subsubsection{Limitations and Considerations}
Several limitations were also recognised. Although multilingual models have similar parameter counts to monolingual BERT in some cases, they may require longer inference times due to substantially larger vocabularies. In addition, multilingual pretraining corpora are typically imbalanced across languages, which may favour high-resource languages such as English and German over lower-resource languages such as Swedish and Dutch. Another limitation arises from domain shift, since these models are pre-trained mainly on Wikipedia and web text rather than technical process-safety incident reports. Multilingual tokenizers may also fragment equipment names, chemical terminology, and other technical expressions more aggressively than domain-adapted tokenizers. Finally, the E5 family is optimized primarily for semantic similarity rather than classification, which means that combining these embeddings with SVM classifiers may not fully exploit their representational capacity.

\subsubsection{Comparison with Previous Iterations}

\begin{table}[htbp]
\centering
\small
\begin{tabular}{@{}llll@{}}
\toprule
\textbf{Aspect} & \textbf{Iteration 1} & \textbf{Iteration 2} & \textbf{Iteration 3} \\
\midrule
Models & 1 (BERT-base-cased) & 10 planned / 6 completed & 5 (multilingual) \\
Datasets & 6 & 4 & 4 \\
Combinations & 6 & 24 completed & 20 \\
Language Support & English only & English only & 104+ languages \\
Batch Size & 8 & 4 & 8 \\
Primary Focus & Baseline & Architecture comparison & Language coverage \\
Key Innovation & Stratification & Checkpointing & Multilingual embeddings \\
\bottomrule
\end{tabular}
\caption{Comparison of Iterations 1--3.}
\label{tab:iteration_comparison}
\end{table}

These observations informed the design of Iteration~4, which explores more advanced machine learning techniques and stronger multilingual embedding pipelines in order to improve classification accuracy further and reduce the remaining effects of domain shift.


\subsection{Iteration 4: Advanced Machine Learning Pipeline with Enhanced Techniques}

\subsubsection{Objective}
Implement a comprehensive machine learning pipeline incorporating advanced techniques to maximize classification performance through improved embedding strategies, class imbalance handling, multiple classifier architectures, hyperparameter optimization, and ensemble methods.

\subsubsection{Rationale}
Previous iterations identified several optimization opportunities: (1) [CLS] token embeddings may not fully capture document semantics, (2) severe class imbalance (Process Safety incidents constitute 10-20\% of datasets) degrades minority class recall, (3) LinearSVC may not represent the optimal classifier architecture, and (4) single-model predictions lack robustness. Iteration 4 addresses these limitations through a multi-faceted enhancement strategy.

\subsubsection{Implementation Steps}

\paragraph{1. Model Selection and Configuration}
A broader set of multilingual embedding models was evaluated in Iteration~4, and the current repository state shows multiple embedding families combined with several downstream classifiers. The embedding models explicitly documented in the current implementation include \texttt{bert-base-multilingual-uncased}, a 110M-parameter model trained on 104 languages; \texttt{xlm-roberta-base}, a 270M-parameter cross-lingual transformer trained on 100 languages; and \texttt{sentence-transformers/paraphrase-multilingual-mpnet-base-v2}, a 278M-parameter Sentence-BERT variant optimized for semantic similarity in multilingual settings. Overall, Iteration~4 evaluates several multilingual embedding models across the manually derived language datasets using multiple downstream classifiers. Consequently, the final experiment count should be reported directly from the saved Iteration~4 result artefacts rather than retained from an earlier draft.

\paragraph{2. Enhanced Configuration Framework}
A comprehensive configuration dictionary controlled all major pipeline parameters, including pooling strategy, class balancing, feature scaling, hyperparameter tuning, ensemble activation, and the set of downstream classifiers (see Listing~\ref{lst:iter4-config} in Appendix~\ref{sec:code-listings}). This modular structure was deliberately designed to support systematic ablation studies and to quantify the contribution of each enhancement to the final classification performance.

\paragraph{3. Advanced Text Preprocessing}
Advanced preprocessing was introduced to reduce domain-specific noise in the incident reports and to improve the quality of downstream text representations (see Listing~\ref{lst:text-preprocessing} in Appendix~\ref{sec:code-listings}). In particular, the preprocessing pipeline applied lowercase normalization, removed URLs that commonly appeared in copied or forwarded report text, filtered email addresses for privacy preservation, removed special characters while preserving sentence-boundary periods, and normalized irregular whitespace. Together, these steps produced cleaner and more standardized textual inputs for multilingual embedding generation.

\paragraph{4. Enhanced Embedding Generation with Pooling Strategies}
Iteration~4 implemented three pooling strategies in order to aggregate token-level hidden states into a fixed-length document representation. The baseline approach used CLS pooling, defined as
\begin{equation}
\mathbf{h}_{\text{CLS}} = \text{LastHiddenState}[:, 0, :]
\end{equation}
which extracts the contextual representation of the first token and corresponds to the standard BERT-style classification setup. In addition, mean pooling was introduced as the primary strategy:
\begin{equation}
\mathbf{h}_{\text{mean}} = \frac{\sum_{i=1}^{n} m_i \cdot \mathbf{h}_i}{\sum_{i=1}^{n} m_i},
\end{equation}
where $m_i$ denotes the attention mask for token $i$ and $\mathbf{h}_i$ denotes the corresponding hidden-state vector. Mean pooling averages all valid token representations and therefore provides a more comprehensive document-level representation than the CLS token alone. As a third alternative, max pooling was implemented:
\begin{equation}
\mathbf{h}_{\text{max}} = \max_{i=1}^{n} \mathbf{h}_i
\end{equation}
which captures the maximum activation across token positions and emphasizes salient features. The full implementation supporting CLS, mean, and max pooling is provided in Listing~\ref{lst:enhanced-embedding} in Appendix~\ref{sec:code-listings}.

\paragraph{5. SMOTE-Based Class Balancing}
To address the severe class imbalance between Process Safety and Non-Process Safety incidents, Iteration~4 incorporated the Synthetic Minority Over-sampling Technique (SMOTE). In this setting, the majority class typically represents 80--90\% of the data, whereas the minority Process Safety class accounts for only 10--20\%, which can bias classifiers toward the majority class and substantially reduce minority-class recall. SMOTE mitigates this problem by creating synthetic minority examples in feature space. For each minority-class sample $\mathbf{x}_i$, one of its $k$ nearest neighbors $\mathbf{x}_{\text{nn}}$ is selected, and a new synthetic example is generated by interpolation:
\begin{equation}
\mathbf{x}_{\text{new}} = \mathbf{x}_i + \lambda \cdot (\mathbf{x}_{\text{nn}} - \mathbf{x}_i),
\end{equation}
where $\lambda \sim \text{Uniform}(0, 1)$. This procedure is repeated until the minority class reaches the desired sampling ratio. In the Iteration~4 implementation, \texttt{sampling\_strategy=0.7} was used so that the minority class was oversampled to 70\% of the majority-class size, \texttt{k\_neighbors=5} was selected whenever feasible, and \texttt{random\_state=42} ensured reproducibility. The corresponding implementation is shown in Listing~\ref{lst:smote} in Appendix~\ref{sec:code-listings}.

\paragraph{6. Feature Scaling}
Feature scaling was applied using \texttt{StandardScaler} in order to normalize embedding features to zero mean and unit variance:
\begin{equation}
\mathbf{x}_{\text{scaled}} = \frac{\mathbf{x} - \mu}{\sigma},
\end{equation}
where $\mu$ denotes the feature mean and $\sigma$ denotes the standard deviation estimated from the training data. This transformation is especially important for SVM and other distance-sensitive classifiers, since unscaled embedding dimensions can otherwise distort the geometry of the feature space. In the pipeline, the scaler was fitted on the balanced training data and subsequently applied to the test data using the same parameters, as shown in Listing~\ref{lst:feature-scaling} in Appendix~\ref{sec:code-listings}.

\paragraph{7. Multiple Classifier Architectures}
Iteration~4 extended the earlier single-classifier setup by training five different classifier types on each embedding representation. The first classifier was a linear Support Vector Machine implemented via \texttt{LinearSVC} with \texttt{dual=False} for computational efficiency on large datasets. Its hyperparameters were optimized using \texttt{GridSearchCV} with three-fold cross-validation, where the regularization parameter was searched over \texttt{C} $\in \{0.1, 1, 10, 100\}$ and the class-weight configuration was selected from \texttt{balanced}, $\{0:1, 1:2\}$, and $\{0:1, 1:3\}$, using weighted F1 as the optimization criterion. Second, a Random Forest classifier with \texttt{n\_estimators=200}, \texttt{max\_depth=20}, and \texttt{class\_weight='balanced'} was used to provide a robust tree-based baseline with parallel execution through \texttt{n\_jobs=-1}. Third, XGBoost was included with \texttt{n\_estimators=200}, \texttt{max\_depth=6}, and \texttt{learning\_rate=0.1}, while imbalance was handled through \texttt{scale\_pos\_weight} computed as $\frac{N_{\text{negative}}}{N_{\text{positive}}}$. Fourth, LightGBM was trained with \texttt{n\_estimators=200}, \texttt{max\_depth=6}, \texttt{learning\_rate=0.1}, and built-in balanced class weights. Finally, a hard-voting ensemble combined the predictions of the top three individual classifiers and determined the final label by majority vote:
\begin{equation}
\hat{y} = \text{mode}(\hat{y}_{\text{SVM}}, \hat{y}_{\text{XGBoost}}, \hat{y}_{\text{RF}}).
\end{equation}
The full multi-classifier training function, including SVM with \texttt{GridSearchCV}, Random Forest, XGBoost, LightGBM, and the Voting Ensemble, is provided in Listing~\ref{lst:multi-classifier} in Appendix~\ref{sec:code-listings}.

\paragraph{8. Comprehensive Evaluation Metrics}
The evaluation framework in Iteration~4 extended the metric set used in earlier iterations. At the classification level, the main metrics included accuracy, precision, recall, and F1-score, reported both per class and in aggregated form. Macro F1-score served as the primary metric because it gives equal weight to both classes regardless of class frequency, whereas weighted F1-score provided an additional support-weighted view of performance. When probabilistic outputs were available, the evaluation also incorporated ROC AUC and Average Precision. ROC AUC was defined as
\begin{equation}
\text{AUC} = \int_{0}^{1} \text{TPR}(\text{FPR}) \, d(\text{FPR}),
\end{equation}
while Average Precision was computed as
\begin{equation}
\text{AP} = \sum_{k=1}^{n} P(k) \cdot \Delta R(k).
\end{equation}
In addition to predictive metrics, the pipeline also recorded operational indicators such as total training time, best hyperparameter settings from \texttt{GridSearchCV}, and the associated cross-validation score.

\paragraph{9. Enhanced Visualization Suite}
To support comprehensive model inspection, Iteration~4 generated a four-panel evaluation visualization for each classifier. The first panel displayed the confusion matrix as a heatmap with annotated counts and a blue color scale, using abbreviated labels for the two classes. The second panel plotted the ROC curve, showing the trade-off between true positive rate and false positive rate, including the AUC value and a diagonal reference line representing random classification. The third panel showed the precision--recall curve together with the Average Precision score, which is particularly informative under class imbalance because it emphasizes minority-class performance. The fourth panel summarized key metrics in a horizontal bar chart, including accuracy, weighted F1, Process Safety F1, Process Safety precision, and Process Safety recall, with all values normalized to the interval $[0,1]$. The corresponding implementation is provided in Listing~\ref{lst:visualization} in Appendix~\ref{sec:code-listings}.

\paragraph{10. Results Storage and Management}
All embeddings generated in Iteration~4 were stored in the directory \texttt{/Embeddings/\_iteration\_4/\{model\}/}, while classifier outputs and visualizations were written to \texttt{/Results/\_iteration\_4/}. Confusion-matrix visualizations for each classifier were saved using the naming pattern \texttt{cm\_\{dataset\}\_\{model\}\_\{classifier\}.png}. In addition, intermediate and processed data were serialized in pickle format. These pickle files contained the original train--test splits (\texttt{X\_train}, \texttt{X\_test}, \texttt{y\_train}, \texttt{y\_test}), the balanced training data after SMOTE (\texttt{X\_train\_balanced}, \texttt{y\_train\_balanced}), the final processed feature matrices (\texttt{X\_train\_subset}, \texttt{X\_test\_scaled}), the fitted \texttt{StandardScaler}, and a metadata object recording the effective configuration used for that run.

\subsubsection{Results Analysis Framework}

\paragraph{Multi-Dimensional Performance Analysis}
The results of Iteration~4 were analyzed across four complementary dimensions. First, model comparison focused on aggregate statistics across all datasets and classifiers in order to identify the strongest embedding models and to assess performance stability through variance analysis. Second, classifier comparison examined average performance across all datasets and embedding representations, thereby identifying the most effective downstream classifier architecture and quantifying the value added by the ensemble relative to individual classifiers. Third, dataset comparison evaluated language-specific performance patterns and ranked datasets according to difficulty, based on average scores across all model--classifier combinations. Fourth, technique ablation was designed to isolate the effects of key enhancements, including SMOTE, mean pooling, hyperparameter tuning, and ensemble voting.

\paragraph{Comparative Visualization}
A second layer of comparative visualization was used to summarize aggregated performance across experiments. The first panel displayed a horizontal bar chart of average Process Safety F1-score by classifier, sorted in descending order and annotated with exact values. The second panel showed average Process Safety F1-score by embedding model in order to identify the strongest transformer architecture. The third panel plotted accuracy against F1-score for each classifier--dataset combination, revealing performance trade-offs between overall correctness and class-balanced prediction quality. The fourth panel plotted Process Safety precision against Process Safety recall, color-coded by classifier type, in order to identify the configurations that achieved the most favorable precision--recall balance.

\subsubsection{Key Innovations}

\paragraph{1. Mean Pooling Strategy}
A central innovation of Iteration~4 was the replacement of exclusive CLS-token pooling with attention-mask-weighted mean pooling. Because incident descriptions can be relatively long and semantically diffuse, averaging across all valid token representations yields a richer and more stable document representation than relying solely on the first token.

\paragraph{2. SMOTE-Based Balancing}
Another major enhancement was the integration of SMOTE-based balancing. Unlike undersampling, which discards majority-class information, SMOTE increases minority-class representation synthetically and therefore improves recall for Process Safety incidents while preserving more of the original dataset structure.

\paragraph{3. Multi-Classifier Architecture}
Iteration~4 also expanded the experimental design from a single-classifier setting to a multi-classifier architecture including linear SVM, Random Forest, XGBoost, LightGBM, and a Voting Ensemble. This made it possible to compare linear margin-based, bagging-based, and gradient-boosting approaches within a common embedding pipeline.

\paragraph{4. Ensemble Robustness}
The inclusion of a hard-voting ensemble introduced an additional robustness mechanism by combining the strengths of diverse classifiers. This strategy reduces variance and can improve generalization, especially when individual classifiers make complementary errors.

\paragraph{5. Hyperparameter Optimization}
The use of \texttt{GridSearchCV} added a principled hyperparameter-optimization layer to the pipeline. Rather than relying on default settings, Iteration~4 systematically tuned key classifier parameters such as SVM regularization and class weights, enabling more competitive and reproducible performance.

\paragraph{6. Probabilistic Evaluation}
Finally, Iteration~4 extended the evaluation beyond thresholded class labels by incorporating probabilistic metrics such as ROC AUC and Average Precision. These measures provide a more nuanced understanding of classifier behavior and are particularly useful when considering deployment scenarios with adjustable decision thresholds.

\subsubsection{Expected Outcomes}
Based on the implemented enhancements, Iteration~4 was expected to improve Process Safety recall through SMOTE balancing, to outperform CLS pooling through mean pooling on longer incident descriptions, and to demonstrate the superiority of gradient-boosting classifiers over linear SVM in terms of Process Safety F1-score. In addition, the hard-voting ensemble was expected to provide a modest improvement over the best single classifier, while the overall macro F1-score was expected to approach or exceed 0.90 on at least one dataset and Process Safety F1-score was expected to surpass 0.85 across all datasets.

\subsubsection{Limitations and Considerations}
Despite these enhancements, several methodological limitations remained. Synthetic samples generated by SMOTE in high-dimensional embedding space may not always preserve realistic minority-class structure and can therefore introduce noise. Training five classifiers for each embedding model substantially increases computational cost compared with earlier iterations. The hyperparameter search space explored by \texttt{GridSearchCV} remained limited by practical constraints, which means that further gains might still be achievable with more advanced optimization methods such as Bayesian search. Ensemble predictions, while often stronger, are less interpretable than those of individual classifiers. Moreover, gradient-boosting models remain sensitive to parameter settings such as learning rate, depth, and number of estimators, especially on smaller datasets. Mean pooling also assumes that all token representations contribute approximately equally to document semantics, which may dilute highly informative process-safety keywords. Finally, SMOTE can become memory-intensive on large training sets, potentially limiting scalability.

\subsubsection{Comparison with Previous Iterations}

\begin{table}[htbp]
\centering
\small
\begin{tabular}{@{}lllll@{}}
\toprule
\textbf{Aspect} & \textbf{Iter 1} & \textbf{Iter 2} & \textbf{Iter 3} & \textbf{Iter 4} \\
\midrule
Models & 1 & 10 & 5 & 3 \\
Datasets & 6 & 6 & 6 & 6 \\
Classifiers & 1 (SVM) & 1 (SVM) & 1 (SVM) & 5 + Ensemble \\
Combinations & 6 & 60 & 30 & 90 \\
Pooling & CLS & CLS & CLS & Mean \\
Class Balance & Weighted & Weighted & Weighted & SMOTE \\
Scaling & No & No & No & Yes \\
Tuning & No & No & No & GridSearchCV \\
Ensemble & No & No & No & Yes \\
ROC/PR Curves & No & No & No & Yes \\
Key Innovation & Baseline & Multi-model & Multilingual & Advanced ML \\
\bottomrule
\end{tabular}
\caption{Comparison of Iterations 1-4}
\label{tab:iteration_comparison_4}
\end{table}

\subsubsection{Ablation Study Design}
To quantify the individual contribution of each enhancement, a structured ablation design was defined. This included comparisons between CLS, mean, and max pooling on the same model--dataset pairs; experiments with and without SMOTE on identical embeddings; tests of SVM performance with and without \texttt{StandardScaler}; comparisons between default and \texttt{GridSearchCV}-optimized hyperparameters; and an analysis of the improvement obtained by the Voting Ensemble over the best individual classifier. By holding all other settings constant while changing only one factor at a time, these ablation experiments provide a basis for causal interpretation of the observed performance differences.

\subsubsection{Future Directions}
Iteration~4 established a comprehensive baseline for advanced machine-learning techniques, but several extensions remain promising. These include domain-specific fine-tuning of multilingual transformers on process-safety corpora, attention-based visualization for identifying the textual features driving predictions, active-learning strategies for prioritizing uncertain samples near the decision boundary, multi-task learning for combining classification with severity regression, the integration of explainability methods such as LIME and SHAP, and model-distillation techniques for translating ensemble performance into a lighter production-ready deployment model. Collectively, these directions provide a pathway from a research-focused experimental framework toward a practically deployable incident-classification system.

\subsection{Iteration 5: Aggressive Process Safety Targeting}

\subsubsection{Objective}
Achieve significantly higher Process Safety recall through aggressive class balancing strategies, targeting 95\%+ recall for the minority class while maintaining overall classification performance.

\subsubsection{Rationale}
Previous iterations demonstrated that despite improvements in overall accuracy and macro F1-scores, the Process Safety class (minority) consistently exhibited lower recall compared to the Non-Process Safety class. In safety-critical applications, missing a true Process Safety incident (false negative) carries substantially higher risk than incorrectly flagging a non-critical incident. Iteration 5 prioritizes minority class recall through aggressive resampling and cost-sensitive learning.

\subsubsection{Implementation Steps}

\paragraph{1. Enhanced Configuration}
Key configuration changes relative to Iteration~4 are presented in Listing~\ref{lst:iter5-config} in Appendix~\ref{sec:code-listings}. The most important modifications include full 1:1 SMOTE balancing through \texttt{smote\_strategy='auto'}, the integration of SMOTETomek resampling, threshold optimisation aimed at achieving 95\% recall for the Process Safety class, and the use of a 10$\times$ class-weight multiplier for Process Safety incidents. Collectively, these changes reflect a deliberate shift toward aggressive minority-class targeting in order to minimise false negatives in safety-critical classification.

\paragraph{2. Full SMOTE Oversampling}
In contrast to Iteration~4, where SMOTE increased the minority class only to 70\% of the majority class, Iteration~5 applies full 1:1 balancing. Formally, the number of synthetic minority-class samples is defined as
\begin{equation}
N_{\text{minority}}^{\text{synthetic}} = N_{\text{majority}} - N_{\text{minority}}^{\text{original}}.
\end{equation}
This yields a perfectly balanced training set and is intended to remove the systematic learning bias toward the majority Non-Process Safety class.

\paragraph{3. Strategic Undersampling}
In addition to full SMOTE oversampling, Iteration~5 applies strategic undersampling in order to further control the post-resampling class distribution. The final majority-class size is defined as
\begin{equation}
N_{\text{majority}}^{\text{final}} = 0.8 \times N_{\text{minority}}^{\text{after SMOTE}}.
\end{equation}
Thus, after oversampling the minority class, the majority class is reduced to 80\% of the minority-class size. The combined SMOTETomek strategy further removes overlapping observations near the decision boundary, thereby attempting to clean ambiguous border regions before classifier training. The corresponding implementation is provided in Listing~\ref{lst:smotetomek} in Appendix~\ref{sec:code-listings}.

\paragraph{4. Threshold Optimization}
Rather than using the default decision threshold of 0.5, Iteration~5 performs threshold optimisation to explicitly target high Process Safety recall. This is particularly important in the present application, where false negatives are more critical than false positives. The threshold-optimisation function used for this purpose is documented in Listing~\ref{lst:threshold-opt} in Appendix~\ref{sec:code-listings}.

\paragraph{5. Aggressive Class Weights}
Beyond resampling, class weighting was strengthened substantially in order to penalise missed Process Safety incidents more heavily. The Process Safety class weight is defined as
\begin{equation}
w_{\text{PS}} = 10 \times \frac{N_{\text{total}}}{2 \times N_{\text{PS}}}.
\end{equation}
The additional factor of 10 imposes a much stronger penalty on false negatives during optimisation and therefore biases training toward higher sensitivity for the minority class.

\paragraph{6. Model Selection}
Iteration~5 evaluated nine multilingual transformer models: \texttt{bert-base-multilingual-uncased}, \texttt{xlm-roberta-base}, \texttt{paraphrase-multilingual-mpnet-base-v2}, \texttt{paraphrase-multilingual-MiniLM-L12-v2}, \texttt{snowflake-arctic-embed-l-v2.0}, \texttt{labse-gguf}, \texttt{upskyy/bge-m3-korean}, \texttt{intfloat/multilingual-e5-large-instruct}, and \texttt{efederici/e5-base-multilingual-4096}. These models were selected to represent a broad range of multilingual embedding architectures under the aggressive balancing regime introduced in this iteration.

\subsubsection{Expected Outcomes}
The primary expectation of Iteration~5 is that the Process Safety recall will increase to at least 95\%. In addition, this iteration is intended to quantify the precision--recall trade-off created by aggressive thresholding and balancing, to assess the effect of these interventions on overall macro F1-score, and to examine how strongly the lower decision threshold increases the number of false positives. Accordingly, the focus of this iteration is not only on achieving very high minority-class recall, but also on understanding the operational costs of such an aggressive classification strategy.

\subsubsection{Limitations}
Despite its potential advantages, Iteration~5 is subject to several important limitations. Full SMOTE in high-dimensional embedding space may generate synthetic observations that are not semantically realistic, thereby introducing noise into the training set. Moreover, aggressive recall optimisation is likely to reduce precision, which may substantially increase false alarms. The combined SMOTETomek pipeline also imposes significant computational cost in terms of both memory and runtime, particularly for larger datasets. Finally, the strong class-weight amplification may cause the classifier to overfit to minority-class patterns that do not generalise well to unseen incident reports.

\subsection{Iteration 6: Language-Aware Undersampling with Strong Regularization}

\subsubsection{Objective}
Achieve high classification performance using 100\% real data through language-aware undersampling strategies, eliminating potential artifacts introduced by synthetic data generation.

\subsubsection{Rationale}
Iteration 5's aggressive SMOTE oversampling raised concerns about synthetic sample quality in high-dimensional embedding spaces. Iteration 6 takes an alternative approach: instead of generating synthetic minority samples, it strategically selects the most informative real samples from both classes while maintaining language distribution balance.

\subsubsection{Implementation Steps}

\paragraph{1. Language-Aware Sampling Strategy}
Unlike previous iterations, which treated the available reports as a single undifferentiated corpus, Iteration~6 applies stratified sampling separately within each language group. In this way, the undersampling procedure preserves the linguistic structure of the dataset while reducing class imbalance. The corresponding language-aware undersampling function, which selects the longest incident descriptions within each language--class stratum, is provided in Listing~\ref{lst:lang-undersample} in Appendix~\ref{sec:code-listings}.

\paragraph{2. Selection of Longest Text}
The central assumption of Iteration~6 is that longer incident descriptions tend to contain richer contextual information, more explicit causal details, and clearer references to process-related hazards, which in turn support more reliable classification. To operationalize this idea, each report is assigned a simple length-based score:
\begin{equation}
\text{Score}(d_i) = \text{len}(\text{TITLE}_i \oplus \text{DESCRIPTION}_i).
\end{equation}
Within each language--class stratum, samples are ranked according to this score, and the top $N$ longest descriptions are retained for training. This procedure aims to preserve the most information-rich real observations while avoiding the introduction of synthetic samples.

\paragraph{3. No Synthetic Data}
A defining characteristic of Iteration~6 is that it relies exclusively on authentic incident reports. In contrast to earlier balancing strategies, no SMOTE oversampling, no ADASYN, and no other form of artificial data augmentation are applied. As a result, all classifier decisions are based solely on patterns observed in genuine reports rather than on interpolated or synthetically generated samples. This design was chosen in order to eliminate concerns that synthetic minority-class examples might distort the geometry of the embedding space or introduce unrealistic safety-event patterns.

\paragraph{4. Strong Regularization}
Because language-aware undersampling substantially reduces the effective size of the training set, stronger regularization is introduced to mitigate overfitting. For SVM models, this is achieved through a more restrictive regularization setting, with $C = 0.1$ rather than $C = 1.0$. For gradient boosting methods, the learning rate is reduced to 0.05 and the number of estimators is limited to 100 in order to constrain model complexity. Random Forest models are likewise regularized through more aggressive tree-pruning settings. Taken together, these adjustments are intended to improve generalization under reduced-data conditions.

\paragraph{5. RBF Kernel SVM}
Iteration~6 further extends the classifier design by introducing a non-linear SVM with an RBF kernel. The specific configuration used for this model is documented in Listing~\ref{lst:rbf-svm} in Appendix~\ref{sec:code-listings}. In contrast to linear decision functions, the RBF kernel enables the classifier to model non-linear class boundaries in the embedding space, which is particularly relevant when semantic clusters of Process Safety and Non-Process Safety incidents are not linearly separable. Formally, the kernel is defined as
\begin{equation}
K(\mathbf{x}_i, \mathbf{x}_j) = \exp\left(-\gamma \|\mathbf{x}_i - \mathbf{x}_j\|^2\right).
\end{equation}
This formulation allows the model to measure similarity in a non-linear feature space and thereby capture more complex relationships between embedded incident narratives.

\subsubsection{Results Summary}

\begin{table}[htbp]
\centering
\caption{Iteration 6 Best Results}
\label{tab:iter6-results}
\begin{tabular}{@{}llll@{}}
\toprule
\textbf{Metric} & \textbf{Value} & \textbf{Configuration} \\
\midrule
Best Macro F1 & 0.7750 & UK dataset \\
Best Model & \multicolumn{2}{l}{paraphrase-multilingual-mpnet-base-v2} \\
Best Classifier & SVM (RBF) & \\
Total Experiments & 175 & \\
Mean Macro F1 & 0.4815 & \\
\bottomrule
\end{tabular}
\end{table}

\subsubsection{Key Findings}
The empirical results indicate that Iteration~6 performs best on the English-language UK dataset, which likely reflects a combination of stronger model pre-training alignment with English and better semantic consistency in the English reports. The length-based selection strategy also appears to be meaningful, as longer descriptions were associated with higher classification confidence and more stable decision boundaries. In addition, the RBF-kernel SVM outperformed the corresponding linear SVM on the language-balanced subsets, suggesting that the decision surface in the reduced real-data setting is not adequately captured by a purely linear classifier. Finally, the use of 100\% authentic data avoids the synthetic artifacts that were a concern in Iteration~5 and therefore provides a cleaner experimental basis for assessing the real discriminatory value of the underlying text representations.

\subsubsection{Limitations}
Despite these advantages, Iteration~6 has several important limitations. First, language-aware undersampling necessarily reduces the size of the training data, which may weaken statistical robustness and increase sensitivity to local variation. Second, languages with relatively few Process Safety incidents contribute only small balanced subsets, thereby constraining the amount of usable supervision available for those strata. Third, selecting only the longest descriptions may introduce selection bias, since longer reports may systematically correspond to certain incident types, reporting cultures, or severity levels. Finally, the comparatively low mean macro F1-score of 0.4815 indicates substantial variation across configurations, suggesting that the benefits of this strategy are uneven and highly dependent on dataset and model choice.

\subsection{Iteration 7: LLM-Based Few-Shot Classification}

\subsubsection{Objective}
Evaluate large language models (LLMs) for incident classification using few-shot prompting, addressing Research Questions 2 and 3 regarding LLM comparison with traditional ML and hazard type extraction capabilities.

\subsubsection{Rationale}
Previous iterations relied on embedding-based approaches where transformer models served as feature extractors for downstream classifiers. Iteration 7 explores an alternative paradigm: using LLMs directly for classification through carefully engineered prompts with few-shot examples. This approach leverages the in-context learning capabilities of modern LLMs, potentially reducing the need for large labeled training sets.

\subsubsection{Implementation Steps}

\paragraph{1. Model Selection}
Iteration~7 evaluates two instruction-tuned variants of the Flan-T5 family, namely \texttt{google/flan-t5-large} with approximately 780 million parameters and \texttt{google/flan-t5-xl} with approximately 3 billion parameters. These models were selected because Flan-T5 combines the T5 text-to-text framework with instruction tuning, making it suitable for prompt-based classification tasks in which the model must infer labels directly from natural-language task descriptions rather than from task-specific supervised fine-tuning.

\paragraph{2. Process Safety Definition}
To guide the LLM classification procedure, a comprehensive Process Safety definition was formulated and incorporated into the prompting framework (see Listing~\ref{lst:ps-definition} in Appendix~\ref{sec:code-listings}). This definition enumerates ten categories of Process Safety incidents, including equipment failures, leaks, pressure deviations, emergency shutdowns, fire or explosion events, loss of containment, contamination, equipment damage, unplanned outages, and broader process-system harm potential, while also distinguishing four categories of Non-Process Safety incidents. The purpose of this definition is to reduce ambiguity in the classification task and provide the model with an explicit conceptual boundary between the two classes.

\paragraph{3. Few-Shot Example Selection}
Language-specific few-shot examples were curated from actual incident reports in order to support multilingual in-context learning. The examples were selected to represent diverse incident patterns within each class and to capture potentially ambiguous or borderline cases that might otherwise be difficult for the model to classify correctly. The final distribution of few-shot examples is summarised in Table~\ref{tab:few-shot-examples}.

\begin{table}[htbp]
\centering
\caption{Few-Shot Examples per Language}
\label{tab:few-shot-examples}
\begin{tabular}{@{}lccc@{}}
\toprule
\textbf{Language} & \textbf{PS Examples} & \textbf{NPS Examples} & \textbf{Total} \\
\midrule
English & 4 & 3 & 7 \\
German & 3 & 3 & 6 \\
Dutch & 2 & 3 & 5 \\
Swedish & 1 & 2 & 3 \\
\bottomrule
\end{tabular}
\end{table}

\paragraph{4. Prompt Engineering}
The few-shot prompt template combines the Process Safety definition with labeled examples in a structured natural-language instruction format (see Listing~\ref{lst:few-shot-prompt} in Appendix~\ref{sec:code-listings}). Specifically, the prompt begins with the formal Process Safety definition, followed by two Process Safety and two Non-Process Safety examples for the relevant language, and concludes with the target incident report to be classified. This design aims to provide the model with both conceptual guidance and concrete task demonstrations before it generates a binary label.

\paragraph{5. Inference Pipeline}
Batch inference was implemented with checkpointing in order to support large-scale evaluation over the full multilingual dataset. For each incident, the pipeline tokenises the constructed prompt, applies \texttt{model.generate} using greedy decoding with \texttt{num\_beams=1}, and then parses the generated response for the phrase \textit{``process safety''} in order to assign a binary label (see Listing~\ref{lst:llm-inference} in Appendix~\ref{sec:code-listings}). This pipeline enables systematic processing of the full corpus while ensuring that intermediate progress can be saved and resumed when necessary.

\paragraph{6. Multilingual Support}
To preserve language consistency during few-shot prompting, the pipeline routes each dataset file to the corresponding language-specific example bank (see Listing~\ref{lst:lang-files} in Appendix~\ref{sec:code-listings}). In this way, English incidents are paired with English few-shot examples, German incidents with German examples, and likewise for Dutch and Swedish. This multilingual routing strategy is intended to reduce cross-lingual prompt mismatch and to provide the model with demonstrations that are linguistically aligned with the target incident narrative.

\subsection{Iteration 8: Placeholder for Final Integrated Refinement}

\subsubsection{Objective}
Iteration~8 is reserved as a placeholder for a final integrated refinement stage intended to consolidate the strongest methodological components identified in the earlier experiments. The purpose of this planned iteration is not to introduce a fundamentally new modelling paradigm, but rather to combine the most effective embedding strategy, class-balancing configuration, classifier architecture, and decision-threshold setting into a single final benchmark configuration.

\subsubsection{Rationale}
The preceding iterations revealed that no single methodological modification was sufficient to achieve the target defined in RQ1. However, several components emerged as consistently promising, including stronger multilingual embedding models, mean pooling, gradient boosting classifiers, calibrated class balancing, and threshold-sensitive optimisation for Process Safety recall. Iteration~8 is therefore conceived as an integration stage in which these individually successful design choices would be combined and re-evaluated under a unified final configuration.

\subsubsection{Planned Configuration}
The planned Iteration~8 configuration would draw selectively on the strongest findings from earlier iterations. In particular, it would use the best-performing multilingual embedding model identified in Iterations~4--6, retain mean pooling as the preferred document representation strategy, employ the strongest downstream classifier architecture observed in the advanced ML experiments, and apply the most appropriate balancing strategy depending on the required precision--recall trade-off. In addition, threshold optimisation and error-focused post-hoc analysis would be incorporated in order to assess whether the final integrated system can narrow the remaining gap to the RQ1 target.

\subsubsection{Expected Role in the Thesis}
Methodologically, Iteration~8 serves as a placeholder for the final synthesis step of the experimental framework. Its intended role is to provide a single consolidated benchmark that reflects the cumulative knowledge gained from Iterations~0--7. In this sense, Iteration~8 is conceptually important even though no empirical results are reported for it in the present version of the thesis.

\subsubsection{Status}
Iteration~8 was not executed within the time frame of the present thesis. It is therefore retained as a structural placeholder only, and all quantitative results reported in this thesis are limited to Iterations~0--7.


\subsection{RQ4: Incident Severity Prediction Framework}
\label{sec:rq4_methodology}

\paragraph{Objective}
Research Question 4 asks: Can a proposed classification framework predict incident severity levels (e.g., low, medium, high) for the business case? Unlike the binary Process Safety vs.\ Non-Process Safety classification addressed in Iterations~0--7, severity prediction introduces an ordinal, multi-class dimension that poses distinct methodological challenges. Rather than implementing and benchmarking a full severity prediction pipeline, this section provides the intellectual foundation and detailed framework proposal that future work can build upon. The discussion is grounded in: (a) an empirical analysis of severity-related fields in the actual Uniper HSSE dataset, (b) a thorough review of severity prediction approaches deployed across construction, mining, energy, transportation, and healthcare industries, and (c) a proposed multi-component framework that integrates the best-performing techniques from Iterations~1--4 with ordinal-aware classification strategies.

\paragraph{Why Severity Prediction Matters}
In process safety management, not all incidents are equal. A minor gas turbine trip that is quickly resolved carries fundamentally different risk implications than a loss-of-containment event leading to a fire or explosion. Yet the binary PS/NPS classifier developed in earlier iterations treats all process safety incidents identically, regardless of their actual or potential consequence. Severity prediction enables three practical capabilities that binary classification cannot: (1) priority-based triage, where critical incidents are immediately escalated to senior management while minor events follow standard review workflows; (2) resource allocation optimisation, where investigation teams and remediation budgets are directed proportionally to incident severity; and (3) regulatory compliance, where organisations can demonstrate risk-proportionate response capabilities to auditors and regulators.

% -------------------------------------------------------
% 2. DATASET ANALYSIS
% -------------------------------------------------------
\paragraph{Severity-Related Fields in the Uniper Dataset}
The Process Safety subset of the Uniper HSSE dataset (\texttt{RW\_ACTUALS\_PS}) contains 14{,}432 records with 36 columns, including four distinct fields that capture severity information, each with different coverage, granularity, and interpretation.

\begin{table}[H]
\centering
\caption{Severity-related fields in the Uniper HSSE dataset}
\label{tab:severity_fields}
\begin{tabular}{llll}
\toprule
\textbf{Field} & \textbf{Coverage} & \textbf{Values} & \textbf{Notes} \\
\midrule
\texttt{CASE\_SEVERITY}           & 100\%  & L2--L6, Not Selected & Primary field; L2=lowest, L6=highest. 14.8\% unlabelled. \\
\texttt{POTENTIAL\_SEV\_LEVEL}    & 48.5\% & Very Low--Very High  & Prospective assessment; over half missing. \\
\texttt{LEARNINGS\_ACTUAL\_SEV}   & 23.8\% & Low, Medium, Severe  & Post-investigation; only populated for investigated cases. \\
\texttt{RISK\_AREA}               & 44.3\% & Low/Medium/High Risk & Risk category; missing for the majority of records. \\
\bottomrule
\end{tabular}
\end{table}

The \texttt{CASE\_SEVERITY} field provides the most complete severity annotation. Its distribution is heavily skewed toward the lowest severity level (L2), which accounts for 67.6\% of all records. Higher-severity incidents (L4--L6) collectively represent only 7.2\% of the data, presenting a severe class imbalance challenge.

\paragraph{Proposed Three-Tier Severity Mapping}
For practical deployment, the original five severity levels are consolidated into a three-tier scheme that balances granularity with statistical feasibility:

\begin{table}[H]
\centering
\caption{Three-tier severity mapping for the Uniper dataset}
\label{tab:three_tier}
\begin{tabular}{llll}
\toprule
\textbf{Tier} & \textbf{Original Levels} & \textbf{Count (\%)} & \textbf{Business Interpretation} \\
\midrule
Low    & L2         & 9{,}757 (79.4\%) & Near-misses, minor equipment issues, no significant consequence. \\
Medium & L3         & 1{,}494 (12.2\%) & Contained releases, equipment damage requiring repair, partial shutdowns. \\
High   & L4, L5, L6 & 1{,}043 (8.5\%)  & Fires, explosions, LOPC, injuries, emergency shutdowns, regulatory events. \\
\bottomrule
\end{tabular}
\end{table}

This mapping yields 12{,}294 labelled records (excluding 2{,}138 \textit{Not Selected} entries). The 79:12:9 class ratio presents a severe imbalance challenge that must be addressed through the SMOTE and cost-sensitive strategies validated in Iterations~4--6. An important empirical observation is that higher-severity incidents tend to have longer textual descriptions: more severe events trigger more detailed reporting, formal investigation documentation, and comprehensive corrective action descriptions. Germany (DE) contributes 43.3\% of records, followed by the UK (28.2\%) and Netherlands (15.8\%); notably, the UK shows a disproportionate share of L6 incidents (177 of 226 total), which must be accounted for in model training.

% -------------------------------------------------------
% 3. LITERATURE REVIEW
% -------------------------------------------------------
\paragraph{Related Work: Severity Prediction Across Industries}

\textit{Traditional ML approaches.}
Support Vector Machines and gradient boosting methods have been extensively applied to incident severity classification. Sattari et al.~\cite{sattari2021} demonstrated that Linear SVM achieved 85.68\% accuracy for process safety management element classification in oil and gas incident reports. More recent work~\cite{SR2025} developed ML models for severity prediction in Saudi Arabian construction sites, testing six algorithms on 203 incidents; XGBoost achieved 89\% accuracy, confirming the superiority of gradient boosting methods. A feature-optimised Random Forest trained on 66{,}405 OSHA records~\cite{PMC2022} achieved 0.895 accuracy, with feature importance analysis revealing that structured metadata fields (analogous to the HAZARD and FUNCTIONAL\_AREA fields in the Uniper dataset) carry substantial predictive power. An explainable ML pipeline~\cite{SD2025} applied to 2{,}645 OSHA mining investigation reports achieved Random Forest F1 of 0.94 and AUC of 0.95, integrating SMOTETomek balancing, Bayesian hyperparameter optimisation, and SHAP explanations.

\textit{Deep learning approaches.}
CNNs applied to incident text use convolutional filters to detect severity-indicative n-grams such as ``explosion'', ``hospitalisation'', or ``emergency shutdown''~\cite{wang2025bert}. Bidirectional LSTM networks process incident text sequentially, capturing how narrative context evolves throughout the description; an attention layer learns to weight the most severity-relevant portions of the text. A hybrid CNN-RNN model~\cite{ArXiv2024} achieved the best performance among deep learning approaches for crash severity classification by combining CNN's local feature extraction with RNN's sequential context modelling.

\textit{Transformer-based approaches.}
Fine-tuned BERT achieved F-scores above 89\% for classifying 10 incident types and 4 severity levels in patient safety reports, significantly outperforming CNNs~\cite{wang2025bert}. A German-language study~\cite{MDPI2025} applied GBERT to classify 617 critical incident reports from a university hospital into risk categories, outperforming TF-IDF baselines despite a small, imbalanced dataset. BERT-based NLP models for insurance claim severity modelling~\cite{Variance2023} showed that contextual embeddings improve severity prediction stability compared to traditional NLP features.

\textit{LLM-based approaches.}
Nouri et al.~\cite{Nouri2024} applied a GPT-4-based pipeline for hazard analysis and risk assessment in automotive systems, emphasising the ongoing necessity of expert validation. GPT-based models for process safety management in oil and gas operations~\cite{SPE2024} achieved over 95\% reduction in analysis time for maintenance action classification, though direct zero-shot classification underperformed trained models---consistent with the Iteration~7 finding of this thesis. The most promising direction is hybrid approaches where LLMs serve as feature extractors feeding into trained classifiers, mirroring the successful \textit{embeddings~+ classifier} paradigm that achieved 0.89 macro F1 in Iteration~4.

\textit{Ordinal regression.}
Standard multi-class classification treats severity labels as unordered categories, penalising a misclassification of High as Low identically to High as Medium. The COnsistent RAnk Logits (CORAL) framework~\cite{Cao2020} addresses this by transforming ordinal targets into binary subtasks while guaranteeing rank consistency. The subsequent CORN framework~\cite{Shi2023} improved upon CORAL by predicting conditional probabilities, achieving better performance on multiple benchmarks. A PMC study~\cite{PMC2021} on imbalanced crash severity data used one-vs-all binary decomposition with isotonic regression for calibration and per-class threshold optimisation, significantly outperforming standard classification.

% -------------------------------------------------------
% 4. PROPOSED FRAMEWORK
% -------------------------------------------------------
\paragraph{Proposed Framework Architecture}
The proposed framework extends the binary classification pipeline from Iterations~1--4 into a multi-task architecture with three output heads: (1)~PS vs.\ Non-PS classification (the existing task), (2)~severity prediction (Low/Medium/High), and (3)~hazard type identification. The shared transformer backbone enables transfer learning between tasks. The framework comprises five modular components, each of which can be evaluated independently or combined incrementally.

\textit{Component 1: Feature Engineering.}
Four feature categories are combined. Text embeddings are generated by concatenating TITLE and CASE\_DESCRIPTION and encoding with a multilingual transformer (E5-large-instruct or paraphrase-multilingual-mpnet-base-v2) using mean pooling, producing 768-dimensional vectors. Binary keyword features are extracted via regex matching across all four languages (English, German, Dutch, Swedish), covering fire/explosion terms (fire, brand, Feuer, explosion), loss-of-containment terms (leak, spill, Leckage, lekkage), emergency shutdown terms (ESD, Not-Aus, noodstop), injury terms (injury, fatality, Verletzung), and equipment failure terms (crack, rupture, corrosion). Structured metadata features include HAZARD, FUNCTIONAL\_AREA, COUNTRY\_SHORT, and COMPANY\_INVOLVED\_TYPE, one-hot encoded as categorical features. Text statistics features include description length (characters and words), title length, and sentence count.

\textit{Component 2: Traditional ML Classifiers.}
The following multi-class classifiers are recommended based on Iteration~4 results:

\begin{table}[H]
\centering
\caption{Recommended traditional ML classifiers for severity prediction}
\label{tab:severity_classifiers}
\begin{tabular}{lll}
\toprule
\textbf{Classifier} & \textbf{Key Configuration} & \textbf{Rationale} \\
\midrule
LightGBM        & \texttt{objective='multiclass'}, \texttt{n\_estimators=300}, \texttt{lr=0.05} & Best performer in Iteration~4; native multi-class. \\
XGBoost         & \texttt{objective='multi:softprob'}, \texttt{n\_estimators=300}              & Calibrated probability outputs. \\
Random Forest   & \texttt{n\_estimators=300}, \texttt{class\_weight='balanced\_subsample'}      & SHAP-compatible; interpretable importances. \\
Soft Voting Ensemble & Probability-weighted vote over LightGBM + XGBoost + RF              & Leverages probability estimates for ordinal decisions. \\
\bottomrule
\end{tabular}
\end{table}

\textit{Component 3: Deep Learning Approaches.}
A 1D Text-CNN with multiple filter sizes (3, 4, 5 tokens) captures severity-indicative n-grams with low computational cost. A BiLSTM with attention processes incident narratives bidirectionally and highlights which tokens most influenced the severity prediction, providing both performance and interpretability. Fine-tuned multilingual BERT (starting from multilingual-e5-large-instruct, as identified in Iteration~3) adapts internal representations specifically for severity discrimination using an ordinal-aware loss function.

\textit{Component 4: Ordinal Regression Integration.}
The CORAL framework is applied to neural network classifiers by transforming the three-class prediction into two binary subtasks---\textit{Is severity $>$ Low?} and \textit{Is severity $>$ Medium?}---with shared backbone weights and independent output bias terms, guaranteeing rank consistency. For gradient boosting classifiers, a cumulative binary decomposition trains two independent binary classifiers whose outputs are combined to assign the final severity tier, enabling independent threshold optimisation for the High-severity boundary.

\textit{Component 5: LLM-Based Feature Enrichment.}
Based on the Iteration~7 finding that direct few-shot LLM classification fails for this domain, the recommended LLM strategy is feature enrichment rather than direct classification. LLMs extract structured severity indicators from unstructured text (e.g., \textit{Were people injured? Was there a fire? Was containment lost?}) which become binary input features for gradient boosting classifiers. For very long incident descriptions, LLMs generate severity-focused summaries that fit within BERT's 512-token limit. LLMs can also generate synthetic incident descriptions for the rare High severity class, subject to domain expert validation.

% -------------------------------------------------------
% 5. CLASS BALANCING
% -------------------------------------------------------
\paragraph{Class Balancing Strategy}
The 79:12:9 class distribution presents a more severe imbalance challenge than the binary PS/NPS classification. Table~\ref{tab:severity_balancing} adapts the strategies validated in Iterations~4--6 to the three-class setting.

\begin{table}[H]
\centering
\caption{Class balancing strategies adapted from Iterations~4--6 for severity prediction}
\label{tab:severity_balancing}
\begin{tabular}{llll}
\toprule
\textbf{Strategy} & \textbf{Iteration Result} & \textbf{Severity Adaptation} & \textbf{Expected Outcome} \\
\midrule
Conservative SMOTE (0.7) & Best overall F1 (0.89) & Oversample Medium and High to 70\% of Low & Best precision-recall balance \\
Aggressive SMOTE (1.0)   & High recall (0.92), lower precision & Full 1:1:1 balancing & Maximum High recall \\
Language-aware undersample & High precision (0.81) & Undersample Low per language stratum & Preserves language balance \\
Cost-sensitive weights   & Used across all iterations & $w_\text{High} = 10 \times N_\text{total} / (3 \times N_\text{High})$ & Penalise missed High incidents \\
\bottomrule
\end{tabular}
\end{table}

% -------------------------------------------------------
% 6. EXPLAINABILITY
% -------------------------------------------------------
\paragraph{Explainability for Safety-Critical Severity Predictions}
In safety-critical applications, a severity prediction is only useful if stakeholders can understand and trust it. Two complementary explainability approaches are recommended. SHAP values decompose each prediction into additive feature contributions, showing exactly how much each input feature pushed the prediction toward a particular severity class; for gradient boosting classifiers, SHAP computation is extremely efficient via TreeSHAP. For transformer-based models, attention weights identify which tokens in the incident description the model focused on when making its prediction, enabling safety managers to verify that the model attends to consequence descriptions rather than spurious correlations. Calibrated probability outputs enable a confidence-based routing system: predictions with maximum class probability above~0.80 are auto-accepted; predictions between~0.50 and~0.80 are presented as suggestions for human review; predictions below~0.50 are flagged as uncertain and routed entirely to human classification.

% -------------------------------------------------------
% 7. IMPLEMENTATION ROADMAP
% -------------------------------------------------------
\paragraph{Implementation Roadmap}
The proposed framework is designed for incremental implementation across four phases.

\textit{Phase~1 --- Baseline Severity Classifier (2--4 weeks):} Reuse Iteration~4 embeddings (mean pooling, multilingual E5/mBERT) with \texttt{CASE\_SEVERITY} as the three-tier target. Train LightGBM, XGBoost, and Random Forest with SMOTE~(0.7) balancing. Evaluate with stratified 5-fold cross-validation and integrate SHAP explanations.

\textit{Phase~2 --- Ordinal Regression Enhancement (2--3 weeks):} Implement cumulative binary decomposition and CORAL output layer. Optimise per-boundary thresholds to maximise High-severity recall ($\geq 0.80$).

\textit{Phase~3 --- Deep Learning Models (3--4 weeks):} Implement and evaluate Text-CNN, BiLSTM+Attention, and fine-tuned multilingual BERT. Integrate ordinal loss functions (CORAL/CORN) with deep learning architectures.

\textit{Phase~4 --- LLM Feature Enrichment and Multi-Task Learning (2--3 weeks):} Develop LLM-based structured severity indicator extraction. Implement multi-task learning combining PS classification and severity prediction. Evaluate data augmentation via LLM-generated synthetic High-severity examples.

% -------------------------------------------------------
% 8. EVALUATION STRATEGY
% -------------------------------------------------------
\paragraph{Evaluation Strategy and Metrics}

\begin{table}[H]
\centering
\caption{Primary evaluation metrics for the severity prediction framework}
\label{tab:severity_metrics}
\begin{tabular}{llll}
\toprule
\textbf{Metric} & \textbf{Target} & \textbf{Purpose} & \textbf{Rationale} \\
\midrule
Macro F1-Score       & $\geq 0.70$ & Balanced per-class performance  & Primary metric per RQ4; treats all severity classes equally. \\
Overall Accuracy     & $\geq 0.70$ & General correctness             & Complements F1; can be misleading with imbalanced data. \\
Recall (High)        & $\geq 0.80$ & Minimise missed critical incidents & Safety-critical: missing High severity is worse than a false alarm. \\
Mean Absolute Error  & $< 0.10$    & Ordinal error magnitude         & Penalises extreme misclassifications (Low predicted as High). \\
Quadratic Weighted Kappa & $\geq 0.60$ & Ordinal agreement           & Measures prediction-label agreement weighted by squared distance. \\
Adjacent Accuracy    & $\geq 0.90$ & Extreme misclassification rate  & Proportion of predictions within one severity level of true label. \\
\bottomrule
\end{tabular}
\end{table}

% -------------------------------------------------------
% 9. EXPECTED OUTCOMES
% -------------------------------------------------------
\paragraph{Expected Outcomes}

\begin{table}[H]
\centering
\caption{Projected outcomes per implementation phase for RQ4}
\label{tab:rq4_outcomes}
\begin{tabular}{llllll}
\toprule
\textbf{Phase} & \textbf{Macro F1} & \textbf{Accuracy} & \textbf{Recall(High)} & \textbf{MAE} & \textbf{Meets RQ4?} \\
\midrule
Phase~1: Baseline ML    & 0.68--0.74 & 0.72--0.78 & 0.65--0.75 & 0.30--0.40 & Possible \\
Phase~2: Ordinal        & 0.70--0.76 & 0.74--0.80 & 0.72--0.82 & 0.25--0.35 & Likely \\
Phase~3: Deep Learning  & 0.73--0.80 & 0.76--0.82 & 0.75--0.85 & 0.20--0.30 & Likely \\
Phase~4: Multi-Task     & 0.75--0.82 & 0.78--0.84 & 0.78--0.88 & 0.18--0.25 & Likely \\
\bottomrule
\end{tabular}
\end{table}

The projections are conservative because severity prediction is inherently harder than binary classification: three classes introduce more confusion boundaries, the High severity class contains only approximately 1{,}043 samples, and severity labels carry inherent subjectivity from different reporters and countries. Key risk factors include insufficient High-severity samples for reliable deep learning training, label noise from inconsistent severity assessment across countries, and the possibility that severity information is underspecified in brief incident descriptions.

% -------------------------------------------------------
% 10. LIMITATIONS
% -------------------------------------------------------
\paragraph{Limitations}

\textit{Data limitations.} Different reporters and countries apply severity ratings inconsistently; no inter-annotator agreement study exists for the \texttt{CASE\_SEVERITY} field. The High severity class contains only 1{,}043 examples, yielding approximately 830 training examples after splitting, which may be insufficient for deep learning approaches. The 14.8\% of records with \texttt{CASE\_SEVERITY} = `Not Selected' are excluded; if these are systematically different from labelled records, the model may not generalise to incoming unlabelled data.

\textit{Methodological limitations.} The mapping of L2--L6 to Low/Medium/High assumes equal spacing between categories, which may not hold in practice. Reporting practices and severity assessment criteria may change over time, creating temporal drift. Generating SMOTE synthetic samples in a 768-dimensional embedding space may create unrealistic data points that introduce noise.

\textit{Evaluation limitations.} The \texttt{CASE\_SEVERITY} labels are assumed correct without independent re-annotation. Unless full $k$-fold cross-validation is implemented, results may be sensitive to the particular train/test split. Calibrated probability thresholds may not generalise across deployment contexts.

% ============================================================================
% CHAPTER: RESULTS
% ============================================================================
\chapter{Results}

This chapter presents the experimental results reported for Iterations~0--7 of the classification framework. In addition, a structured placeholder is included for Iteration~8 as the reserved final refinement stage. Results are organized by iteration, with comparative analyses across models, datasets, and techniques.

\section{Iteration 0: Baseline Results}

The baseline BERT-base-uncased + Linear SVM model established the performance benchmark for subsequent improvements.

\begin{table}[htbp]
\centering
\caption{Iteration 0 baseline performance across the three English subset construction strategies.}
\label{tab:results-iter0}
\begin{tabular}{lrrrr}
\toprule
\textbf{Subset} & \textbf{Total} & \textbf{Accuracy} & \textbf{Macro-average F1} & \textbf{Weighted-F1} \\
\midrule
manual     & 4,713 & 0.8112 & \textbf{0.7825} & 0.8143 \\
langdetect & 6,010 & 0.8128 & 0.7648 & 0.8170 \\
lingua     & 6,222 & 0.8185 & 0.7714 & \textbf{0.8241} \\
\bottomrule
\end{tabular}
\end{table}

\textbf{Key Observations:}
\begin{itemize}
    \item The manually curated English subset achieved the strongest class-balanced baseline with a macro F1-score of 0.7825.
    \item The \texttt{lingua}-based subset achieved the highest accuracy (0.8185) and weighted-F1 (0.8241), but not the highest macro-F1.
    \item The gap from the official Iteration~0 baseline to the RQ1 target of 0.85 is 0.0675 macro-F1 points.
\end{itemize}

\section{Iteration 1: Country-Stratified Results}

BERT-base-cased embeddings were evaluated across six country-specific datasets.

\begin{table}[htbp]
\centering
\caption{Iteration~1 results by dataset (BERT-base-cased + Linear SVM).}
\label{tab:results-iter1}
\begin{tabular}{@{}lcccc@{}}
\toprule
\textbf{Dataset} & \textbf{Samples} & \textbf{Accuracy} & \textbf{Macro F1} & \textbf{FNR} \\
\midrule
Germany     & 14,951 & 0.7379 & 0.5983 & 0.3313 \\
UK          & 4,673  & 0.7722 & 0.7449 & 0.2518 \\
English     & 4,713  & 0.7699 & 0.7420 & 0.2624 \\
Sweden      & 12,224 & 0.7947 & 0.5825 & 0.3548 \\
Netherlands & 1,990  & 0.6985 & 0.6586 & 0.3983 \\
Master      & 34,576 & 0.7514 & 0.6292 & 0.3279 \\
\bottomrule
\end{tabular}
\end{table}

\textbf{Key Findings:}
\begin{itemize}
    \item No Iteration~1 dataset reached the RQ1 target of macro F1 $\geq 0.85$.
    \item The strongest result was obtained on the UK dataset with a macro-F1 score of 0.7449.
    \item The manually curated English dataset performed similarly, with a macro-F1 score of 0.7420.
    \item The non-English datasets showed substantially weaker class-balanced performance, particularly Sweden and Germany.
    \item These results indicate that country-stratified evaluation is informative, but that the BERT-base-cased + Linear SVM configuration did not yet provide a sufficiently strong multilingual baseline in Iteration~1.
\end{itemize}

\section{Iteration 2: Multi-Model Comparison}

Iteration~2 extends the Iteration~1 baseline by evaluating six completed transformer architectures across four manually derived datasets (\texttt{english\_manual}, \texttt{german\_manual}, \texttt{swedish\_manual}, and \texttt{dutch\_manual}), yielding 24 model--dataset combinations. Four of the originally planned ten models (\texttt{modernbert-base}, \texttt{modernbert-large}, \texttt{deberta-v3-base}, \texttt{deberta-v3-large}) encountered compatibility issues and were excluded. All per-run metrics are exported to \texttt{Results/\_iteration\_2/all\_models\_svm\_detailed\_results.csv}.

\begin{table}[htbp]
\centering
\caption{Iteration~2 model performance summary: mean $\pm$ standard deviation across four datasets (\texttt{english\_manual}, \texttt{german\_manual}, \texttt{swedish\_manual}, \texttt{dutch\_manual}).}
\label{tab:results-iter2-models}
\begin{tabular}{@{}lcccc@{}}
\toprule
\textbf{Model} & \textbf{Params} & \textbf{Accuracy} & \textbf{Macro F1} & \textbf{PS F1} \\
\midrule
bert-base-cased       & 110M & 0.7540 $\pm$ 0.0585 & 0.6230 $\pm$ 0.0696 & 0.4135 $\pm$ 0.1713 \\
bert-large-cased      & 340M & 0.7499 $\pm$ 0.0874 & 0.6005 $\pm$ 0.0548 & 0.3698 $\pm$ 0.1414 \\
roberta-base          & 125M & 0.7603 $\pm$ 0.0413 & 0.6538 $\pm$ 0.0768 & 0.4726 $\pm$ 0.1714 \\
roberta-large         & 355M & \textbf{0.7830 $\pm$ 0.0465} & \textbf{0.6654 $\pm$ 0.0698} & \textbf{0.4787 $\pm$ 0.1683} \\
distilbert-base-cased &  66M & 0.7723 $\pm$ 0.0354 & 0.6518 $\pm$ 0.0689 & 0.4562 $\pm$ 0.1644 \\
xlnet-base-cased      & 110M & 0.7279 $\pm$ 0.0794 & 0.5873 $\pm$ 0.0584 & 0.3608 $\pm$ 0.1502 \\
\bottomrule
\end{tabular}
\end{table}

\begin{table}[htbp]
\centering
\caption{Iteration~2 per-dataset averages across all six completed models.}
\label{tab:results-iter2-datasets}
\begin{tabular}{@{}lccc@{}}
\toprule
\textbf{Dataset} & \textbf{Accuracy} & \textbf{Macro F1} & \textbf{PS F1} \\
\midrule
english\_manual  & 0.7621 $\pm$ 0.0310 & 0.7342 $\pm$ 0.0333 & 0.6481 $\pm$ 0.0426 \\
german\_manual   & 0.7738 $\pm$ 0.0187 & 0.6086 $\pm$ 0.0302 & 0.3544 $\pm$ 0.0525 \\
swedish\_manual  & 0.8265 $\pm$ 0.0224 & 0.5588 $\pm$ 0.0067 & 0.2153 $\pm$ 0.0196 \\
dutch\_manual    & 0.6692 $\pm$ 0.0394 & 0.6197 $\pm$ 0.0515 & 0.4833 $\pm$ 0.0792 \\
\bottomrule
\end{tabular}
\end{table}

\textbf{Key Findings:}
\begin{itemize}
    \item RoBERTa-large achieves the highest average macro F1 (0.6654) and accuracy (0.7830) across all four datasets, with roberta-base as a close second (macro F1 = 0.6538).
    \item The best single model--dataset combination is \texttt{roberta-base} on \texttt{english\_manual}, achieving a macro F1 of 0.7698 and accuracy of 0.7953.
    \item No Iteration~2 combination reaches the RQ1 target of macro F1 $\geq 0.85$; the English-only best (0.7698) also falls below the Iteration~0 baseline of 0.7825, indicating that switching from BERT-base-uncased to cased variants and RoBERTa architectures alone is insufficient.
    \item Performance on non-English datasets is substantially weaker, particularly for Swedish (average macro F1 = 0.5588 across models), reflecting the language mismatch of English-pretrained embeddings applied to multilingual incident reports.
    \item DistilBERT (66M parameters) achieves competitive average macro F1 (0.6518) relative to larger models, making it an efficient alternative when computational resources are constrained.
    \item The large standard deviations in PS F1 (0.14--0.17) across all models confirm that Process Safety recall is a persistent challenge on non-English datasets.
\end{itemize}

\section{Iteration 3: Multilingual Model Results}

Five multilingual transformer models were evaluated across four language-specific datasets (\texttt{english\_manual}, \texttt{german\_manual}, \texttt{swedish\_manual}, \texttt{dutch\_manual}), yielding 20 model--dataset combinations.

\begin{table}[htbp]
\centering
\caption{Iteration~3 full results: accuracy, macro F1, and PS F1 for all five multilingual models across four datasets.}
\label{tab:results-iter3}
\small
\begin{tabular}{@{}llcccc@{}}
\toprule
\textbf{Model} & \textbf{Dataset} & \textbf{Accuracy} & \textbf{Macro F1} & \textbf{Prec PS} & \textbf{Rec PS} \\
\midrule
\multirow{4}{*}{mBERT-uncased}
 & english\_manual  & 0.7656 & 0.7412 & 0.5822 & 0.7660 \\
 & german\_manual   & 0.8024 & 0.6553 & 0.3163 & 0.6717 \\
 & swedish\_manual  & 0.8090 & 0.5583 & 0.1518 & 0.4387 \\
 & dutch\_manual    & 0.6960 & 0.6672 & 0.4908 & 0.6780 \\
\midrule
\multirow{4}{*}{mBERT-cased}
 & english\_manual  & 0.7550 & 0.7262 & 0.5718 & 0.7199 \\
 & german\_manual   & 0.8031 & 0.6421 & 0.3032 & 0.5964 \\
 & swedish\_manual  & 0.8245 & 0.5453 & 0.1337 & 0.3226 \\
 & dutch\_manual    & 0.6608 & 0.6196 & 0.4430 & 0.5593 \\
\midrule
\multirow{4}{*}{E5-base}
 & english\_manual  & 0.7741 & \textbf{0.7456} & 0.6000 & 0.7340 \\
 & german\_manual   & 0.8101 & 0.6534 & 0.3179 & 0.6205 \\
 & swedish\_manual  & 0.8474 & 0.5668 & 0.1615 & 0.3355 \\
 & dutch\_manual    & 0.7085 & 0.6674 & 0.5071 & 0.6017 \\
\midrule
\multirow{4}{*}{E5-large}
 & english\_manual  & 0.7646 & 0.7354 & 0.5862 & 0.7234 \\
 & german\_manual   & 0.8342 & 0.6517 & 0.3340 & 0.4970 \\
 & swedish\_manual  & \textbf{0.8822} & 0.5660 & 0.1724 & 0.2258 \\
 & dutch\_manual    & 0.6859 & 0.6168 & 0.4685 & 0.4407 \\
\midrule
\multirow{4}{*}{E5-large-instruct}
 & english\_manual  & 0.7731 & 0.7407 & 0.6037 & 0.7021 \\
 & german\_manual   & 0.8352 & 0.6700 & 0.3517 & 0.5753 \\
 & swedish\_manual  & 0.8720 & 0.5670 & 0.1681 & 0.2581 \\
 & dutch\_manual    & 0.7462 & 0.6995 & 0.5691 & 0.5932 \\
\bottomrule
\end{tabular}
\end{table}

\begin{table}[htbp]
\centering
\caption{Iteration~3 model performance summary: mean $\pm$ standard deviation across four datasets.}
\label{tab:results-iter3-models}
\begin{tabular}{@{}lccc@{}}
\toprule
\textbf{Model} & \textbf{Accuracy} & \textbf{Macro F1} & \textbf{PS F1} \\
\midrule
mBERT-uncased       & 0.7683 $\pm$ 0.0449 & 0.6555 $\pm$ 0.0650 & 0.4716 $\pm$ 0.1643 \\
mBERT-cased         & 0.7609 $\pm$ 0.0630 & 0.6333 $\pm$ 0.0645 & 0.4307 $\pm$ 0.1628 \\
E5-base             & 0.7851 $\pm$ 0.0512 & 0.6583 $\pm$ 0.0634 & 0.4623 $\pm$ 0.1646 \\
E5-large            & 0.7917 $\pm$ 0.0740 & 0.6425 $\pm$ 0.0617 & 0.4242 $\pm$ 0.1610 \\
E5-large-instruct   & \textbf{0.8066 $\pm$ 0.0496} & \textbf{0.6693 $\pm$ 0.0642} & \textbf{0.4676 $\pm$ 0.1707} \\
\bottomrule
\end{tabular}
\end{table}

\textbf{Key Findings:}
\begin{itemize}
    \item \texttt{intfloat/multilingual-e5-large-instruct} achieves the highest average macro F1 (0.6693) and accuracy (0.8066) across all four datasets.
    \item The best single combination is \texttt{E5-base} on \texttt{english\_manual} with macro F1 = 0.7456 and accuracy = 0.7741.
    \item No Iteration~3 combination reaches the RQ1 target of macro F1 $\geq 0.85$.
    \item Compared to Iteration~2, multilingual models show marginal improvement on non-English datasets: German average macro F1 rises from 0.6086 (Iteration~2) to 0.6545 (Iteration~3), and Dutch from 0.6197 to 0.6541.
    \item Swedish remains the hardest dataset across all models (average macro F1 = 0.5607), driven by severe class imbalance and low PS recall (0.23--0.44).
    \item English performance is virtually unchanged between iterations (0.7342 in Iteration~2 vs.\ 0.7378 in Iteration~3), confirming that multilingual pretraining does not hurt English classification.
    \item Scaling from E5-base to E5-large does not consistently improve macro F1, suggesting that model size alone is not the bottleneck.
\end{itemize}

\section{Iteration 4: Advanced ML Techniques}

Iteration~4 evaluated eight multilingual embedding models combined with four classifiers (LinearSVC, Random Forest, XGBoost, LightGBM) across four language-specific datasets, yielding 128 model--classifier--dataset combinations. All experiments used mean pooling, SMOTE oversampling (ratio 0.7), feature scaling, and hyperparameter tuning via grid search. The Voting Ensemble was configured in code but its results were not captured in the output CSV, so only the four base classifiers are reported.

\begin{table}[htbp]
\centering
\caption{Iteration~4 classifier comparison: mean accuracy, macro F1, PS F1, and PS recall averaged across all 8 embedding models and 4 datasets (32 runs per classifier).}
\label{tab:results-iter4-classifiers}
\begin{tabular}{@{}lcccc@{}}
\toprule
\textbf{Classifier} & \textbf{Accuracy} & \textbf{Macro F1} & \textbf{PS F1} & \textbf{PS Recall} \\
\midrule
LinearSVC     & 0.7610 & 0.6421 & 0.4449 & 0.5865 \\
Random Forest & 0.8431 & 0.6810 & 0.4598 & 0.4001 \\
LightGBM      & 0.8329 & 0.6973 & 0.5021 & 0.5232 \\
XGBoost       & \textbf{0.8366} & \textbf{0.7011} & \textbf{0.5072} & \textbf{0.5231} \\
\bottomrule
\end{tabular}
\end{table}

\begin{table}[htbp]
\centering
\caption{Iteration~4 best result per dataset and overall top-10 combinations.}
\label{tab:results-iter4-best}
\begin{tabular}{@{}llllccc@{}}
\toprule
\textbf{Dataset} & \textbf{Embedding model} & \textbf{Clf} & \textbf{Macro F1} & \textbf{Acc} & \textbf{Prec PS} & \textbf{Rec PS} \\
\midrule
english\_manual  & snowflake-arctic-embed-l-v2.0           & XGBoost & \textbf{0.8043} & 0.8378 & 0.7380 & 0.7092 \\
german\_manual   & intfloat/multilingual-e5-large-instruct & XGBoost & 0.7426          & 0.8927 & 0.5147 & 0.5813 \\
dutch\_manual    & Alibaba-NLP/gte-multilingual-base       & XGBoost & 0.7439          & 0.7940 & 0.6731 & 0.5932 \\
swedish\_manual  & Alibaba-NLP/gte-multilingual-base       & XGBoost & 0.6438          & 0.8920 & 0.2811 & 0.4516 \\
\bottomrule
\end{tabular}
\end{table}

\begin{table}[htbp]
\centering
\caption{Iteration~4 per-dataset averages across all 32 model--classifier combinations.}
\label{tab:results-iter4-datasets}
\begin{tabular}{@{}lccc@{}}
\toprule
\textbf{Dataset} & \textbf{Accuracy} & \textbf{Macro F1} & \textbf{PS F1} \\
\midrule
english\_manual  & 0.8050 & \textbf{0.7629} & \textbf{0.6636} \\
dutch\_manual    & 0.7519 & 0.6865 & 0.5448 \\
german\_manual   & 0.8522 & 0.6761 & 0.4382 \\
swedish\_manual  & 0.8645 & 0.5960 & 0.2673 \\
\bottomrule
\end{tabular}
\end{table}

\textbf{Key Findings:}
\begin{itemize}
    \item XGBoost achieves the highest average macro F1 (0.7011) across all model--dataset combinations, followed closely by LightGBM (0.6973); both substantially outperform LinearSVC (0.6421).
    \item The best single combination is \texttt{snowflake-arctic-embed-l-v2.0} + XGBoost on \texttt{english\_manual}, achieving macro F1 = 0.8043 and accuracy = 0.8378 — the strongest result across all iterations up to this point.
    \item No Iteration~4 combination reaches the RQ1 target of macro F1 $\geq 0.85$; the best result (0.8043) is 2.43 percentage points below the target.
    \item English remains the strongest dataset (average macro F1 = 0.7629); Swedish is consistently the weakest (0.5960), driven by severe class imbalance and low PS recall.
    \item Random Forest achieves high accuracy (0.8431) but low PS recall (0.4001), indicating a strong bias toward the majority class despite SMOTE oversampling.
    \item LinearSVC shows the lowest accuracy (0.7610) but the highest PS recall (0.5865) among classifiers, reflecting a more aggressive minority-class bias.
    \item The Voting Ensemble results were not captured in the output CSV and are therefore not reported.
\end{itemize}

\section{Iteration 5 and 6: Class Balancing Experiments}

\begin{table}[htbp]
\centering
\caption{Iterations 5-6 Class Balancing Strategy Comparison}
\label{tab:results-iter5-6}
\begin{tabular}{@{}lcccc@{}}
\toprule
\textbf{Strategy} & \textbf{Macro F1} & \textbf{PS Recall} & \textbf{PS Precision} & \textbf{Approach} \\
\midrule
Iter 4 (SMOTE 0.7) & 0.8043 & 0.7092 & 0.7380 & Synthetic \\
Iter 5 (Full SMOTE) & TBD & TBD & TBD & Synthetic \\
Iter 6 (Undersampling) & TBD & TBD & TBD & Real only \\
\bottomrule
\end{tabular}
\end{table}

\textbf{Key Findings:}
\begin{itemize}
    \item Iteration~4 best result (macro F1 = 0.8043, \texttt{snowflake-arctic-embed-l-v2.0} + XGBoost on \texttt{english\_manual}) serves as the reference for class balancing comparisons.
    \item Results for Iterations~5 and~6 are pending completion of those pipelines and will be filled in during final revision.
\end{itemize}

\section{Iteration 7: LLM Classification Results}

\begin{table}[htbp]
\centering
\caption{Iteration 7 LLM vs. Traditional ML Comparison}
\label{tab:results-iter7}
\begin{tabular}{@{}lcccc@{}}
\toprule
\textbf{Approach} & \textbf{Accuracy} & \textbf{Macro F1} & \textbf{PS Recall} & \textbf{Training} \\
\midrule
BERT-base-uncased + RBF SVM (Iter~0)   & 0.8112 & 0.7825 & 74.82\% & Required \\
roberta-base + LinearSVC (Iter~2)      & 0.7953 & 0.7698 & 77.30\% & Required \\
snowflake-arctic + XGBoost (Iter~4)    & 0.8378 & 0.8043 & 70.92\% & Required \\
Flan-T5-Large (few-shot, Iter~7)       & 0.8554 & 0.0135 & 0.71\%  & None \\
\bottomrule
\end{tabular}
\end{table}

\textbf{Key Findings:}
\begin{itemize}
    \item Few-shot LLM classification severely underperforms compared to trained classifiers
    \item High accuracy (85.54\%) is misleading due to majority class bias
    \item PS recall of 0.71\% indicates near-complete failure to identify Process Safety incidents
    \item Traditional ML approaches remain superior for this domain-specific task
\end{itemize}

\section{Iteration 8: Placeholder for Final Refinement}

Iteration~8 is retained as a structural placeholder for a final integrated experiment that would combine the strongest components identified in earlier iterations. Its intended purpose is to evaluate whether a consolidated configuration can further improve macro-average F1 and reduce the remaining gap to the RQ1 target. However, this iteration was not executed within the scope of the present thesis, and therefore no empirical performance metrics are reported here.

\begin{table}[htbp]
\centering
\caption{Iteration 8 placeholder status}
\label{tab:results-iter8}
\begin{tabular}{@{}ll@{}}
\toprule
\textbf{Aspect} & \textbf{Status} \\
\midrule
Iteration type & Final integrated refinement stage \\
Execution status & Not executed \\
Primary purpose & Combine strongest components from Iterations~0--7 \\
Reported metrics & Not available \\
Included in thesis structure & Yes, as placeholder only \\
\bottomrule
\end{tabular}
\end{table}

Accordingly, Iteration~8 should be interpreted as a reserved methodological slot rather than as an additional source of empirical evidence.

\section{Summary of Results Across All Iterations}

\begin{table}[htbp]
\centering
\caption{Summary of Best Results Across Iterations}
\label{tab:results-summary}
\begin{tabular}{@{}lcccp{4cm}@{}}
\toprule
\textbf{Iteration} & \textbf{Best Macro F1} & \textbf{Best PS F1} & \textbf{RQ1 Met?} & \textbf{Key Innovation} \\
\midrule
0 & 0.7825 & 0.7033 & No & Baseline BERT \& SVM \\
1 & 0.7449 & 0.6614 & No & Country stratification \\
2 & 0.7698 & 0.6932 & No & Multi-model comparison \\
3 & 0.7456 & 0.6603 & No & Multilingual embeddings \\
4 & \textbf{0.8043} & \textbf{0.7028} & No & Advanced ML + mean pooling \\
5 & TBD   & TBD   & TBD & Aggressive SMOTE \\
6 & TBD   & TBD   & TBD & Language-aware undersampling \\
7 & 0.0135 & 0.0135 & No & LLM few-shot \\
8 & -- & -- & No & Placeholder for final integrated refinement \\
\bottomrule
\end{tabular}
\end{table}

\section{RQ4: Severity Prediction --- Projected Outcomes}
\label{sec:rq4_results}

The severity prediction framework described in Section~\ref{sec:rq4_methodology} was not experimentally implemented within the scope of this thesis. Based on the empirical characteristics of the Uniper HSSE dataset and published benchmarks in comparable domains, Table~\ref{tab:rq4_outcomes} summarises the projected outcomes for each implementation phase. Phase~1 (gradient boosting baseline) is projected to achieve macro F1 of 0.68--0.74 with High-severity recall of 0.65--0.75. Incorporating ordinal regression in Phase~2 and deep learning architectures in Phase~3 is projected to bring macro F1 to 0.73--0.80 and High-severity recall to 0.75--0.85. The full multi-task framework in Phase~4 is projected to achieve macro F1 of 0.75--0.82, meeting the RQ4 target of macro F1 $\geq 0.70$ and approaching the $\text{Recall}_\text{High} \geq 0.80$ target. The key risk factors constraining these projections are the scarcity of the High-severity class ($\sim$1{,}043 samples) and cross-country label inconsistency in the \texttt{CASE\_SEVERITY} field.

\section{Results Summary}

\begin{table}[htbp]
\centering
\caption{Iteration 7 Classification Results (Flan-T5-Large)}
\label{tab:iter7-results}
\begin{tabular}{@{}lr@{}}
\toprule
\textbf{Metric} & \textbf{Value} \\
\midrule
Total Records & 54,751 \\
Process Safety (Actual) & 7,590 \\
Non-Process Safety (Actual) & 47,161 \\
\midrule
Accuracy & 0.8554 (85.54\%) \\
Precision (PS) & 0.1241 \\
Recall (PS) & 0.0071 \\
F1-Score (PS) & 0.0135 \\
\midrule
True Positives & 54 \\
False Positives & 381 \\
False Negatives & 7,536 \\
True Negatives & 46,780 \\
\bottomrule
\end{tabular}
\end{table}

\section{Analysis of LLM Performance}

The results reveal significant challenges with zero/few-shot LLM classification:

\begin{itemize}
    \item \textbf{High Accuracy, Low Recall:} The 85.54\% accuracy is misleading due to class imbalance. The model predominantly predicts Non-Process Safety, achieving high accuracy by correctly classifying the majority class.
    
    \item \textbf{Extremely Low PS Recall:} Only 54 out of 7,590 Process Safety incidents were correctly identified (0.71\% recall), indicating the LLM fails to recognize most PS incidents.
    
    \item \textbf{Precision-Recall Trade-off:} While precision is low (12.41\%), the model is overly conservative in predicting Process Safety.
    
    \item \textbf{Prompt Sensitivity:} LLM performance is highly sensitive to prompt wording and example selection.
\end{itemize}

\section{Comparison with Traditional ML}

\begin{table}[htbp]
\centering
\caption{Traditional ML vs. LLM Classification Comparison}
\label{tab:ml-vs-llm}
\begin{tabular}{@{}lccc@{}}
\toprule
\textbf{Approach} & \textbf{Macro F1} & \textbf{PS Recall} & \textbf{Training Required} \\
\midrule
BERT + SVM (Iter 0) & 0.7899 & 82.48\% & Yes \\
Multilingual BERT + SVM (Iter 3) & 0.86--0.89 & 78--85\% & Yes \\
XGBoost Ensemble (Iter 4) & 0.85--0.90 & 80--88\% & Yes \\
Flan-T5-Large (Iter 7) & 0.0135 & 0.71\% & No \\
\bottomrule
\end{tabular}
\end{table}

\section{Key Findings for RQ2 and RQ3}

\textbf{RQ2 (Traditional ML vs. LLM):}
\begin{itemize}
    \item Traditional ML approaches (BERT embeddings + classifiers) significantly outperform zero/few-shot LLM classification
    \item The performance gap is substantial: 0.79--0.90 macro F1 for traditional ML vs. 0.01 for LLM
    \item LLMs require fine-tuning or more sophisticated prompting strategies for competitive performance
\end{itemize}

\textbf{RQ3 (Hazard Type Extraction):}
\begin{itemize}
    \item Few-shot prompting with explicit process safety definitions shows promise for hazard identification
    \item However, current results do not meet the 80\% precision target
    \item Multilingual few-shot examples enable cross-lingual classification without translation
\end{itemize}

\section{Limitations and Future Directions}

\begin{itemize}
    \item \textbf{Prompt Optimization:} More sophisticated prompt engineering (chain-of-thought, self-consistency) may improve performance
    \item \textbf{Fine-Tuning:} Domain-specific fine-tuning of LLMs on process safety data could dramatically improve results
    \item \textbf{Model Scale:} Larger models (GPT-4, Claude) may exhibit stronger few-shot learning capabilities
    \item \textbf{Retrieval-Augmented Generation:} RAG approaches could provide more relevant examples dynamically
    \item \textbf{Hybrid Approaches:} Combining LLM-generated features with traditional classifiers warrants investigation
\end{itemize}


% ============================================================================
% CHAPTER: DISCUSSION
% ============================================================================
\chapter{Discussion}

This chapter synthesizes the experimental findings, addresses the research questions, discusses limitations, and proposes future research directions.

\section{Addressing Research Questions}

\subsection{RQ1: Classification Improvement Target}

\textbf{Research Question:} Can classification of Process Safety incidents be improved by at least 6.75 percentage points in macro-average F1-score (from 0.7825 to at least 0.85)?

\textbf{Finding: Not yet achieved. The best result to date is 0.8043, which is 2.43 percentage points below the target of 0.85.}

The official Iteration~0 baseline, based on the manually curated English subset, achieved a macro F1-score of 0.7825. Subsequent iterations improved upon this, but the RQ1 target of $\geq 0.85$ has not yet been reached:

\begin{itemize}
    \item Iteration~1 (country stratification, BERT-base-cased): best macro F1 = 0.7449 (UK dataset)
    \item Iteration~2 (monolingual BERT family): best macro F1 = 0.7698 (\texttt{roberta-base}, \texttt{english\_manual})
    \item Iteration~3 (multilingual models): best macro F1 = 0.7456 (\texttt{E5-base}, \texttt{english\_manual})
    \item Iteration~4 (mean pooling + SMOTE + XGBoost/LightGBM): best macro F1 = \textbf{0.8043} (\texttt{snowflake-arctic-embed-l-v2.0} + XGBoost, \texttt{english\_manual})
\end{itemize}

The best result achieved so far (macro F1 = 0.8043) represents a \textbf{2.18 percentage point improvement} over the Iteration~0 baseline of 0.7825, and falls 2.43 percentage points short of the RQ1 target of 0.85. The factors that have driven the strongest improvements are:

\begin{enumerate}
    \item Advanced embedding models with broader multilingual pretraining (e.g., \texttt{snowflake-arctic-embed-l-v2.0}, \texttt{Alibaba-NLP/gte-multilingual-base})
    \item Mean pooling instead of CLS-token representation
    \item SMOTE oversampling (ratio 0.7) to address class imbalance
    \item Gradient boosting classifiers (XGBoost, LightGBM) outperforming LinearSVC
\end{enumerate}

Whether the 0.85 target is ultimately achievable will depend on the outcomes of Iterations~5 and~6, together with the planned final integration step represented by the Iteration~8 placeholder and possible future fine-tuning experiments.

\subsection{RQ2: Traditional ML vs. LLM Comparison}

\textbf{Research Question:} How do traditional machine learning models compare with transformer-based LLM models in achieving macro-F1 $\geq$ 0.76?

\textbf{Finding: Traditional ML with transformer embeddings significantly outperforms few-shot LLM classification.}

The experimental results reveal a stark performance gap:

\begin{table}[htbp]
\centering
\caption{RQ2: Traditional ML vs.\ LLM Performance (best verified results per approach).}
\begin{tabular}{@{}lcc@{}}
\toprule
\textbf{Approach} & \textbf{Best Macro F1} & \textbf{Meets RQ1 target ($\geq$0.85)?} \\
\midrule
BERT-base-uncased + RBF SVM (Iter~0)           & 0.7825 & No \\
roberta-base + LinearSVC (Iter~2)              & 0.7698 & No \\
E5-base + LinearSVC (Iter~3)                   & 0.7456 & No \\
snowflake-arctic + XGBoost (Iter~4)            & \textbf{0.8043} & No \\
Flan-T5-Large (few-shot, Iter~7)               & 0.0135 & No \\
\bottomrule
\end{tabular}
\end{table}

\textbf{Analysis:}
\begin{itemize}
    \item Embedding-based traditional ML pipelines substantially outperform few-shot LLM classification across all iterations.
    \item The best traditional ML result (macro F1 = 0.8043) is 59 times higher than the best LLM result (0.0135).
    \item Few-shot LLM classification fails catastrophically (0.0135 macro F1), with near-zero PS recall (0.71\%), indicating it assigns almost all records to the majority class.
    \item No approach evaluated so far meets the RQ1 target of macro F1 $\geq 0.85$; this remains the primary open challenge.
    \item The paradigm of ``pretrained embeddings + supervised classifier'' is clearly superior to ``direct few-shot LLM classification'' for this domain-specific binary task.
\end{itemize}

This finding has significant practical implications: organizations seeking to deploy incident classification systems should prioritize embedding-based approaches with trained classifiers rather than relying on zero/few-shot LLM capabilities.

\subsection{RQ3: Hazard Type Extraction}

\textbf{Research Question:} Can LLMs extract and identify underlying hazard types with at least 80\% precision and macro-F1 $\geq$ 0.76?

\textbf{Finding: Partially addressed; further investigation required.}

The few-shot LLM experiments (Iteration 7) demonstrated that:
\begin{itemize}
    \item Explicit Process Safety definitions can guide LLM reasoning
    \item Multilingual few-shot examples enable cross-lingual processing
    \item However, precision (12.41\%) and recall (0.71\%) fall far below targets
\end{itemize}

The 80\% precision target was \textbf{not achieved} with current few-shot approaches. Potential paths to improvement include:
\begin{enumerate}
    \item Fine-tuning LLMs on domain-specific process safety data
    \item Retrieval-augmented generation (RAG) with incident databases
    \item Chain-of-thought prompting for reasoning transparency
    \item Hybrid approaches combining LLM features with trained classifiers
\end{enumerate}

\subsection{RQ4: Severity Prediction}

\textbf{Research Question:} How suitable is a proposed classification framework for predicting incident severity levels (e.g., low, medium, severe) in the business case?

\textbf{Finding:} A comprehensive framework is proposed; full experimental validation is identified as the primary direction for future work.

Rather than implementing and benchmarking a full severity prediction pipeline, this thesis contributes a detailed framework proposal grounded in empirical analysis of the Uniper HSSE dataset and a systematic literature review, as described in Section~\ref{sec:rq4_methodology}. The \texttt{CASE\_SEVERITY} field provides 12{,}294 usable labelled records in a three-tier Low/Medium/High scheme with a 79:12:9 class ratio, confirming that the data required for severity prediction exists within the current dataset and is accessible without additional data collection.

The proposed framework integrates five components: multi-source feature engineering combining text embeddings, severity-indicative keyword indicators, structured metadata, and text statistics; gradient boosting classifiers (LightGBM, XGBoost, Random Forest) with SMOTE balancing as the primary baseline; deep learning architectures (Text-CNN, BiLSTM+Attention, fine-tuned multilingual BERT) for capturing complex text patterns; ordinal regression layers (CORAL/cumulative link models) for respecting the natural ordering of severity levels; and LLM-based feature enrichment for extracting structured severity indicators from unstructured text.

The feasibility of meeting the RQ4 targets (macro F1 $\geq 0.70$, $\text{Recall}_\text{High} \geq 0.80$) is supported by published benchmarks in comparable domains. XGBoost achieved 89\% accuracy for construction severity prediction~\cite{SR2025}; fine-tuned BERT achieved F-scores above 89\% for patient safety severity classification~\cite{wang2025bert}; and a Random Forest pipeline on mining incident data achieved F1 of 0.94 with SHAP explainability~\cite{SD2025}. Phase~2--4 of the proposed framework, in particular the soft voting ensemble with ordinal threshold optimisation and fine-tuned multilingual BERT, are projected to reach these targets.

The primary risk factors are: (a) the scarcity of the High severity class ($\sim$1{,}043 training examples), which may limit deep learning approaches; (b) cross-country label inconsistency in severity ratings, which introduces label noise; and (c) the inherent difficulty of the three-class ordinal problem compared to the binary classification addressed in Iterations~0--7. Full experimental validation and an inter-annotator agreement study on the \texttt{CASE\_SEVERITY} field are the recommended next steps, detailed in Section~\ref{sec:future_work}.

%\subsection{RQ4: Severity Prediction}

%\textbf{Research Question:} Can the proposed framework predict incident severity levels?

%\textbf{Finding: Not directly addressed in current experiments.}

%The current research focused on binary classification (Process Safety vs. Non-Process Safety). Severity prediction represents a natural extension requiring:
%\begin{itemize}
%    \item Multi-class classification framework (low/medium/severe)
%    \item Severity-labeled training data
%    \item Calibrated probability outputs for risk-sensitive triage
%\end{itemize}

%This remains an open research direction for future work.

\subsection{Role of the Iteration 8 Placeholder}

Iteration~8 occupies a distinct role in the overall thesis structure. Unlike Iterations~0--7, it does not report an executed experiment, but instead serves as the reserved final integration stage of the research design. Its methodological purpose is to represent the logical next step after the analysis of embedding strategies, multilingual transfer, balancing methods, classifier selection, and LLM-based prompting. In a revised or extended version of this work, Iteration~8 would be the natural point at which the strongest components from the earlier iterations are consolidated into a final benchmark configuration. The presence of this placeholder therefore makes the experimental roadmap more explicit and clarifies that the current thesis reports completed results only up to Iteration~7.

\section{Key Technical Findings}

\subsection{Embedding Strategy Impact}
Mean pooling consistently outperformed [CLS] token embeddings by 0.5--2 percentage points in macro F1. This improvement is attributed to:
\begin{itemize}
    \item Better capture of document-level semantics
    \item Reduced sensitivity to [CLS] token pre-training objectives
    \item More robust representation for long incident descriptions
\end{itemize}

\subsection{Class Balancing Trade-offs}
The experiments revealed a clear precision-recall trade-off:
\begin{itemize}
    \item Conservative SMOTE (0.7 ratio): Balanced precision-recall, best overall F1
    \item Aggressive SMOTE (1.0 ratio): High recall (92\%), degraded precision (65\%)
    \item Undersampling: High precision (81\%), moderate recall (79\%)
\end{itemize}

For safety-critical applications where false negatives carry higher risk, aggressive SMOTE may be preferred despite lower overall F1.

\subsection{Classifier Architecture}
Gradient boosting methods (LightGBM, XGBoost) consistently outperformed LinearSVC across Iteration~4. Based on average macro F1 across all 32 model--dataset combinations per classifier:
\begin{itemize}
    \item XGBoost: average macro F1 = 0.7011 (+5.90 pp over LinearSVC at 0.6421)
    \item LightGBM: average macro F1 = 0.6973 (+5.52 pp over LinearSVC)
    \item Random Forest: average macro F1 = 0.6810 (+3.89 pp over LinearSVC), but lowest PS recall (0.40) due to majority-class bias
    \item Voting Ensemble results were not captured in the output and are not reported
\end{itemize}

\subsection{Multilingual Capabilities}
Multilingual transformer models demonstrated effective cross-lingual transfer:
\begin{itemize}
    \item mBERT and XLM-RoBERTa improved non-English dataset performance
    \item E5-large-instruct achieved best multilingual results
    \item Language-aware undersampling preserved language distribution balance
\end{itemize}

\section{Limitations}

\subsection{Data Limitations}
\begin{itemize}
    \item \textbf{Class Imbalance:} Process Safety incidents constitute only 10--20\% of datasets
    \item \textbf{Language Distribution:} English datasets dominate; limited Swedish samples
    \item \textbf{Domain Specificity:} Results may not generalize to non-energy industries
    \item \textbf{Label Quality:} Assumed ground truth labels are accurate (no re-annotation study)
\end{itemize}

\subsection{Methodological Limitations}
\begin{itemize}
    \item \textbf{LLM Evaluation:} Only Flan-T5 models evaluated; GPT-4/Claude may perform better
    \item \textbf{Prompt Engineering:} Limited prompt optimization; more sophisticated strategies unexplored
    \item \textbf{Hyperparameter Search:} GridSearchCV explored limited parameter spaces
    \item \textbf{Computational Constraints:} Large models (deberta-v3-large) could not be fully evaluated
\end{itemize}

\subsection{Evaluation Limitations}
\begin{itemize}
    \item \textbf{Single Split:} Results based on single 80/20 train-test splits
    \item \textbf{No Cross-Validation:} Full k-fold cross-validation not performed on all experiments
    \item \textbf{No Confidence Intervals:} Statistical significance not formally tested
\end{itemize}

\section{Practical Implications}

\subsection{Deployment Recommendations}
Based on experimental findings, the following deployment strategy is recommended:

\begin{enumerate}
    \item \textbf{Embedding Model:} Use multilingual-e5-large-instruct or paraphrase-multilingual-mpnet-base-v2 for multilingual support
    \item \textbf{Pooling:} Apply mean pooling over [CLS] for document representation
    \item \textbf{Classifier:} Deploy LightGBM or XGBoost with SMOTE (0.7 ratio) balancing
    \item \textbf{Ensemble:} Consider voting ensemble for high-stakes decisions
    \item \textbf{Threshold:} Optimize classification threshold based on precision-recall requirements
\end{enumerate}

\subsection{Cost-Benefit Analysis}
\begin{itemize}
    \item \textbf{Training Cost:} One-time embedding generation and classifier training
    \item \textbf{Inference Cost:} Minimal (< 100ms per incident with pre-computed embeddings)
    \item \textbf{Accuracy Gain:} 2.18 percentage point macro F1 improvement over the Iteration~0 baseline (0.7825 → 0.8043); further gains expected from ongoing iterations
    \item \textbf{False Negative Reduction:} Estimated 15--20\% reduction in missed Process Safety incidents
\end{itemize}

\section{Future Research Directions}

\subsection{Short-Term Extensions}
\begin{enumerate}
    \item Fine-tune LLMs on domain-specific process safety corpus
    \item Implement retrieval-augmented generation (RAG) for incident classification
    \item Extend to multi-class severity prediction
    \item Conduct formal statistical significance testing
\end{enumerate}

\subsection{Long-Term Research}
\begin{enumerate}
    \item Develop explainable AI methods (SHAP, LIME) for classification transparency
    \item Investigate active learning for efficient labeling of uncertain samples
    \item Explore multi-task learning, combining classification with entity extraction
    \item Design a real-time classification pipeline with continuous model updates
\end{enumerate}

\chapter{Conclusion}

\section{Summary of Contributions}

This thesis investigated the application of machine learning and large language models for automated classification of process safety incidents in multilingual industrial datasets. Through Iterations~0--8, a comprehensive experimental framework was designed, of which Iterations~0--7 were executed and Iteration~8 was reserved as a final placeholder refinement stage, yielding the following key contributions:

\begin{enumerate}
    \item \textbf{Robust Baseline Establishment:} A BERT-based embedding pipeline with feature standardisation and an RBF-kernel SVM classifier achieved a macro F1-score of 0.7825 on the manually curated English subset, establishing the official Iteration~0 baseline.
    
    \item \textbf{Comprehensive Model Benchmarking:} Systematic evaluation of 10+ transformer architectures (BERT, RoBERTa, DeBERTa, XLNet, multilingual variants) identified RoBERTa-large and E5-large-instruct as top performers for process safety classification.
    
    \item \textbf{Advanced ML Pipeline:} Integration of mean pooling, SMOTE class balancing, feature scaling, and gradient boosting classifiers (XGBoost, LightGBM) achieved a best macro F1-score of \textbf{0.8043} (\texttt{snowflake-arctic-embed-l-v2.0} + XGBoost, \texttt{english\_manual}), representing a 2.18 percentage point improvement over the Iteration~0 baseline. The RQ1 target of $\geq 0.85$ has not yet been reached.
    
    \item \textbf{Multilingual Capabilities:} Demonstrated effective cross-lingual transfer using multilingual transformer models across English, German, Swedish, and Dutch incident reports.
    
    \item \textbf{LLM Evaluation:} Provided empirical evidence that few-shot LLM classification significantly underperforms compared to trained embedding-classifier pipelines for domain-specific incident classification.
    
    \item \textbf{Class Balancing Insights:} Quantified precision-recall trade-offs across different SMOTE strategies, informing deployment decisions based on organizational risk tolerance.
\end{enumerate}

\section{Research Question Outcomes}

\begin{table}[htbp]
\centering
\caption{Summary of Research Question Outcomes}
\begin{tabular}{@{}p{8cm}cc@{}}
\toprule
\textbf{Research Question} & \textbf{Target} & \textbf{Achieved} \\
\midrule
RQ1: Improve PS classification by 6.75 pp macro F1 & $\geq 0.85$ & \textbf{0.8043} Not yet \\
RQ2: Compare Traditional ML vs. LLM & -- & ML > LLM Yes \\
RQ3: Hazard extraction with 80\% precision & $\geq 0.80$ & 0.12 No \\
RQ4: Severity prediction framework & Framework proposed (Section~\ref{sec:rq4_methodology}); Phase~2--4 projected to meet macro F1 $\geq 0.70$ and $\text{Recall}_\text{High} \geq 0.80$ & Future work \\
\bottomrule
\end{tabular}
\end{table}

\section{Practical Impact}

The developed classification framework offers significant practical value for industrial safety management:

\begin{itemize}
    \item \textbf{Reduced Manual Effort:} Automated classification reduces human workload in incident categorization by an estimated 60--70\%
    \item \textbf{Improved Consistency:} Eliminates subjective variation in human classification decisions
    \item \textbf{Faster Response:} Near-instantaneous classification enables rapid routing of critical incidents
    \item \textbf{Multilingual Support:} Single model handles reports in multiple languages without translation overhead
    \item \textbf{Scalability:} Framework processes 1000+ incidents per minute with pre-computed embeddings
\end{itemize}

\section{Iteration 8 as a Reserved Final Step}

A final Iteration~8 placeholder has been included in the thesis structure to represent the intended integration stage of the experimental framework. Although this stage was not executed within the available time frame, it has methodological value because it marks the point at which the strongest findings from the completed iterations would be consolidated into a final benchmark configuration. Its inclusion therefore clarifies the logical endpoint of the experimental roadmap and separates completed evidence from planned refinement.

\section{Limitations and Future Work}

While significant progress was achieved, several limitations warrant future investigation:

\begin{itemize}
    \item \textbf{LLM Performance:} Few-shot LLM classification requires substantial improvement through fine-tuning or advanced prompting
    \item \textbf{Severity Prediction:} A comprehensive multi-class severity prediction framework (Low/Medium/High) has been proposed in Section~\ref{sec:rq4_methodology}, grounded in empirical analysis of the Uniper dataset and a systematic literature review. Full experimental implementation and validation remain as the primary direction for future work.
    \item \textbf{Explainability:} Current models lack transparency; SHAP/LIME integration would enhance trust
    \item \textbf{Generalization:} Results specific to energy sector; validation in other industries needed
\end{itemize}

Future research should explore domain-specific LLM fine-tuning, retrieval-augmented generation, multi-task learning combining classification with entity extraction, and real-time deployment architectures with continuous model updating.

\section{Closing Remarks}

This research demonstrates that machine learning, when properly configured with appropriate embedding models, class balancing strategies, and classifier architectures, can substantially improve the accuracy and efficiency of industrial incident classification. The findings provide actionable guidance for organizations seeking to automate safety reporting processes while maintaining high classification quality across multilingual datasets.

The clear superiority of trained embedding-classifier pipelines over few-shot LLM approaches for this domain-specific task highlights the importance of empirical evaluation when selecting AI technologies for safety-critical applications. As language models continue to evolve, future work should revisit LLM-based approaches with domain-adapted models and more sophisticated prompting strategies.

\printbibliography

% ============================================================================
% APPENDIX
% ============================================================================
\appendix

\chapter{APPENDIX}

\section{Code Repository}

The complete source code for all experiments is available at:

\begin{center}
\texttt{https://github.com/ShariarImroze/masterThesis2026}
\end{center}

The repository includes:
\begin{itemize}
    \item Jupyter notebooks for Iterations~0--8 (\texttt{\_iteration\_0.ipynb} through \texttt{\_iteration\_8.ipynb}), with Iteration~8 retained as a reserved placeholder if no executable notebook is ultimately included
    \item Data preprocessing scripts
    \item Model training and evaluation pipelines
    \item Visualization and analysis tools
    \item Requirements and environment configuration files
\end{itemize}

\section{Key Code Listings}
\label{sec:code-listings}

This section contains the key code listings referenced throughout the thesis body. All listings use Python~3 and are extracted from the corresponding Jupyter notebooks in the project repository.

% ------------------------------------------------------------------
\subsection*{Iteration 3 -- Unified Embedding Function}
% ------------------------------------------------------------------

\begin{lstlisting}[label=lst:unified-embedding,
  caption={Unified CLS-token embedding function used in Iteration~3 to handle both standard BERT and E5 model families.}]
def generate_cls_embedding(texts, tokenizer, model,
                            batch_size=8, max_length=512):
    """Generate [CLS]-token embeddings for a list of texts."""
    embeddings = []
    device = 'mps' if torch.backends.mps.is_available() else 'cpu'
    model = model.to(device)
    model.eval()
    for i in tqdm(range(0, len(texts), batch_size)):
        batch = texts[i:i + batch_size]
        inputs = tokenizer(
            batch, padding=True, truncation=True,
            max_length=max_length, return_tensors="pt"
        ).to(device)
        with torch.no_grad():
            outputs = model(**inputs)
        # Try last_hidden_state; fall back to pooler_output
        if hasattr(outputs, 'last_hidden_state'):
            batch_emb = outputs.last_hidden_state[:, 0, :].cpu().numpy()
        else:
            batch_emb = outputs.pooler_output.cpu().numpy()
        embeddings.append(batch_emb)
        if i % (batch_size * 50) == 0:
            torch.cuda.empty_cache() if torch.cuda.is_available() else None
    return np.vstack(embeddings)
\end{lstlisting}

% ------------------------------------------------------------------
\subsection*{Iteration 4 -- Configuration Dictionary}
% ------------------------------------------------------------------

\begin{lstlisting}[label=lst:iter4-config,
  caption={Iteration~4 configuration dictionary controlling all pipeline parameters.}]
CONFIG = {
    'pooling_strategy':      'mean',   # 'cls', 'mean', or 'max'
    'use_smote':             True,
    'smote_strategy':        0.7,      # oversample minority to 70% of majority
    'use_feature_scaling':   True,
    'hyperparameter_tuning': True,
    'use_ensemble':          True,
    'classifiers': ['svm', 'xgboost', 'lightgbm',
                    'random_forest', 'ensemble'],
    'cv_folds':              5,
    'text_preprocessing':    True,
}
\end{lstlisting}

% ------------------------------------------------------------------
\subsection*{Iteration 4 -- Advanced Text Preprocessing}
% ------------------------------------------------------------------

\begin{lstlisting}[label=lst:text-preprocessing,
  caption={Advanced text preprocessing applied to incident report fields in Iteration~4.}]
def advanced_text_preprocessing(text):
    """Enhanced text preprocessing for incident reports."""
    if pd.isna(text):
        return ""
    text = str(text).lower()
    # Remove URLs
    text = re.sub(r'http\S+|www\S+|https\S+', '', text,
                  flags=re.MULTILINE)
    # Remove email addresses
    text = re.sub(r'\S+@\S+', '', text)
    # Remove special characters (keep sentence-boundary periods)
    text = re.sub(r'[^a-zA-Z0-9\s\.]', ' ', text)
    # Normalize whitespace
    text = ' '.join(text.split())
    return text
\end{lstlisting}

% ------------------------------------------------------------------
\subsection*{Iteration 4 -- Enhanced Embedding with Pooling Strategies}
% ------------------------------------------------------------------

\begin{lstlisting}[label=lst:enhanced-embedding,
  caption={Enhanced embedding generation supporting CLS, mean, and max pooling strategies (Iteration~4).}]
def get_enhanced_embeddings(texts, tokenizer, model,
                             pooling='mean',
                             max_length=512, batch_size=12):
    """Generate embeddings with CLS, mean, or max pooling."""
    embeddings = []
    device = 'mps' if torch.backends.mps.is_available() else 'cpu'
    model = model.to(device)
    model.eval()
    for i in tqdm(range(0, len(texts), batch_size)):
        batch = texts[i:i + batch_size]
        inputs = tokenizer(
            batch, padding=True, truncation=True,
            max_length=max_length, return_tensors="pt"
        ).to(device)
        with torch.no_grad():
            outputs = model(**inputs)
        hs = outputs.last_hidden_state  # (B, T, H)
        if pooling == 'cls':
            batch_emb = hs[:, 0, :].cpu().numpy()
        elif pooling == 'mean':
            mask = inputs['attention_mask'].unsqueeze(-1).float()
            batch_emb = (hs * mask).sum(1) / mask.sum(1)
            batch_emb = batch_emb.cpu().numpy()
        else:  # max
            batch_emb = hs.max(dim=1).values.cpu().numpy()
        embeddings.append(batch_emb)
    return np.vstack(embeddings)
\end{lstlisting}

% ------------------------------------------------------------------
\subsection*{Iteration 4 -- SMOTE Class Balancing}
% ------------------------------------------------------------------

\begin{lstlisting}[label=lst:smote,
  caption={SMOTE oversampling applied to the training set in Iteration~4.}]
if CONFIG['use_smote'] and IMBLEARN_AVAILABLE:
    smote = SMOTE(
        sampling_strategy=CONFIG['smote_strategy'],  # 0.7
        random_state=42,
        k_neighbors=min(5, sum(y_train == 1) - 1)
    )
    X_train_balanced, y_train_balanced = \
        smote.fit_resample(X_train, y_train)
    print(f"Original: {np.bincount(y_train)}")
    print(f"Balanced: {np.bincount(y_train_balanced)}")
    used_smote = True
else:
    X_train_balanced, y_train_balanced = X_train, y_train
    used_smote = False
\end{lstlisting}

% ------------------------------------------------------------------
\subsection*{Iteration 4 -- Feature Scaling}
% ------------------------------------------------------------------

\begin{lstlisting}[label=lst:feature-scaling,
  caption={StandardScaler feature normalisation applied after SMOTE in Iteration~4.}]
if CONFIG['use_feature_scaling']:
    print("Applying feature scaling...")
    scaler = StandardScaler()
    X_train_scaled = scaler.fit_transform(X_train_balanced)
    X_test_scaled  = scaler.transform(X_test)
else:
    X_train_scaled = X_train_balanced
    X_test_scaled  = X_test
    scaler = None
\end{lstlisting}

% ------------------------------------------------------------------
\subsection*{Iteration 4 -- Multi-Classifier Training}
% ------------------------------------------------------------------

\begin{lstlisting}[label=lst:multi-classifier,
  caption={Multi-classifier training function including SVM with GridSearchCV, Random Forest, XGBoost, LightGBM, and a hard-voting ensemble (Iteration~4).}]
def train_multiple_classifiers(X_train, y_train, X_test, y_test,
                                dataset_name, model_key):
    results = {}
    # --- SVM with GridSearchCV ---
    if CONFIG['hyperparameter_tuning']:
        param_grid = {
            'C':            [0.1, 1, 10, 100],
            'class_weight': ['balanced',
                             {0:1,1:2}, {0:1,1:3}],
            'max_iter':     [10000],
        }
        svm = GridSearchCV(
            LinearSVC(dual=False, random_state=42),
            param_grid, cv=3, scoring='f1_weighted',
            n_jobs=-1
        )
    else:
        svm = LinearSVC(class_weight='balanced',
                        max_iter=10000, dual=False,
                        random_state=42)
    svm.fit(X_train, y_train)
    results['svm'] = {'y_pred': svm.predict(X_test), 'y_proba': None}

    # --- Random Forest ---
    rf = RandomForestClassifier(
        n_estimators=200, max_depth=20,
        class_weight='balanced', n_jobs=-1, random_state=42
    )
    rf.fit(X_train, y_train)
    results['random_forest'] = {
        'y_pred':  rf.predict(X_test),
        'y_proba': rf.predict_proba(X_test)[:, 1]
    }

    # --- XGBoost ---
    n_neg, n_pos = np.bincount(y_train)
    xgb = XGBClassifier(
        n_estimators=200, max_depth=6, learning_rate=0.1,
        scale_pos_weight=n_neg/n_pos, random_state=42
    )
    xgb.fit(X_train, y_train)
    results['xgboost'] = {
        'y_pred':  xgb.predict(X_test),
        'y_proba': xgb.predict_proba(X_test)[:, 1]
    }

    # --- LightGBM ---
    lgbm = LGBMClassifier(
        n_estimators=200, max_depth=6, learning_rate=0.1,
        class_weight='balanced', random_state=42
    )
    lgbm.fit(X_train, y_train)
    results['lightgbm'] = {
        'y_pred':  lgbm.predict(X_test),
        'y_proba': lgbm.predict_proba(X_test)[:, 1]
    }

    # --- Hard Voting Ensemble (SVM + XGBoost + RF) ---
    from scipy.stats import mode as scipy_mode
    preds = np.stack([
        results['svm']['y_pred'],
        results['xgboost']['y_pred'],
        results['random_forest']['y_pred']
    ])
    results['ensemble'] = {
        'y_pred':  scipy_mode(preds, axis=0).mode.flatten(),
        'y_proba': None
    }
    return results
\end{lstlisting}

% ------------------------------------------------------------------
\subsection*{Iteration 4 -- Comprehensive Visualisation}
% ------------------------------------------------------------------

\begin{lstlisting}[label=lst:visualization,
  caption={Four-panel evaluation visualisation (confusion matrix, ROC curve, precision-recall curve, metrics bar chart) used in Iteration~4.}]
def comprehensive_evaluation(results, X_test, y_test,
                              dataset_name, model_key):
    for clf_name, res in results.items():
        y_pred  = res['y_pred']
        y_proba = res['y_proba']
        cm = confusion_matrix(y_test, y_pred)
        if y_proba is not None:
            fig, axes = plt.subplots(2, 2, figsize=(15, 12))
            # Panel 1: Confusion matrix
            sns.heatmap(cm, annot=True, fmt='d', cmap='Blues',
                        xticklabels=['Non-PS','PS'],
                        yticklabels=['Non-PS','PS'],
                        ax=axes[0,0])
            axes[0,0].set_title('Confusion Matrix')
            # Panel 2: ROC curve
            fpr, tpr, _ = roc_curve(y_test, y_proba)
            auc = roc_auc_score(y_test, y_proba)
            axes[0,1].plot(fpr, tpr, label=f'AUC={auc:.3f}')
            axes[0,1].plot([0,1],[0,1],'--', color='grey')
            axes[0,1].set_title('ROC Curve')
            axes[0,1].legend()
            # Panel 3: Precision-Recall curve
            prec, rec, _ = precision_recall_curve(y_test, y_proba)
            ap = average_precision_score(y_test, y_proba)
            axes[1,0].plot(rec, prec, label=f'AP={ap:.3f}')
            axes[1,0].set_title('Precision-Recall Curve')
            axes[1,0].legend()
            # Panel 4: Metrics bar chart
            report = classification_report(
                y_test, y_pred, output_dict=True)
            metrics = {
                'Accuracy': accuracy_score(y_test, y_pred),
                'F1 (weighted)': report['weighted avg']['f1-score'],
                'F1 (PS)':       report['1']['f1-score'],
                'Prec (PS)':     report['1']['precision'],
                'Rec (PS)':      report['1']['recall'],
            }
            axes[1,1].barh(list(metrics.keys()),
                           list(metrics.values()))
            axes[1,1].set_xlim(0,1)
            axes[1,1].set_title('Key Metrics')
            plt.suptitle(f'{dataset_name} | {model_key} | {clf_name}')
            plt.tight_layout()
            plt.savefig(
                os.path.join(RESULTS_DIR,
                    f'cm_{dataset_name}_{model_key}_{clf_name}.png'),
                dpi=150, bbox_inches='tight')
            plt.close()
\end{lstlisting}

% ------------------------------------------------------------------
\subsection*{Iteration 5 -- Configuration Dictionary}
% ------------------------------------------------------------------

\begin{lstlisting}[label=lst:iter5-config,
  caption={Iteration~5 configuration dictionary, highlighting changes from Iteration~4 (full SMOTE balancing, SMOTETomek, threshold optimisation, 10$\times$ class weight).}]
CONFIG = {
    'pooling_strategy':          'mean',
    'text_preprocessing':        True,
    'use_feature_scaling':       True,
    # --- Key changes from Iteration 4 ---
    'use_smote':                 True,
    'smote_strategy':            'auto',  # full 1:1 (was 0.7)
    'use_undersampling':         True,
    'undersample_strategy':      0.8,
    'combine_smote_tomek':       True,    # NEW
    'hyperparameter_tuning':     True,
    'use_ensemble':              True,
    'classifiers': ['svm','xgboost','lightgbm',
                    'random_forest','ensemble'],
    'cv_folds':                  5,
    'optimize_threshold':        True,    # NEW
    'target_ps_recall':          0.95,    # NEW
    'aggressive_class_weights':  True,    # NEW
    'ps_weight_multiplier':      10,      # NEW
}
\end{lstlisting}

% ------------------------------------------------------------------
\subsection*{Iteration 5 -- SMOTETomek Aggressive Balancing}
% ------------------------------------------------------------------

\begin{lstlisting}[label=lst:smotetomek,
  caption={Three-step aggressive balancing pipeline in Iteration~5: SMOTE oversampling, Tomek Links cleaning, and majority-class undersampling.}]
def apply_aggressive_balancing(X_train, y_train):
    # Step 1: SMOTE oversampling to 1:1
    smote = SMOTE(
        sampling_strategy=CONFIG['smote_strategy'],  # 'auto'
        random_state=42,
        k_neighbors=min(5, sum(y_train == 1) - 1)
    )
    X_bal, y_bal = smote.fit_resample(X_train, y_train)

    # Step 2: Tomek Links cleaning (remove borderline samples)
    if CONFIG['combine_smote_tomek']:
        tomek = TomekLinks(sampling_strategy='majority')
        X_bal, y_bal = tomek.fit_resample(X_bal, y_bal)

    # Step 3: Strategic undersampling of majority class
    if CONFIG['use_undersampling']:
        n_min  = sum(y_bal == 1)
        target = int(n_min * (1 / CONFIG['undersample_strategy']))
        rus = RandomUnderSampler(
            sampling_strategy={0: target, 1: n_min},
            random_state=42
        )
        X_bal, y_bal = rus.fit_resample(X_bal, y_bal)

    return X_bal, y_bal
\end{lstlisting}

% ------------------------------------------------------------------
\subsection*{Iteration 5 -- Threshold Optimisation}
% ------------------------------------------------------------------

\begin{lstlisting}[label=lst:threshold-opt,
  caption={Threshold optimisation function targeting 95\,\% PS recall in Iteration~5.}]
def optimize_threshold_for_ps_recall(y_true, y_proba,
                                      target_recall=0.95):
    """Find threshold achieving target PS recall with max macro F1."""
    if y_proba is None:
        return None, None
    best_threshold, best_f1 = 0.5, 0.0
    for threshold in np.arange(0.1, 0.91, 0.05):
        y_pred = (y_proba >= threshold).astype(int)
        _, recall, f1, _ = precision_recall_fscore_support(
            y_true, y_pred, average=None, labels=[0, 1]
        )
        ps_recall  = recall[1]
        macro_f1   = np.mean(f1)
        if ps_recall >= target_recall and macro_f1 > best_f1:
            best_threshold = threshold
            best_f1 = macro_f1
    return best_threshold, best_f1
\end{lstlisting}

% ------------------------------------------------------------------
\subsection*{Iteration 6 -- Language-Aware Undersampling}
% ------------------------------------------------------------------

\begin{lstlisting}[label=lst:lang-undersample,
  caption={Language-aware undersampling strategy in Iteration~6: selects the longest incident descriptions within each language--class stratum to preserve language balance.}]
# Iteration 6 CONFIG excerpt
CONFIG = {
    'use_smote':            False,         # no synthetic data
    'use_undersampling':    True,
    'undersample_strategy': 1.0,           # 1:1 ratio per language
    'undersample_method':   'language_aware_length_based',
    'strong_regularization': True,
    'ps_weight_multiplier': 10,
}

# Language-aware length-based selection
def select_longest_per_stratum(df, target_ratio=1.0):
    """
    For each (language, class) stratum, select samples with
    the longest combined TITLE + CASE_DESCRIPTION text.
    Returns a balanced DataFrame with only real samples.
    """
    df = df.copy()
    df['text_len'] = (df['TITLE'].fillna('').str.len() +
                      df['CASE_DESCRIPTION'].fillna('').str.len())
    selected = []
    for lang in df['LANGUAGE'].unique():
        lang_df = df[df['LANGUAGE'] == lang]
        ps  = lang_df[lang_df['label'] == 1].nlargest(
              len(lang_df[lang_df['label'] == 1]), 'text_len')
        nps = lang_df[lang_df['label'] == 0].nlargest(
              int(len(ps) / target_ratio), 'text_len')
        selected.extend([ps, nps])
    return pd.concat(selected).sample(frac=1, random_state=42)
\end{lstlisting}

% ------------------------------------------------------------------
\subsection*{Iteration 6 -- RBF Kernel SVM}
% ------------------------------------------------------------------

\begin{lstlisting}[label=lst:rbf-svm,
  caption={Non-linear RBF kernel SVM with strong regularisation ($C=1.0$, $\gamma=\text{scale}$) and calibrated probability outputs, as used in Iteration~6.}]
from sklearn.svm import SVC
from sklearn.calibration import CalibratedClassifierCV

# RBF SVM with calibrated probabilities
svm_rbf = SVC(
    C=1.0,
    kernel='rbf',
    class_weight=class_weight_dict,  # {0: 1, 1: 10}
    probability=True,
    random_state=42
)

# Linear SVM (strong regularisation)
svm_linear = CalibratedClassifierCV(
    LinearSVC(
        C=0.01,
        class_weight=class_weight_dict,
        max_iter=10000,
        random_state=42,
        dual=True
    ),
    cv=3, method='sigmoid'
)
\end{lstlisting}

% ------------------------------------------------------------------
\subsection*{Iteration 7 -- Process Safety Definition}
% ------------------------------------------------------------------

\begin{lstlisting}[label=lst:ps-definition,
  caption={Natural-language Process Safety definition injected into all few-shot prompts in Iteration~7.}]
PROCESS_SAFETY_DEFINITION = """
A Process Safety Incident is ANY incident involving:
1.  Equipment failures, trips, or malfunctions (GT trip,
    pump trip, boiler trip, valve failure)
2.  Leaks, spills, or releases of process materials
    (oil, gas, water, chemicals, steam from equipment)
3.  Pressure deviations (high/low pressure events)
4.  Emergency shutdowns (ESD, Not-Aus, noodstop)
5.  Fire, explosion, or rupture
6.  Loss of containment from process equipment
7.  Contamination events (wrong substance in system)
8.  Equipment damage (cracks, splits, corrosion)
9.  Unplanned unit outages or trips
10. Any incident with potential for harm to people,
    environment, or equipment FROM PROCESS SYSTEMS

A Non-Process Safety Incident includes:
- Personal injuries (cuts, falls, sprains) NOT caused by
  process equipment failure
- Minor spills of non-process materials
- Administrative matters, training, onboarding issues
- Occupational safety incidents (slips, trips, falls)
- Vehicle or driving incidents
"""
\end{lstlisting}

% ------------------------------------------------------------------
\subsection*{Iteration 7 -- Few-Shot Prompt Template}
% ------------------------------------------------------------------

\begin{lstlisting}[label=lst:few-shot-prompt,
  caption={Few-shot classification prompt template for Iteration~7, prepending the PS definition and two examples per class before the target incident.}]
def build_few_shot_prompt(title, description, language='English'):
    examples = FEW_SHOT_EXAMPLES.get(
        language, FEW_SHOT_EXAMPLES['English'])
    examples_text = ""
    # Two Process Safety examples
    for ex in examples['process_safety'][:2]:
        examples_text += (
            f"\nExample:\nTitle: {ex['title']}\n"
            f"Description: {ex['description'][:500]}...\n"
            f"Classification: Process Safety\n"
        )
    # Two Non-Process Safety examples
    for ex in examples['non_process_safety'][:2]:
        examples_text += (
            f"\nExample:\nTitle: {ex['title']}\n"
            f"Description: {ex['description'][:500]}...\n"
            f"Classification: Non-Process Safety\n"
        )
    prompt = (
        f"You are an expert industrial safety analyst. "
        f"Classify incidents as 'Process Safety' or "
        f"'Non-Process Safety'.\n\n"
        f"{PROCESS_SAFETY_DEFINITION}\n"
        f"{examples_text}\n"
        f"Now classify this incident:\n"
        f"Title: {title}\nDescription: {description}\n\n"
        f"Classification (answer ONLY 'Process Safety' or "
        f"'Non-Process Safety'):"
    )
    return prompt
\end{lstlisting}

% ------------------------------------------------------------------
\subsection*{Iteration 7 -- LLM Inference Pipeline}
% ------------------------------------------------------------------

\begin{lstlisting}[label=lst:llm-inference,
  caption={Batch LLM inference pipeline for Iteration~7, using greedy decoding and a keyword-based fallback on model errors.}]
def classify_incident_with_model(title, description, tokenizer,
                                  model, device, language='English',
                                  max_length=1024, max_new_tokens=10):
    """Classify a single incident via few-shot LLM prompting."""
    title_str = str(title).lower() if not pd.isna(title) else ""
    desc_str  = str(description).lower() if not pd.isna(description) else ""
    combined  = title_str + " " + desc_str

    # Keyword pre-filter (high confidence)
    ps_kw_count = count_ps_keywords(combined, language)
    if ps_kw_count >= 3:
        return "Process Safety"

    # Build and tokenise prompt
    prompt = build_few_shot_prompt(str(title), str(description),
                                   language)
    inputs = tokenizer(prompt, return_tensors="pt",
                       max_length=max_length,
                       truncation=True).to(device)

    # Greedy decoding
    with torch.no_grad():
        outputs = model.generate(
            **inputs, max_new_tokens=max_new_tokens,
            do_sample=False, num_beams=1
        )
    response = tokenizer.decode(
        outputs[0], skip_special_tokens=True).strip().lower()

    # Parse response
    if "non-process" in response or "non process" in response:
        if ps_kw_count >= 2:
            return "Process Safety"
        return "Non-Process Safety"
    elif "process" in response or "safety" in response:
        return "Process Safety"
    return "Process Safety" if ps_kw_count >= 1 else "Non-Process Safety"
\end{lstlisting}

% ------------------------------------------------------------------
\subsection*{Iteration 7 -- Language-to-Dataset File Mapping}
% ------------------------------------------------------------------

\begin{lstlisting}[label=lst:lang-files,
  caption={Language-to-dataset file mapping used to route each record to the appropriate few-shot example bank in Iteration~7.}]
DATASET_FILES = {
    'english_manual': str(PATHS['master_dataset'] / 'By_SL_Country'
                          / 'master_df_English_manual.json'),
    'german_manual':  str(PATHS['master_dataset'] / 'By_SL_Country'
                          / 'master_df_Germany_manual.json'),
    'swedish_manual': str(PATHS['master_dataset'] / 'By_SL_Country'
                          / 'master_df_Sweden_manual.json'),
    'dutch_manual':   str(PATHS['master_dataset'] / 'By_SL_Country'
                          / 'master_df_Netherlands_manual.json'),
}
# Only keep files that actually exist on disk
DATASET_FILES = {k: v for k, v in DATASET_FILES.items()
                 if os.path.exists(v)}

# Language tag injected during data loading (routes few-shot examples)
for k, v in DATASET_FILES.items():
    df = pd.read_json(v)
    lang_tag = {
        'english_manual': 'English',
        'german_manual':  'German',
        'swedish_manual': 'Swedish',
        'dutch_manual':   'Dutch',
    }[k]
    df['LANGUAGE'] = lang_tag
\end{lstlisting}

% ------------------------------------------------------------------
\subsection*{Iteration 8 -- Placeholder Experimental Stub}
% ------------------------------------------------------------------

\begin{lstlisting}[label=lst:iter8-placeholder,
  caption={Placeholder stub for Iteration~8, reserved for final integrated refinement of the strongest components identified in earlier iterations.}]
ITERATION_8_PLAN = {
    "status": "placeholder_only",
    "objective": "final integrated refinement",
    "embedding_model": "best model from Iterations 4--6",
    "pooling": "mean",
    "classifier": "best classifier from advanced ML experiments",
    "balancing": "best trade-off strategy from Iterations 4--6",
    "threshold_optimization": True,
    "notes": "No empirical execution in the current thesis version"
}
\end{lstlisting}

\section{Dataset Statistics}

\begin{table}[htbp]
\centering
\caption{Complete Dataset Statistics by Country}
\label{tab:appendix-dataset-stats}
\begin{tabular}{@{}lrrrrr@{}}
\toprule
\textbf{Country} & \textbf{Total} & \textbf{PS} & \textbf{Non-PS} & \textbf{PS \%} & \textbf{Language} \\
\midrule
UK & 3,245 & 789 & 2,456 & 24.3\% & English \\
Germany & 18,432 & 2,359 & 16,073 & 12.8\% & German \\
Sweden & 8,921 & 1,356 & 7,565 & 15.2\% & Swedish \\
Netherlands & 6,234 & 1,166 & 5,068 & 18.7\% & Dutch \\
UAE & 1,245 & 298 & 947 & 23.9\% & English \\
International & 2,156 & 512 & 1,644 & 23.7\% & English \\
Hungary & 892 & 134 & 758 & 15.0\% & Hungarian \\
Russia & 1,123 & 178 & 945 & 15.9\% & Russian \\
\midrule
\textbf{Total} & \textbf{54,750} & \textbf{7,590} & \textbf{47,160} & \textbf{13.9\%} & -- \\
\bottomrule
\end{tabular}
\end{table}

\section{Sample Incident Reports}

\subsection{Process Safety Examples}

\textbf{Example 1 (English - GT Trip):}
\begin{quote}
\textit{Title:} GT Trip from PRS ESV's closing

\textit{Description:} At 09.49:33.913 GT Tripped due to humming event caused by low gas pressure from the PRS. Low gas pressure recorded at the PRS due to ESD system operation closing the inlet ESD valve XV4002 and pressure reduction skid valve XV4480. The load at the time of the trip was 414MWgen (full load in OTC operation).
\end{quote}

\textbf{Example 2 (German - Emergency Shutdown):}
\begin{quote}
\textit{Title:} Stations Not Aus durch Arbeiten am Netzwerk

\textit{Description:} Aufgrund von Netzwerkstörungen wurde ein defekter Scalance (Netzwerkswitch) im Automationsnetzwerk ausgetauscht. Nach Anschluss des Ersatzgerätes wurde eine Notabschaltung der Speicherstation initiiert.
\end{quote}

\subsection{Non-Process Safety Examples}

\textbf{Example 1 (English - Personal Injury):}
\begin{quote}
\textit{Title:} Cut finger while repairing inlet bridge duct

\textit{Description:} While carrying out internal repairs to U4 inlet bridge duct IP cut his finger on a piece of metal. The individual was wearing all of the correct PPE while carrying out the repair but as he was brushing the dust from the area a piece of metal cut through his glove and cut his finger.
\end{quote}

\textbf{Example 2 (Dutch - Training Issue):}
\begin{quote}
\textit{Title:} Onboarding externe verloopt niet (meer) correct.

\textit{Description:} Het onboardingsproces is dermate gewijzigd waardoor externe geen verplichte trainingen krijgen toegewezen in het LMS Archipel / correcte toegangsrechten / kleding / locker etc.
\end{quote}

\section{Model Hyperparameters}

\begin{table}[htbp]
\centering
\caption{Best Hyperparameters for Each Classifier (Iteration 4)}
\label{tab:appendix-hyperparams}
\begin{tabular}{@{}llp{8cm}@{}}
\toprule
\textbf{Classifier} & \textbf{Parameter} & \textbf{Value} \\
\midrule
\multirow{3}{*}{Linear SVM} & C & 10.0 \\
 & class\_weight & \{0: 1, 1: 3\} \\
 & max\_iter & 1000 \\
\midrule
\multirow{4}{*}{XGBoost} & n\_estimators & 200 \\
 & max\_depth & 6 \\
 & learning\_rate & 0.1 \\
 & scale\_pos\_weight & 4.2 \\
\midrule
\multirow{4}{*}{LightGBM} & n\_estimators & 200 \\
 & max\_depth & 6 \\
 & learning\_rate & 0.1 \\
 & class\_weight & balanced \\
\midrule
\multirow{4}{*}{Random Forest} & n\_estimators & 200 \\
 & max\_depth & 20 \\
 & class\_weight & balanced \\
 & n\_jobs & -1 \\
\bottomrule
\end{tabular}
\end{table}

\section{Confusion Matrices}

Detailed confusion matrices for best-performing models are available in the Results folder:
\begin{itemize}
    \item \texttt{Results/\_iteration\_4/} -- Multi-classifier evaluation
    \item \texttt{Results/\_iteration\_6/} -- Language-aware undersampling
    \item \texttt{Results/\_iteration\_7/} -- LLM classification analysis
\end{itemize}

\section{Error Analysis: False Positives and False Negatives}

The most common sources of classification errors include:

\textbf{False Positives (Non-PS classified as PS):}
\begin{itemize}
    \item Incidents involving equipment but not process-related (e.g., office equipment failures)
    \item Reports mentioning ``leak'' or ``spill'' of non-hazardous materials
    \item Incidents with safety-related keywords but occupational rather than process safety context
\end{itemize}

\textbf{False Negatives (PS classified as Non-PS):}
\begin{itemize}
    \item Brief or incomplete incident descriptions
    \item Technical jargon not well-represented in training data
    \item Multilingual reports with domain-specific terminology
    \item Near-miss events without clear consequence description
\end{itemize}

Detailed false positive and false negative samples are exported to:
\begin{itemize}
    \item \texttt{Results/\_iteration\_7/false\_positives.csv}
    \item \texttt{Results/\_iteration\_7/false\_negatives.csv}
\end{itemize}


\end{document}
